\documentclass[14pt, a5paper, twoside, extrafontsizes]{memoir}

\def \sederTitle{סדר קדשים}
\def \tractateTitle{\sederTitle} % for running heads until the first tractate
\usepackage{lettrine}

% indentation of the first word on the first line after the drop capital
\setlength{\DefaultFindent}{6pt}

% extra indentation for the first word on the next line
\setlength{\DefaultNindent}{0pt}

\renewcommand*{\DefaultLoversize}{-.2}
\renewcommand*{\DefaultLraise}{.6}
\renewcommand{\LettrineTextFont}{\relax}

\renewcommand*{\LettrineFont}{%
 \LettrineFontHook
 \fontsize{18pt}{18pt}%
 \selectfont
}

% \leadword{first word [to be enlarged]}{second word}
\newcommand{\leadword}[2]{%
  \lettrine[realheight=true]{#1}{#2}
}


\usepackage{graphicx}
\usepackage{pgffor}
\usepackage{anysize}
\usepackage[all]{nowidow}
\usepackage{relsize}
\usepackage{ragged2e}

\usepackage[pdfauthor={יהודה הנשיא},
  pdftitle=\sederTitle,
  pdfsubject={תורה שבעל פה},
  pdfproducer={סיימון מונטגיו, הוצאת חטף סגול},
  pdfdirection=R2L,
  pdfcreator=xelatex,
  colorlinks=false, linkcolor=blue,
  linkbordercolor={.6 .6 .6},
  plainpages=true]{hyperref}
\makeatletter
\Hy@AtBeginDocument{%
  % Override border definition set with colorlinks=true
  \def\@pdfborder{0 0 1}
  % Overrides border style set with colorlinks=true
  % Hyperlink border style will be underline of width 1pt
  \def\@pdfborderstyle{/S/U/W 1}
}
\makeatother

\usepackage{polyglossia}
\setmainlanguage[numerals=hebrew]{hebrew}
\setmainfont{SBL Hebrew}
\setotherlanguage{english}

% General formatting
\marginsize{1.2in}{.8in}{.5in}{.5in}
\renewcommand{\baselinestretch}{1.1}
\parskip 8pt
\parindent 0pt
\makeevenhead{myheadings}{\thepage}{}{\texthebrew{\sederTitle}}
\makeoddhead{myheadings}{\texthebrew{\tractateTitle}}{}{\thepage}
\makepagestyle{mytitlepage}

\newcounter{perekCount}
\newcounter{halachaCount}[perekCount]
\newcounter{tractateCount}
\newcounter{perakim}

% \tractatehead{title for printing}{hyperlink reference}
\newcommand{\tractatehead}[2]{%
 \cleardoublepage
 \setcounter{perekCount}{0}
 \thispagestyle{empty}
 \def \tractateTitle{#1}
{\centering\LARGE\bfseries \hypertarget{#2}{#1}

  }}

% \summarycontentsline{tractate name}{hyperlink reference}
%                     {number of chapters [numeric]}
\newcommand{\summarycontentsline}[3]{%
\setcounter{perakim}{#3}%
\stepcounter{tractateCount}%
\textlarger{\alph{tractateCount}} \hyperlink{#2}{#1} \alph{perakim} פרקים\hdot%
}

\newcommand{\summarycontents}[1]{
  \centerlastline
  \foreach \linkText / \linkName / \numChapters in #1 {
      \summarycontentsline{\linkText}{\linkName}{\numChapters}
  } 
}

% \sederContents{number of tractates [hebrew]}
%               {list of tractates}{total number of chapters [hebrew]}
\newcommand{\sederContents}[3]{
  {
    \centering
    {\Large משניות \sederTitle}

  יש בו #1 מסכתות וזהו סדורן

    \begin{minipage}{0.9\textwidth}
    \small
    \summarycontents{#2}

    בין הכל \textlarger[2]{#3} פרקים׃
    \end{minipage}

  }
}

% centered rubric
\newcommand{\crubric}[1]{{\centering\small\bfseries{#1}

  }}

% centered header (keep with next paragraph)
\newcommand{\chead}[1]{%
  \vspace{0.25em}
  \needspace{4\baselineskip}
  \crubric{#1}
}

% centered footer
\newcommand{\cfoot}[1]{%
  \vspace{0.25em}
  \crubric{#1}
}

\newcommand{\mispar}[1]{{\footnotesize\textsuperscript{\alph{#1}}}}

% There is no Unicode character for the raised dot used in many Hebrew books
% as mid-sentence punctuation. This is a hack to replace it using
% U+005C4 HEBREW MARK UPPER DOT
\sfcode`\^^^^05c43000
\newcommand{\hdot}{\raisebox{-0.5\height}{^^^^2009^^^^2009^^^^05c4} }

\tolerance=8192

% \perek{first word}{rest of first halacha}
% First halacha in a perek
\newcommand{\perek}[2]{%
 \stepcounter{perekCount}
 \perekbar{פרק \alph{perekCount}}{#1}{#2}
}

% \baraita[subhead]{first word}{rest of first halacha}
% First halacha in a baraita (e.g. after the last perek of Bikkurim)
\newcommand{\baraita}[3][]{%
 \setcounter{halachaCount}{0}
 \perekbar[#1]{ברייתא}{#2}{#3}
}

% \perekbar[subhead]{heading}{first word}{rest of first halacha}
% Called by \perek and \baraita
\newcommand{\perekbar}[4][]{%
 \vspace{0.5em}
 \stepcounter{halachaCount}
 \needspace{2\baselineskip}
 \crubric{#2}

 \centerlastline
 #1

 \leadword{#3}#4
}

% \halacha{text}
% Prints a halacha
\newcommand{\halacha}[1]{%
 \stepcounter{halachaCount}
 \mispar{halachaCount}
 \centerlastline
 #1\nowidow
}

\newcommand{\selik}{~~{\footnotesize סליק}}

% \titlepage{Hebrew chronogram for year of issue}
%           {Gregorian year of issue}
\newcommand{\titlepage}[2]{%
\Centering

  {\huge\bfseries משניות}

  {\LARGE\bfseries\sederTitle}

  \vspace{1em}

  \vspace{0.5em}
  על פי מהדורת הרב דוד אלטאראס ז״ל\\שהודפס בשנת ה׳תק״ע בעיר פיסא

    \vspace{8em}
    פה {\LARGE ירושלים} יע״א

    {\small שנת} #1 {\small לפ״ק}

    \vspace{2em}
      \includegraphics[width=10mm]{hatafSegolLogoNoText.png}\\

    \vspace{.5em}
  {
    \bfseries הוצאת חטף־סגול\\#2

 {\footnotesize  הובא לדפוס באמצעות \XeLaTeX\ ע״י שמעון מונטגיו}

  }

}

% code for special handling of specific numbers by "Circumscribe" copied from https://tex.stackexchange.com/questions/300008/modify-specific-hebrew-alpha-numerals-on-page-number
% with a few modifications to the function names
\makeatletter
\let\@hebrew@numeral@usual\@hebrew@numeral  %% <- store old definition
\newcommand*\@hebrew@numeral@special[1]{%   %% <- new definition
  \ifcsdef{specialnum@\number#1}            %% <- if this number is listed
    {\csuse{specialnum@\number#1}}          %% <- give it special handling
    {\@hebrew@numeral@usual{#1}}%           %% <- otherwise use usual handling
}
\let\@hebrew@numeral\@hebrew@numeral@special %% <- replace old definition
\makeatother

%% Declare numbers for special handling:
\newcommand*\newspecialnum[2]{\csdef{specialnum@#1}{#2}}
\newspecialnum{16}{יו}    % At least some Livorno printings don't use טז for 16
\newspecialnum{116}{קיו}
\newspecialnum{216}{ריו}
\newspecialnum{316}{שיו}
\newspecialnum{416}{תיו}
\newspecialnum{516}{תקיו}
\newspecialnum{616}{תריו}
\newspecialnum{716}{תשיו}
\newspecialnum{816}{תתיו}
\newspecialnum{916}{תתקיו}
\newspecialnum{18}{חי}    % Add some other "specials" while we're here
\newspecialnum{270}{ער}
\newspecialnum{272}{ערב}
\newspecialnum{275}{ערה}



\begin{document}
\frontmatter
\pagestyle{myheadings}
\thispagestyle{mytitlepage}
\begin{hebrew}
  \pagenumbering{arabic}
  \titlepage{ונשלמה פרים {\Large שׄפׄתׄ}נ{\Large וׄ}}
    {2025}
\clearpage
\thispagestyle{empty}

\centerlastline
{
\small
\parskip 2pt
{\centering
  \scriptsize
  תפלת ר׳ נחוניא בן הקנה שחייב לאומרה כל הנכנס לבית המדרש ללמוד

}

יְהִי רָצוֹן מִלְּפָנֶיךָ יְיָ אֱלֹהַי שֶׁלֹּא יֶאֱרַע דְּבַר תְּקָלָה עַל יָדִי וְלֹא אֶכָּשֵׁל בִּדְבַר הֲלָכָה
וְיִשְׂמְחוּ בִי חֲבֵירַי וְלֹא אוֹמַר עַל טָהוֹר טָמֵא וְלֹא עַל טָמֵא טָהוֹר וְלֹא עַל מֻתָּר אָסוּר
וְלֹא עַל אָסוּר מֻתָּר וְלֹא יִכָּֽשְׁלוּ חֲבֵירַי בִּדְבַר הֲלָכָה וְאֶשְׂמַח בָּהֶם׃

{\centering
  {\scriptsize
    והאר״י זצוק״ל היה אומר אחר תפלה זו שני פסוקים אלו

    }
    כִּי יְיָ יִתֵּן חָכְמָה מִפִּיו דַּעַת וּתְבוּנָה׃\\
    גַּל עֵינַי וְאַבִּיטָה נִפְלָאוֹת מִתּוֹרָתֶךָ׃

    \vspace{1em}
   {\larger קודם הלימוד}

    {\scriptsize כתוב בספר תוצאות חיים כשפותח הספר ללמוד בו יאמר}
}

\begin{tp}{396}
הִנְּנִי רוֹצֶה לִלְמוֹד כְּדֵי שֶׁיְּבִיאֵנִי תַּלְמוּד זֶה לִידֵי מַעֲשֶׂה וְלִידֵי מִדּוֹת יְשָׁרוֹת\hdot
וְלִידֵי יְדִיעַת הַתּוֹרָה\hdot
וַהֲרֵינִי עוֹשֶׂה לְשֵׁם יִחוּד קֻדְשָׁא בְרִיךְ הוּא וּשְׁכִינְתֵּיהּ בִּדְחִילוּ וּרְחִימוּ\hdot
וִיהִי נוֹעַם יְהֹוָה אֱלֹהֵינוּ עָלֵינוּ וּמַעֲשֵׂה יָדֵינוּ כּוֹנְנָה עָלֵינוּ וּמַעֲשֵׂה יָדֵינוּ כּוֹנְנֵהוּ׃
\end{tp}

}

\vspace{1em}
{
\parskip 2pt
\sederContents{י״א}{
  {זבחים/zebahim/14,
    מנחות/menahot/13,
    חולין//12,
    בכורות//9,
    ערכין//9,
    תמורה//7,
    כריתות//6,
    מעילה//6,
    תמיד//7,
    מדות//5,
    קנים//3}}{צ״א}


\vspace{4pt}
 { \centering
   \small
       {\larger  ויש} מסדרים קנים קודם מדות והכל עולה בקנה אחד}

}
\mainmatter
\thispagestyle{empty}


\tractatehead{מסכת זבחים}{zebahim}

\chead{יום יח}

\perek{כָּל}{הַזְּבָחִים שֶׁנִּזְבְּחוּ שֶׁלֹּא לִשְׁמָן כְּשֵׁרִים\hdot
אֶלָּא שֶׁלֹּא עָלוּ לַבְּעָלִים לְשֵׁם חוֹבָה\hdot
חוּץ מִן הַפֶּסַח וּמִן הַחַטָּאת\hdot
הַפֶּסַח בִּזְמַנּוֹ\hdot
וְהַחַטָּאת בְּכָל זְמָן\hdot
רִבִּי אֱלִיעֶזֶר אוֹמֵר אַף הָאָשָׁם\hdot
הַפֶּסַח בִּזְמַנּוֹ\hdot
וְהַחַטָּאת וְהָאָשָׁם בְּכָל זְמָן\hdot
אָמַר רִבִּי אֱלִיעֶזֶר הַחַטָּאת בָּאָה עַל חֵטְא\hdot
וְהָאָשָׁם בָּא עַל חֵטְא\hdot
מַה חַטָּאת פְּסוּלָה שֶׁלֹּא לִשְׁמָהּ\hdot
אַף הָאָשָׁם פָּסוּל שֶׁלֹּא לִשְׁמוֹ׃}

\halacha{יוֹסֵי בֶּן חוֹנִי אוֹמֵר הַנִּשְׁחָטִים לְשֵׁם פֶּסַח וּלְשֵׁם חַטָּאת פְּסוּלִין\hdot
שִׁמְעוֹן אֲחִי עֲזַרְיָה אוֹמֵר שְׁחָטָן לְשֵׁם גָּבוֹהַּ מֵהֶם כְּשֵׁרִין\hdot
לְשֵׁם נָמוּךְ מֵהֶם פְּסוּלִין\hdot
כֵּיצַד קָדְשֵׁי קֳדָשִׁים שֶׁשְּׁחָטָן לְשֵׁם קֳדָשִׁים קַלִּים\hdot
פּסוּלִין\hdot
קֳדָשִׁים קַלִּים שֶׁשְּׁחָטָן לְשֵׁם קָדְשֵׁי קֳדָשִׁים\hdot
כְּשֵׁרִין\hdot
הַבְּכוֹר וְהַמַּעֲשֵׂר שֶׁשְּׁחָטָן לְשֵׁם שְׁלָמִים כְּשֵׁרִין\hdot
וּשְׁלָמִים שֶׁשְּׁחָטָן לְשֵׁם בְּכוֹר וּלְשֵׁם מַעֲשֵׂר פְּסוּלִין׃}

\halacha{הַפֶּסַח שֶׁשְּׁחָטוֹ בְּשַׁחְרִית בְּאַרְבָּעָה עָשָׂר שֶׁלֹּא לִשְׁמוֹ\hdot
רִבִּי יְהוֹשֻׁעַ מַכְשִׁיר כְּאִלּוּ נִשְׁחַט בִּשְׁלֹשָׁה עָשָׂר\hdot
בֶּן בְּתֵירָא פּוֹסֵל כְּאִלּוּ נִשְׁחַט בֵּין הָאַרְבָּעִים\hdot
אָמַר שִׁמְעוֹן בֶּן עֲזָאי מְקֻבָּל אֲנִי מִפִּי שִׁבְעִים וּשְׁנַיִם זְקֵנִים\hdot
בְּיוֹם שֶׁהוֹשִׁיבוּ רִבִּי אֶלְעָזָר בֶּן עֲזַרְיָה בִּישִׁיבָה שֶׁכָּל הַזְּבָחִים הַנֶּאֱכָלִים שֶׁנִּזְבְּחוּ שֶׁלֹּא לִשְׁמָן כְּשֵׁרִין\hdot
אֶלָּא שֶׁלֹּא עָלוּ לַבְּעָלִים לְשֵׁם חוֹבָה\hdot
חוּץ מִן הַפֶּסַח וּמִן הַחַטָּאת\hdot
וְלֹא הוֹסִיף בֶּן עֲזָאי אֶלָּא הָעוֹלָה\hdot
וְלֹא הוֹדוּ לוֹ חֲכָמִים׃}

\halacha{הַפֶּסַח וְהַחַטָּאת שֶׁשְּׁחָטָן שֶׁלֹּא לִשְׁמָן\hdot
קִבֵּל וְהָלַךְ וְזָרַק שֶׁלֹּא לִשְׁמָן\hdot
אוֹ לִשְׁמָן וְשֶׁלֹּא לִשְׁמָן\hdot
אוֹ שֶׁלֹּא לִשְׁמָן וְלִשְׁמָן\hdot
פְּסוּלִין\hdot
כֵּיצַד לִשְׁמָן וְשֶׁלֹּא לִשְׁמָן\hdot
לְשֵׁם פֶּסַח וּלְשֵׁם שְׁלָמִים\hdot
שֶׁלֹּא לִשְׁמָן וְלִשְׁמָן\hdot
לְשֵׁם שְׁלָמִים וּלְשֵׁם פֶּסַח\hdot
שֶׁהַזֶּבַח נִפְסָל בְּאַרְבָּעָה דְבָרִים\hdot
בִּשְׁחִיטָה\hdot
וּבְקִבּוּל\hdot
וּבְהִלּוּךְ\hdot
וּבִזְרִיקָה\hdot
רִבִּי שִׁמְעוֹן מַכְשִׁיר בְּהִלּוּךְ\hdot
שֶׁהָיָה רִבִּי שִׁמְעוֹן אוֹמֵר אֵי אֶפְשָׁר שֶׁלֹּא בִשְׁחִיטָה\hdot
וְשֶׁלֹּא בְקַבָּלָה\hdot
וְשֶׁלֹּא בִזְרִיקָה\hdot
אֲבָל אֶפְשָׁר שֶׁלֹּא בְהִלּוּךְ\hdot
שֶׁשּׁוֹחֵט בְּצַד הַמִּזְבֵּחַ וְזוֹרֵק\hdot
רִבִּי אֱלִיעֶזֶר אוֹמֵר הַמְּהַלֵּךְ בְּמָקוֹם שֶׁהוּא צָרִיךְ לְהַלֵּךְ הַמַּחֲשָׁבָה פוֹסֶלֶת\hdot
וּבִמְקוֹם שֶׁאֵינוֹ צָרְיךְ לְהַלֵּךְ\hdot
אֵין הַמַּחֲשָׁבָה פוֹסֶלֶת׃}

\perek{כָּל}{הַזְּבָחִים שֶׁקִּבֵּל דָּמָן זָר\hdot
אוֹנֵן טְבוּל יוֹם\hdot
מְחֻסַּר בְּגָדִים\hdot
מְחֻסַּר כִּפּוּרִים\hdot
שֶׁלֹּא רָחַץ יָדַיִם וְרַגְלַיִם\hdot
עָרֵל\hdot
טָמֵא\hdot
יוֹשֵׁב\hdot
עוֹמֵד עַל גַּבֵּי כֵלִים\hdot
עַל גּבֵּי בְהֵמָה\hdot
עַל גַּבֵּי חֲבֵרוֹ פָּסַל\hdot
קִבֵּל בַּשְּׂמֹאל פָּסַל\hdot
רִבִּי שִׁמְעוֹן מַכְשִׁיר\hdot
נִשְׁפַּךְ עַל הָרִצְפָה וַאֲסָפוֹ פָּסוּל\hdot
נְתָנוֹ עַל גַּבֵּי הַכֶּבֶשׁ שֶׁלֹּא כְנֶגֶד הַיְּסוֹד\hdot
נָתַן אֶת הַנִּתָּנִין לְמַטָּן לְמַעְלָן\hdot
וְאֶת הַנִּתָּנִין לְמַעְלָן לְמַטָּן\hdot
אֶת הַנִּתָּנִים בִּפְנִים בַּחוץ\hdot
וְאֶת הַנִּתָּנִים בַּחוּץ בִּפְנִים\hdot
פָּסוּל וְאֵין בּוֹ כָרֵת׃}

\halacha{הַשּׁוֹחֵט אֶת הַזֶּבַח לִזְרוֹק דָּמוֹ בַחוּץ\hdot
לְהַקְטִיר אֶת אֵמוּרָיו בַּחוּץ\hdot
אוֹ מִקְצַת אֵמוּרָיו בַּחוּץ\hdot
לֶאֱכוֹל בְּשָׂרוֹ בַחוּץ\hdot
אוֹ כַּזַּיִת מִבְּשָׂרוֹ בַחוּץ\hdot
אוֹ לֶאֱכוֹל כַּזַּיִת מֵעוֹר הָאַלְיָה בַחוּץ\hdot
פָּסוּל וְאֵין בּוֹ כָרֵת\hdot
לִזְרוֹק דָּמוֹ לְמָחָר\hdot
אוֹ מִקְצַת דָּמוֹ לְמָחָר\hdot
לְהַקְטִיר אֵמוּרָיו לְמָחָר\hdot
אוֹ מִקְצַת אֵמוּרָיו לְמָחָר\hdot
לֶאֱכוֹל בְּשָׂרוֹ לְמָחָר\hdot
אוֹ כַּזַּיִת מִבְּשָׂרוֹ לְמָחָר\hdot
אוֹ לֶאֱכוֹל כַּזַּיִת מֵעוֹר הָאַלְיָה לְמָחָר\hdot
פִּגּוּל וְחַיָּיבִין עָלָיו כָּרֵת׃}

\halacha{זֶה הַכְּלָל\hdot
כָּל הַשּׁוֹחֵט וְהַמְּקַבֵּל וְהַמְּהַלֵּךְ וְהַזּוֹרֵק\hdot
לֶאֱכוֹל דָּבָר שֶׁדַּרְכּוֹ לֶאֱכוֹל\hdot
לְהַקְטִיר דָּבָר שֶׁדַּרְכּוֹ לְהַקְטִיר\hdot
חוּץ לִמְקוֹמוֹ\hdot
פָּסוּל וְאֵין בּוֹ כָרֵת\hdot
חוּץ לִזְמַנּוֹ פִּגּוּל וְחַיָּבִין עָלָיו כָּרֵת\hdot
וּבִלְבָד שֶׁיִּקְרַב הַמַּתִּיר כְּמִצְוָתוֹ׃}

\halacha{כֵּיצַד קָרַב הַמַּתִּיר כְּמִצְוָתוֹ\hdot
שָׁחַט בִּשְׁתִּיקָה\hdot
קִבֵּל וְהָלַךְ וְזָרַק חוּץ לִזְמַנּוֹ\hdot
אוֹ שֶׁשָּׁחַט חוּץ לִזְמַנּוֹ\hdot
קִבֵּל וְהָלַךְ וְזָרַק בִּשְׁתִּיקָה\hdot
אוֹ שֶׁשָּׁחַט וְקִבֵּל וְהָלַךְ וְזָרַק חוּץ לִזְמַנּוֹ\hdot
זֶה הוּא שֶׁקָּרַב הַמַּתִּיר כְּמִצְוָתוֹ\hdot
כֵּיצַד לֹא קָרַב הַמַּתִּיר כְּמִצְוָתוֹ\hdot
שָׁחַט חוּץ לִמְקוֹמוֹ\hdot
קִבֵּל וְהָלַךְ וְזָרַק חוּץ לִזְמַנּוֹ\hdot
אוֹ שֶׁשָּׁחַט חוּץ לִזְמַנּוֹ\hdot
קִבֵּל וְהָלַךְ וְזָרַק חוּץ  לִמְקוֹמוֹ\hdot
אוֹ שֶׁשָּׁחַט וְקִבֵּל וְזָרַק וְהָלַךְ חוּץ לִמְקוֹמוֹ\hdot
הַפֶּסַח וְהַחַטָּאת שֶׁשְּׁחָטָן שֶׁלֹּא לִשְׁמָן\hdot
קִבֵּל וְהָלַךְ וְזָרַק חוּץ לִזְמַנּוֹ\hdot
אוֹ שֶׁשָּׁחַט חוּץ לִזְמַנּוֹ\hdot
קִבֵּל וְהָלַךְ וְזָרַק שֶׁלֹּא לִשְׁמָן\hdot
אוֹ שֶׁשָּׁחַט קִבֵּל וְהָלַךְ וְזָרַק שֶׁלֹּא לִשְׁמָן\hdot
זֶה הוּא שֶׁלֹּא קָרַב הַמַּתִּיר כְּמִצְוָתוֹ׃}

\halacha{לֶאֱכוֹל כַּזַּיִת בַּחוּץ וְכַזַּיִת לְמָחָר\hdot
כַּזַּיִת לְמָחָר וְכַזַּיִת בַּחוּץ\hdot
כַחֲצִי זַיִת בַּחוּץ וְכַחֲצִי זַיִת לְמָחָר\hdot
כַּחֲצִי זַיִת לְמָחָר וְכַחֲצִי זַיִת בַּחוּץ פָּסוּל וְאֵין בּוֹ כָרֵת\hdot
אָמַר רִבִּי יְהוּדָה זֶה הַכְּלָל אִם מַחֲשֶׁבֶת הַזְּמָן קָֽדְמָה לְמַחֲשֶׁבֶת הַמָּקוֹם\hdot
פִּגּוּל וְחַיָּבִין עָלָיו כָּרֵת\hdot
וְאִם מַחֲשֶׁבֶת הַמָּקוֹם קָֽדְמָה לְמַחֲשֶׁבֶת הַזְּמָן\hdot
פָּסוּל וְאֵין בּוֹ כָרֵת\hdot
וַחֲכָמִים אוֹמְרִים זֶה וְזֶה פָּסוּל וְאֵין בּוֹ כָרֵת\hdot
לֶאֱכוֹל כַּחֲצִי זַיִת וּלְהַקְטִיר כַּחֲצִי זַיִת כָּשֵׁר\hdot
שֶׁאֵין אֲכִילָה וְהַקְטָרָה מִצְטָֽרְפִין׃}

\perek{כָּל}{הַפְּסוּלִין שֶׁשָׁחֲטוּ\hdot
שְׁחִיטָתָן כְּשֵׁרָה\hdot
שֶׁהַשְּׁחִיטָה כְּשֵׁרָה בַּנַּשִׁים וּבַעֲבָדִים וּבַטְּמֵאִים אֲפִילוּ בְּקָדְשֵׁ קֳדָשִׁים\hdot
וּבִלְבַד שֶׁלֹּא יִהְיוּ טְמֵאִים נוֹגְעִים בַּבָּשָׂר\hdot
לְפִיכַךְ הֵם פּוֹסְלִים בְּמַחֲשָׁבָה\hdot
וְכֻלָּן שֶׁקִּבְּלוּ אֶת הַדָּם חוּץ לִזְמַנּוֹ\hdot
וְחוּץ לִמְקוֹמוֹ\hdot
אִם יֵשׁ דַּם הַנֶּפֶשׁ\hdot
יַחֲזוֹר הַכָּשֵׁר וִיקַבֵּל׃}

\halacha{קִבֵּל הַכָּשֵׁר וְנָתַן לַפָּסוּל\hdot
יַחֲזִיר לַכָּשֵׁר\hdot
קִבֵּל בִּימִינוֹ וְנָתַן לִשְׂמֹאלוֹ\hdot
יַחֲזִיר לִימִינוֹ\hdot
קִבֵּל בִּכְלִי קֹדֶשׁ וְנָתַן בִּכְלִי חוֹל\hdot
יַחֲזִיר לִכְלִי קֹדֶשׁ\hdot
נִשְׁפַּךְ מִן הַכֶּלִי עַל הָרִצְפָה וַאֲסָפוֹ\hdot
כָּשֵׁר\hdot
נְתָנוֹ עַל גַּבֵּי הַכֶּבֶשׁ שֶׁלֹּא כְנֶגֶד הַיְּסוֹד\hdot
נָתַן אֶת הַנִּתָּנִין לְמַטָּה לְמַעְלָה\hdot
וְאֶת הַנִּתָּנִן לְמַעְלָה לְמַטָּה\hdot
אֶת הַנִּתָּנִין בִּפְנִים בַּחוּץ\hdot
וְאֶת הַנִּתָּנִין בַּחוּץ בִּפְנִים\hdot
אִם יֵשׁ דַּם הַנֶּפֶשׁ\hdot
יַחֲזוֹר הַכָּשֵׁר וִיקַבֵּל׃}

\halacha{הַשּׁוֹחֵט אֶת הַזֶּבַח\hdot
לֶאֱכוֹל דָּבָר שֶׁאֵין דַּרְכּוֹ לֶאֱכוֹל\hdot
וּלְהַקְטִיר דָּבָר שֶׁאֵין דַּרְכּוֹ לְהַקְטִיר\hdot
כָּשֵׁר\hdot
רִבִּי אֱלִיעֶזֶר פּוֹסֵל\hdot
לֶאֱכוֹל דָּבָר שֶׁדַּרְכּוֹ לֶאֱכוֹל\hdot
וּלְהַקְטִיר דָּבָר שֶׁדַּרְכּוֹ לְהַקְטִיר\hdot
פָּחוּת מִכַּזַּיִת כָּשֵׁר\hdot
לֶאֱכוֹל כַּחֲצִי זַיִּת\hdot
וּלְהַקְטִיר כַּחֲצִי זַיִת כָּשֵׁר\hdot
שֶׁאֵין אֲכִילָה וְהַקְטָרָה מִצְטָֽרְפִין׃}

\halacha{הַשּׁוֹחֵט אֶת הַזֶּבַח לֶאֱכוֹל כַּזַּיִת מִן הָעוֹר\hdot
מִן הָרֹטֶב\hdot
מִן הַקִּפָּה\hdot
מִן הָאָלָל\hdot
מִן הָעֲצָמוֹת\hdot
מִן הַגִּידִים\hdot
מִן הַטְּלָפַיִם\hdot
מִן הַקְּרָנַיִם\hdot
חוּץ לִזְמַנּוֹ\hdot
אוֹ חוּץ לִמְקוֹמוֹ כָּשֵׁר\hdot
וְאֵין חַיָּבִין עֲלֵיהֶם\hdot
מִשּׁוּם פִּגּוּל וְנוֹתָר וְטָמֵא׃}

\halacha{הַשּׁוֹחֵט אֶת הַמֻּקְדָּשִׁין\hdot
לֶאֱכוֹל שָׁלִיל אוֹ שִׁלְיָא בַחוּץ\hdot
לֹא פִגֵּל\hdot
הַמּוֹלֵק תּוֹרִים בִּפְנִים\hdot
לֶאֱכוֹל בֵּיצֵיהֶם בַּחוּץ\hdot
לֹא פִגֵּל\hdot
חֲלֵב הַמֻּקְדָּשִׁין וּבֵיצֵי תּוֹרִים\hdot
אֵין חַיָּבִין עֲלֵיהֶן מִשּׁוּם פִּגּוּל וְנוֹתָר וְטָמֵא׃}

\halacha{שְׁחָטוֹ עַל מְנַת לְהַנִּיחַ אֶת דָּמוֹ וְאֶת אֵמוּרָיו לְמָחָר\hdot
אוֹ לְהוֹצִיאָן לַחוּץ\hdot
רִבִּי יְהוּדָה פוֹסֵל\hdot
וַחֲכָמִים מַכְשִׁירִין\hdot
שְׁחָטוֹ עַל מְנַת לִתְּנוֹ עַל גַּבֵּי הַכֶּבֶשׁ שֶׁלֹּא כְנֶגֶד הַיְּסוֹד\hdot
לִתֵּן אֶת הַנִּתָּנִין לְמַטָּן לְמַעְלָן\hdot
וְאֶת הַנִּתָּנִין לְמַעְלָן לְמַטָּן\hdot
וְאֶת הַנִּתָּנִין בִּפְנִים בַּחוּץ\hdot
וְאֶת הַנִּתָּנִין בַּחוּץ בִּפְנִים\hdot
שֶׁיֹּאכְלוּהוּ טְמֵאִים\hdot
שֶׁיַּקְרִיבוּהוּ טְמֵאִים\hdot
שֶׁיֹּאכְלוּהוּ עֲרֵלִים\hdot
שֶׁיַּקְרִיבוּהוּ עֲרֵלִים\hdot
לְשַׁבֵּר עַצְמוֹת הַפֶּסַח\hdot
וְלֶאֱכוֹל מִמֶּנּוּ נָא\hdot
לְעָרֵב דָּמוֹ בְּדַם פְּסוּלִין כָּשֵׁר\hdot
שֶׁאֵין הַמַּחֲשָׁבָה פוֹסֶלֶת\hdot
אֶלָּא חוּץ לִזְמַנּוֹ וְחוּץ לִמְקוֹמוֹ\hdot
וְהַפֶּסַח וְהַחַטָּאת שֶׁלֹּא לִשְׁמָן׃}

\perek{בֵּית}{שַׁמָּאי אוֹמְרִים כָּל הַנִּתָּנִין עַל מִזְבֵּחַ הַחִיצוֹן\hdot
שֶׁאִם נְתָנָן מַתָּנָה אַחַת כִּפֵּר\hdot
וּבְחַטָּאת שְׁתֵּי מַתָּנוֹת\hdot
וּבֵית הִלֵּל אוֹמְרִים אַף חַטָּאת שֶׁנִּתְּנָה מַתָּנָה אַחַת כִּפֵּר\hdot
לְפִיכַךְ אִם נָתַן הָרִאשׁוֹנָה כְּתִקּוּנָהּ\hdot
וְאֶת הַשְּׁנִיָּה חוּץ לִזְמַנָּהּ כִּפֵּר\hdot
נָתַן אֶת הָרִאשׁוֹנָה חוּץ לִזְמַנָּהּ\hdot
וְאֶת הַשְּׁנִיָּה חוּץ לִמְקוֹמָהּ\hdot
פּגּוּל וְחַיָּבִין עָלָיו כָּרֵת׃}

\halacha{כָּל הַנִּתָּנִין עַל מִזְבֵּחַ הַפְּנִימִי\hdot
שֶׁאִם חִסֵּר אַחַת מִן הַמַּתָּנוֹת לֹא כִפֵּר\hdot
לְפִיכַךְ אִם נָתַן כֻּלָּן כְּתִקּוּנָן וְאַחַת שֶׁלֹּא כְתִקּוּנָהּ\hdot
פָּסוּל וְאֵין בּוֹ כָּרֵת׃}

\halacha{אֵלּוּ דְבָרִים שֶׁאֵין חַיָּבִין עֲלֵיהֶם מִשּׁוּם פִּגּוּל\hdot
הַקֹּמֶץ\hdot
וְהַלְּבוֹנָה\hdot
וְהַקְּטֹרֶת\hdot
וּמִנְחַת כֹּהֲנִים\hdot
וּמִנְחַת כֹּהֵן מָשִׁיחַ\hdot
(וּמִנְחַת נְסָכִים) וְהַדָּם\hdot
וְהַנְּסָכִים הַבָּאִים בִּפְנֵי עַצְמָן\hdot
דִּבְרֵי רִבִּי מֵאִיר\hdot
וַחֲכָמִים אוֹמְרִים אַף הַבָּאִים עִם הַבְּהֵמָה\hdot
לוֹג שֶׁמֶן שֶׁל מְצוֹרָע\hdot
רִבִּי שִׁמְעוֹן אוֹמֵר אֵין חַיָּבִים עָלָיו מִשּׁוּם פִּגּוּל\hdot
וְרִבִּי מֵאִיר אוֹמֵר חַיָּבִין עָלָיו מִשּׁוּם פִּגּוּל\hdot
שֶׁדַּם הָאָשָׁם מַתִּירוֹ\hdot
וְכָל שֶׁיֵּשׁ לוֹ מַתִּירִים\hdot
בֵּין לָאָדָם בֵּין לַמִּזְבֵּחַ\hdot
חַיָּבִין עָלָין מִשּׁוּם פִּגּוּל׃}

\halacha{הָעוֹלָה דָּמָהּ מַתִּיר אֶת בְּשָׂרָהּ לַמִּזְבֵּחַ\hdot
וְעוֹרָהּ לַכֹּהֲנִים\hdot
עוֹלַת הָעוֹף דָּמָהּ מַתִּיר אֶת בְּשָׂרָהּ לַמִּזְבֵּחַ\hdot
חַטָּאת הָעוֹף דָּמָהּ מַתִּיר אֶת בְּשָׂרָהּ לַכֹּהֲנִים\hdot
פָּרִים הַנִּשְׂרָפִים וּשְׂעִירִים הַנִּשְׂרָפִים\hdot
דָּמָן מַתִּיר אֶת אֵימוּרֵיהֶן לִיקָרֵב\hdot
רִבִּי שִׁמְעוֹן אוֹמֵר כָּל שֶׁאֵינוֹ עַל מִזְבֵּחַ הַחִיצוֹן כַּשְּׁלָמִים\hdot
אֵין חַיָּבִין עָלָיו מִשּׁוּם פִּגּוּל׃}

\halacha{קָדְשֵׁי עוֹבְדֵי אֱלִילִים\hdot
אֵין חַיָּבִין עֲלֵיהֶם מִשּׁוּם פִּגּוּל נוֹתָר וְטָמֵא\hdot
וְהַשּׁוֹחֲטָן בַּחוּץ פָּטוּר\hdot
דִּבְרֵי רִבִּי מֵאִיר רִבִּי יוֹסֵי מְחַיֵּב\hdot
דְּבָרִים שֶׁאֵין חַיָּבִין עֲלֵיהֶם מִשּׁוּם פִּגּוּל\hdot
חַיָּבִים עֲלֵיהֶם מִשּׁוּם נוֹתָר\hdot
מִשּׁוּם טָמֵא\hdot
חוּץ מִן הַדָּם\hdot
רִבִּי שִׁמְעוֹן אוֹמֵר בְּדָבָר שֶׁדַּרְכָּן לְהֵאָכֵל\hdot
אֲבָל כְּגוֹן הָעֵצִים וְהַלְּבוֹנָה וְהַקְּטֹרֶת אֵין חַיָּבִין עֲלֵיהֶן מִשּׁוּם טֻמְאָה׃}

\halacha{לְשֵׁם שֵׁשָּׁה דְבָרִים הַזֶּבַח נִזְבַּח\hdot
לְשֵׁם זֶבַח\hdot
לְשֵׁם זוֹבֵחַ\hdot
לְשֵׁם הַשֵּׁם\hdot
לְשֵׁם אִישִׁים\hdot
לְשֵׁם רֵיחַ\hdot
לְשֵׁם נִיחֹחַ\hdot
וְהַחַטָּאת וְהָאָשָׁם לְשֵׁם חֵטְא\hdot
אָמַר רִבִּי יוֹסֵי אַף מִי שֶׁלֹּא הָיָה בְּלִבּוֹ לְשֵׁם אַחַד מִכָּל אֵלּוּ כָּשֵׁר\hdot
שֶׁהוּא תְּנָאי בֵּית דִּין\hdot
שֶׁאֵין הַמַּחֲשָׁבָה הוֹלֶכֶת אֶלָּא אַחַר הָעוֹבֵד׃}

\perek{אֵיזֶהוּ}{מְקוֹמָן שֶׁל זְבָחִים\hdot
קָדְשֵׁי קָדָשִׁים שְׁחִיטָתָן בַּצָּפוֹן\hdot
פַּר וְשָׂעִיר שֶׁל יוֹם הַכִּפּוּרִים\hdot
שְׁחִיטָתָן בַּצָּפוֹן\hdot
וְקִבּוּל דָּמָן בִּכְלֵי שָׁרֵת בַּצָּפוֹן\hdot
וְדָמָן טָעוּן הֲזָיָה עַל בֵּין הַבַּדִּים\hdot
וְעַל הַפָּרֹכֶת\hdot
וְעַל מִזְבַּח הַזָּהָב\hdot
מַתָּנָה אַחַת מֵהֶן מְעַכֶּבֶת\hdot
שִׁיְּרֵי הַדָּם הָיָה שׁוֹפֵךְ עַל יְסוֹד מַעֲרָבִי שֶׁל מִזְבֵּחַ הַחִיצוֹן\hdot
אִם לֹא נָתַן לֹא עִכֵּב׃}

\halacha{פָּרִים הַנִּשְׂרָפִים וּשְׂעִירִים הַנִּשְׂרָפִים שְׁחִיטָתָן בַּצָּפוֹן\hdot
וְקִבּוּל דָּמָן בִּכְלֵי שָׁרֵת בַּצָּפוֹן\hdot
וְדָמָן טָעוּן הֲזָיָה עַל הַפָּרֹכֶת\hdot
וְעַל מִזְבַּח הַזָּהָב\hdot
מַתָּנָה אַחַת מֵהֶן מְעַכֶּבֶת\hdot
שִׁיְּרֵי הַדָּם הָיָה שׁוֹפֵךְ עַל יְסוֹד מַעֲרָבִי שֶׁל מִזְבֵּחַ הַחִיצוֹן\hdot
אִם לֹא נָתַן לֹא עִכֵּב\hdot
אֵלּוּ וָאֵלּוּ נִשְׂרָפִין בְּבֵית הַדֶּשֶׁן׃}

\halacha{חַטּאוֹת הַצִּבּוּר וְהַיָּחִיד\hdot
אֵלּוּ הֵן חַטּאוֹת הַצִּבּוּר\hdot
שְׂעִירֵי רָאשֵׁי חֳדָשִׁים\hdot
וְשֶׁל מוֹעֲדוֹת\hdot
שְׁחִיטָתָן בַּצָּפוֹן\hdot
וְקִבּוּל דָּמָן בִּכְלֵי שָׁרֵת בַּצָּפוֹן\hdot
וְדָמָן טָעוּן אַרְבַּע מַתָּנוֹת עַל אַרְבַּע קְרָנוֹת\hdot
כֵּיצַד עָלָה בַּכֶּבֶשׁ וּפָנָה לַסּוֹבֵב\hdot
וּבָא לוֹ לְקֶרֶן דְּרוֹמִית מִזְרָחִית\hdot
מִזְרָחִית צְפוֹנִית\hdot
צְפוֹנִית מַעֲרָבִית\hdot
מַעֲרָבִית דְּרוֹמִית\hdot
שִׁיְּרֵי הַדָּם הָיָה שׁוֹפֵךְ עַל יְסוֹד דְּרוֹמִית\hdot
וְנֶאֱכָלִין לִפְנִים מִן הַקְּלָעִים\hdot
לְזִכְרֵי כְהוּנָה\hdot
בְּכָל מַאֲכָל\hdot
לְיוֹם וָלַיְלָה עַד חֲצוֹת׃}

\halacha{הָעוֹלָה קֹדֶשׁ קָדָשִׁים\hdot
שְׁחִיטָתָהּ בַּצָּפוֹן\hdot
וְקִבּוּל דָּמָהּ בִּכְלֵי שָׁרֵת בַּצָּפוֹן\hdot
וְדָמָהּ טָעוּן שְׁתֵּי מַתָּנוֹת שֶׁהֵן אַרְבַּע\hdot
וּטְעוּנָה הֶפְשֵׁט וְנִתּוּחַ\hdot
וְכָלִיל לָאִשִּׁים׃}

\halacha{זִבְחֵי שַׁלְמֵי צִבּוּר וַאֲשָׁמוֹת\hdot
אֵלּוּ הֵן אֲשָׁמוֹת\hdot
אֲשַׁם גְּזֵילוֹת\hdot
אֲשַׁם מְעִילוֹת\hdot
אֲשַׁם שִׁפְחָה חֲרוּפָה\hdot
אֲשַׁם נָזִיר\hdot
אֲשַׁם מְצוֹרָע\hdot
אָשָׁם תָּלוּי\hdot
שְׁחִיטָתָן בַּצָּפוֹן\hdot
וְקִבּוּל דָּמָן בִּכְלֵי שָׁרֵת בַּצָּפוֹן\hdot
וְדָמָן טָעוּן שְׁתֵּי מַתָּנוֹת שֶׁהֵן אַרְבַּע\hdot
וְנֶאֱכָלִין לִפְנִים מִן הַקְּלָעִים\hdot
לְזִכְרֵי כְהוּנָה\hdot
בְּכָל מַאֲכָל\hdot
לְיוֹם וָלַיְלָה עַד חֲצוֹת׃}

\halacha{הַתּוֹדָה וְאֵיל נָזִיר קָדָשִׁים קַלִּים\hdot
שְׁחִיטָתָן בְּכָל מָקוֹם בַּעֲזָרָה\hdot
וְדָמָן טָעוּן שְׁתֵּי מַתָּנוֹת שֶׁהֵן אַרְבַּע\hdot
וְנֶאֱכָלִין בְּכָל הָעִיר\hdot
לְכָל אָדָם\hdot
בְּכָל מַאֲכָל\hdot
לְיוֹם וָלַיְלָה עַד חֲצוֹת\hdot
הַמּוּרָם מֵהֶם כַּיּוֹצֵא בָהֶם\hdot
אֶלָּא שֶׁהַמּוּרָם נֶאֱכָל לַכֹּהֲנִים\hdot
לִנְשֵׁיהֶם וְלִבְנֵיהֶם וּלְעַבְדֵיהֶם׃}

\halacha{שְׁלָמִים קָדָשִׁים קַלִּים\hdot
שְׁחִיטָתָן בְּכָל מָקוֹם בַּעֲזָרָה\hdot
וְדָמָן טָעוּן שְׁתֵּי מַתָּנוֹת שֶׁהֵן אַרְבַּע\hdot
וְנֶאֱכָלִין בְּכָל הָעִיר\hdot
לְכָל אָדָם\hdot
בְּכָל מַאֲכָל\hdot
לִשְׁנֵי יָמִים וְלַיְלָה אֶחָד\hdot
הַמּוּרָם מֵהֶם כַּיּוֹצֵא בָהֶם\hdot
אֶלָּא שֶׁהַמּוּרָם נֶאֱכָל לַכֹּהֲנִים\hdot
לִנְשֵׁיהֶם וְלִבְנֵיהֶם וּלְעַבְדֵיהֶם׃}

\halacha{הַבְּכוֹר וְהַמַּעֲשֵׂר וְהַפֶּסַח קָדָשִׁים קַלִּים\hdot
שְׁחִיטָתָן בְּכָל מָקוֹם בַּעֲזָרָה\hdot
וְדָמָן טָעוּן מַתָּנָה אֶחָת\hdot
וּבִלְבַד שֶׁיִּתֵּן כְּנֶגֶד הַיְּסוֹד\hdot
שִׁנָּה בַאֲכִילָתָן\hdot
הַבְּכוֹר נֶאֱכָל לַכֹּהֲנִים\hdot
וְהַמַּעֲשֵׂר לְכָל אָדָם\hdot
וְנֶאֱכָלִין בְּכָל הָעִיר\hdot
בְּכָל מַאֲכָל\hdot
לִשְׁנֵי יָמִים וְלַיְלָה אֶחָד\hdot
הַפֶּסַח אֵינוֹ נֶאֱכָל אֶלָּא בַלַּיְלָה\hdot
וְאֵינוֹ נֶאֱכָל אֶלָּא עַד חֲצוֹת\hdot
וְאֵינוֹ נֶאֱכָל אֶלָּא לִמְנוּיָו\hdot
וְאֵינוֹ נֶאֱכָל אֶלָּא צָלִי׃}

\perek{קָדְשֵׁי}{קָדָשִׁים שֶׁשְּׁחָטָן בְּרֹאשׁ הַמִּזְבֵּחַ\hdot
רִבִּי יוֹסֵי אוֹמֵר כְּאִלּוּ נִשְׁחֲטוּ בַצָּפוֹן\hdot
רִבִּי יוֹסֵי בַּר יְהוּדָה אוֹמֵר מֵחֲצִי הַמִּזְבֵּחַ וְלַצָּפוֹן כַּצָּפוֹן\hdot
מֵחֲצִי הַמִּזְבֵּחַ וְלַדָּרוֹם כַּדָּרוֹם\hdot
הַמְּנָחוֹת הָיוּ נִקְמָצוֹת בְּכָל מָקוֹם בָּעֲזָרָה\hdot
וְנֶאֱכָלִין לִפְנִים מִן הַקְּלָעִים\hdot
לְזִכְרֵי כְהוּנָה\hdot
בְּכָל מַאֲכָל\hdot
לְיוֹם וָלַיְלָה עַד חֲצוֹת׃}

\halacha{חַטָּאת הָעוֹף הָֽיְתָה נַעֲשֵׂית עַל קֶרֶן דְּרוֹמִית מַעֲרָבִית\hdot
בְּכָל מָקוֹם הָֽיְתָה כְשֵׁרָה\hdot
אֶלָּא שֶׁזֶּה הָיָה מְקוֹמָהּ\hdot
שְׁלֹשָׁה דְבָרִים הָֽיְתָה אוֹתָהּ הַקֶּרֶן מְשַׁמֶּשֶׁת מִלְּמַטָּן וּשְׁלֹשָׁה מִלְּמַעְלָן\hdot
מִלְּמַטָּן חַטָּאת הָעוֹף וְהַהַגָּשׁוֹת וְשִׁיְּרֵי הַדָּם\hdot
מִלְּמַעְלָן נִסּוּךְ הַמַּיִם וְהַיַּיִן\hdot
וְעוֹלַת הָעוֹף כְּשֶׁהִיא רַבָּה בַּמִּזְרָח׃}

\halacha{כָּל הָעוֹלִין לַמִּזְבֵּחַ עוֹלִין דֶּרֶךְ יָמִין\hdot
וּמַקִּיפִין וְיוֹרְדִין דֶּרֶךְ שְׂמֹאל\hdot
חוּץ מִן הָעוֹלֶה לִשְׁלֹשָׁה דְבָרִים אֵלּוּ\hdot
שֶׁהָיוּ עוֹלִים וְחוֹזְרִים לָעָקֵב׃}

\halacha{חַטַּאת הָעוֹף כֵּיצַד הָֽיְתָה נַעֲשֵׂית\hdot
הָיָה מוֹלֵק אֶת רֹאשָׁהּ\hdot
מִמּוּל עָרְפָּהּ\hdot
וְאֵינוֹ מַבְדִּיל\hdot
וּמַזֶּה מִדָּמָהּ\hdot
עַל קִיר הַמִּזְבֵּחַ\hdot
שִׁיְּרֵי הַדָּם הָיָה מִתְמַצֶּה עַל הַיְּסוֹד\hdot
אֵין לַמִּזְבֵּחַ אֶלָּא דָמָהּ\hdot
וְכֻלָּהּ לַכֹּהֲנִים׃}

\halacha{עוֹלַת הָעוֹף כֵּיצַד הָֽיְתָה נַעֲשֵׂית\hdot
עָלָה בַכֶּבֶשׁ וּפָנָה לַסּוֹבֵב\hdot
בָּא לוֹ לְקֶרֶן דְּרוֹמִית מִזְרָחִית\hdot
הָיָה מוֹלֵק אֶת רֹאשָׁהּ מִמּוּל עָרְפָּהּ וּמַבְדִּיל\hdot
וּמְמַצֶּה אֶת דָּמָהּ עַל קִיר הַמִּזְבֵּחַ\hdot
נָטַל אֶת הָרֹאשׁ\hdot
וְהִקִּיף בֵּית מְלִיקָתוֹ לַמִּזְבֵּחַ\hdot
וּסְפָגוֹ בַמֶּלַח\hdot
וּזְרָקוֹ עַל גַּבֵּי הָאִישִׁים\hdot
בָּא לוֹ לַגּוּף\hdot
וְהֵסִיר אֶת הַמֻּרְאָה וְאֶת הַנּוֹצָה\hdot
וְאֶת בְּנֵי מֵעַיִם הַיּוֹצְאִים עִמָּהּ\hdot
וְהִשְׁלִיכָן לְבֵית הַדֶּשֶׁן\hdot
שִׁסַּע וְלֹא הִבְדִּיל\hdot
וְאִם הִבְדִּיל כָּשֵׁר\hdot
וּסְפָגוֹ בַמֶּלַח\hdot
וּזְרָקוֹ עַל גַּבֵּי הָאִישִׁים׃}

\halacha{לֹא הֵסִיר לֹא אֶת הַמֻּרְאָה\hdot
וְלֹא אֶת הַנּוֹצָה\hdot
וְלֹא אֶת בְּנֵי מֵעַיִם הַיּוֹצְאִים עִמָּהּ\hdot
וְלֹא סְפָגוֹ בַמֶּלַח\hdot
כָּל שֶׁשִּׁנָּה בָהּ מֵאַחַר שֶׁמִּצָּה אֶת דָּמָהּ כְּשֵׁרָה\hdot
הִבְדִּיל בַּחַטָּאת\hdot
וְלֹא הִבְדִּיל בָּעוֹלָה\hdot
פָּסַל\hdot
מִצָּה דַם הָרֹאשׁ וְלֹא מִצָּה דַם הַגּוּף\hdot
פְּסוּלָה\hdot
מִצָּה דַם הַגּוּף\hdot
וְלֹא מִצָּה דַם הָרֹאשׁ\hdot
כְּשֵׁרָה׃}

\halacha{חַטַּאת הָעוֹף שֶׁמְּלָקָהּ שֶׁלֹּא לִשְׁמָהּ\hdot
מִצָּה דָּמָהּ שֶׁלֹּא לִשְׁמָהּ\hdot
אוֹ לִשְׁמָהּ וְשֶׁלֹּא לִשְׁמָהּ\hdot
אוֹ שֶׁלֹּא לִשְׁמָהּ וְלִשְׁמָהּ\hdot
פְּסוּלָה\hdot
עוֹלַת הָעוֹף כְּשֵׁרָה\hdot
וּבִלְבַד שֶׁלֹּא עָֽלְתָה לַבְּעָלִים\hdot
אֶחָד חַטָּאת הָעוֹף\hdot
וְאֶחָד עוֹלַת הָעוֹף\hdot
שֶׁמְּלָקָן וְשֶׁמִּצָה דָמָן\hdot
לֶאֱכוֹל דָּבָר שֶׁדַּרְכּוֹ לֶאֱכוֹל\hdot
לְהַקְטִיר דָּבָר שֶׁדַּרְכּוֹ לְהַקְטִיר\hdot
חוּץ לִמְקוֹמוֹ\hdot
פָּסוּל וְאֵין בּוֹ כָרֵת\hdot
חוּץ לִזְמַנּוֹ\hdot
פִּגּוּל וְחַיָּבִין עָלָיו כָּרֵת\hdot
וּבִלְבַד שֶׁיִּקָּרֵב הַמַּתִּיר כְּמִצְוָתוֹ\hdot
כֵּיצַד קָרַב הַמַּתִּיר כְּמִצְוָתוֹ\hdot
מָלַק בִּשְׁתִּיקָה\hdot
וּמִצָּה הַדָּם חוּץ לִזְמַנּוֹ\hdot
אוֹ שֶׁמָּלַק חוּץ לִזְמַנּוֹ\hdot
וּמִצָּה הַדָּם בִּשְׁתִּיקָה\hdot
אוֹ שֶׁמָּלַק וּמִצָּה הַדָּם חוּץ לִזְמַנּוֹ\hdot
זֶה הוּא שֶׁקָּרַב הַמַּתִּיר כְּמִצְוָתוֹ\hdot
כֵּיצַד לֹא קָרַב הַמַּתִּיר כְּמִצְוָתוֹ\hdot
מָלַק חוּץ לִמְקוֹמוֹ וּמִצָּה הַדָּם חוּץ לִזְמַנּוֹ\hdot
אוֹ שֶׁמָּלַק חוּץ לִזְמַנּוֹ\hdot
וּמִצָּה הַדָּם חוּץ לִמְקוֹמוֹ\hdot
(אוֹ שֶׁמָּלַק וּמִצָּה אֶת הַדָּם חוּץ לִמְקוֹמוֹ)
חַטַּאת הָעוֹף שֶׁמְּלָקָהּ שֶׁלֹּא לִשְׁמָהּ\hdot
וּמִצָּה דָמָהּ חוּץ לִזְמַנָּהּ\hdot
אוֹ שֶׁמְּלָקָהּ חוּץ לִזְמַנָּהּ\hdot
וּמִצָּה דָמָהּ שֶׁלֹּא לִשְׁמָהּ\hdot
אוֹ שֶׁמְּלָקָהּ\hdot
וּמִצָּה דָמָהּ שֶׁלֹּא לִשְׁמָהּ\hdot
זֶה הוּא שֶׁלֹּא קָרַב הַמַּתִּיר כְּמִצְוָתוֹ\hdot
לֶאֱכוֹל כַּזַּיִת בַּחוּץ\hdot
וְכַזַּיִת לְמָחָר\hdot
כַּזַּיִת לְמָחָר וְכַזַּיִת בַּחוּץ\hdot
כַּחֲצִי זַיִת בַּחוּץ\hdot
וְכַחֲצִי זַיִת לְמָחָר\hdot
כַּחֲצִי זַיִת לְמָחָר\hdot
וְכַחֲצִי זַיִת בַּחוּץ\hdot
פָּסוּל וְאֵין בּוֹ כָרֵת\hdot
אָמַר רִבִּי יְהוּדָה זֶה הַכְּלָל\hdot
אִם מַחֲשֶׁבֶת הַזְּמָן קָֽדְמָה לְמַחֲשֶׁבֶת הַמָּקוֹם\hdot
פִּגּוּל וְחַיָּבִין עָלָיו כָּרֵת\hdot
וְאִם מַחֲשֶׁבֶת הַמָּקוֹם קָֽדְמָה לְמַחֲשֶׁבֶת הַזְּמָן\hdot
פָּסוּל וְאֵין בּוֹ כָרֵת\hdot
וַחֲכָמִים אוֹמְרִים זֶה וְזֶה פָּסוּל וְאֵין בּוֹ כָרֵת\hdot
לֶאֱכוֹל כַּחֲצִי זַיִת וּלְהַקְטִיר כַּחֲצִי זַיִת כָּשֵׁר\hdot
שֶׁאֵין אֲכִילָה וְהַקְטָרָה מִצְטָֽרְפִין׃}

\perek{חַטַּאת}{הָעוֹף שֶׁעֲשָׂאָהּ לְמַטָּה\hdot
כְּמַעֲשֵׂה חַטָּאת לְשֵׁם חַטָּאת\hdot
כְּשֵׁרָה\hdot
כְּמַעֲשֵׂה חַטָּאת לְשֵׁם עוֹלָה\hdot
כְּמַעֲשֵׂה עוֹלָה לְשֵׁם חַטָּאת\hdot
כְּמַעֲשֵׂה עוֹלָה לְשֵׁם עוֹלָה\hdot
פְּסוּלָה\hdot
עֲשָׂאָהּ לְמַעְלָה כְּמַעֲשֵׂה כֻלָּם\hdot
פְּסוּלָה׃}

\halacha{עוֹלַת הָעוֹף שֶׁעֲשָׂאָהּ לְמַעְלָן\hdot
כְּמַעֲשֵׂה עוֹלָה לְשֵׁם עוֹלָה כְּשֵׁרָה\hdot
כְּמַעֲשֵׂה עוֹלָה לְשֵׁם חַטָּאת כְּשֵׁרָה\hdot
וּבִלְבַד שֶׁלֹּא עָֽלְתָה לַבְּעָלִים\hdot
כְּמַעֲשֵׂה חַטָּאת לְשֵׁם עוֹלָה\hdot
כְּמַעֲשֵׁה חַטָּאת לְשֵׁם חַטָּאת\hdot
פְּסוּלָה\hdot
עֲשָׂאָן לְמַטָּן כְּמַעֲשֵׂה כֻלָּן\hdot
פְּסוּלָה׃}

\halacha{וְכֻלָּן אֵינָן מְטַמְּאוֹת בְּבֵית הַבְּלִיעָה וּמוֹעֲלִין בָּהֵן\hdot
חוּץ מֵחַטַּאת הָעוֹף שֶׁעֲשָׂאָהּ לְמַטָּן\hdot
כְּמַעֲשֵׁה חַטָּאת לְשֵׁם חַטָּאת׃}

\halacha{עוֹלַת הָעוֹף שֶׁעֲשָׂאָהּ לְמַטָּן\hdot
כְּמַעֲשֵׂה חַטָּאת לְשֵׁם חַטָּאת\hdot
רִבִּי אֱלִיעֶזֶר אוֹמֵר מוֹעֲלִין בָּהּ\hdot
רִבִּי יְהוֹשֻׁעַ אוֹמֵר אֵין מוֹעֲלִין בָּהּ\hdot
אָמַר רִבִּי אֱלִיעֶזֶר מָה אִם חַטָּאת שֶׁאֵין מוֹעֲלִין בָּהּ לִשְׁמָהּ\hdot
כְּשֶׁשִּׁנָּה אֶת שְׁמָהּ מוֹעֲלִין בָּהּ\hdot
עוֹלָה שֶׁמּוֹעֲלִין בָּהּ לִשְׁמָהּ\hdot
כְּשֶׁשִּׁנָּה אֶת שְׁמָהּ אֵינוֹ דִין שֶׁיִּמְעֲלוּ בָהּ\hdot
אָמַר לוֹ רִבִּי יְהוֹשֻׁעַ לֹא\hdot
אִם אָמַרְתָּ בַחַטָּאת שֶׁשִּׁנָּה אֶת שְׁמָהּ לְשֵׁם עוֹלָה\hdot
שֶׁכֵּן שִׁנָּה (אֶת) שְׁמָהּ לְדָבָר שֶׁיֵּשׁ בּוֹ מְעִילָה\hdot
תֹּאמַר בָּעוֹלָה שֶׁשִּׁנָּה אֶת שְׁמָהּ לְשֵׁם חַטָּאת\hdot
שֶׁכֵּן שִׁנָּה אֶת שְׁמָהּ לְדָבָר שֶׁאֵין בּוֹ מְעִילָה\hdot
אָמַר לוֹ רִבִּי אֱלִיעֶזֶר וַהֲרֵי קָדְשֵׁי קָדָשִׁים שֶׁשְּׁחָטָן בַּדָּרוֹם\hdot
וּשְׁחָטָן לְשֵׁם קָדָשִׁים קַלִּים יוֹכִיחוּ\hdot
שֶׁכֵּן שִׁנָּה אֶת שְׁמָן לְדָבָר שֶׁאֵין בּוֹ מְעִילָה\hdot
וּמוֹעֲלִין בָּהֵן\hdot
אַף אַתָּה אַל תִּתְמַהּ עַל הָעוֹלָה\hdot
שֶׁאַף עַל פִּי שֶׁשִּׁנָּה אֶת שְׁמָהּ\hdot
לְדָבָר שֶׁאֵין בּוֹ מְעִילָה שֶׁיִּמְעֲלוּ בָהּ\hdot
אָמַר לוֹ רִבִּי יְהוֹשֻׁעַ לֹא\hdot
אִם אָמַרְתָּ בְּקָדְשֵׁי קָדָשִׁים שֶׁשְּׁחָטָן בַּדָּרוֹם\hdot
וּשְׁחָטָן לְשֵׁם קָדָשִׁים קַלִּים\hdot
שֶׁכֵּן שִׁנָּה אֶת שְׁמָן לְדָבָר שֶׁיֵּשׁ בּוֹ אָסוּר וְהֶתֵּר\hdot
תֹּאמַר בָּעוֹלָה\hdot
שֶׁשִּׁנָּה אֶת שְׁמָהּ בְּדָבָר שֶׁכֻּלּוֹ הֶתֵּר׃}

\halacha{מָלַק בַּשְּׂמֹאל אוֹ בַלַּיְלָה\hdot
שָׁחַט חוֹלִין בִּפְנִים\hdot
וְקָדָשִׁים בַּחוּץ אֵינָן מְטַמְּאִין בְּבֵית הַבְּלִיעָה\hdot
מָלַק בְּסַכִּין\hdot
מָלַק חוֹלִין בִּפְנִים\hdot
וְקָדָשִׁים בַּחוּץ\hdot
תּוֹרִים שֶׁלֹּא הִגִּיעַ זְמַנָּן\hdot
וּבְנֵי יוֹנָה שֶׁעָבַר זְמַנָּן\hdot
שֶׁיָּבֵשׁ גַּפָּהּ\hdot
וְשֶׁנִּסְמֵית עֵינָהּ\hdot
וְשֶׁנִּקְטְעָה רַגְלָהּ\hdot
מְטַמְּאָה בְּבֵית הַבְּלִיעָה\hdot
זֶה הַכְּלָל כָּל שֶׁהָיָה פִסּוּלָהּ בַּקֹּדֶשׁ\hdot
אֵינָהּ מְטַמְּאָה בְּבֵית הַבְּלִיעָה\hdot
לֹא הָיָה פִסּוּלָהּ בַּקֹּדֶשׁ\hdot
מְטַמְּאָה בְּבֵית הַבְּלִיעָה\hdot
וְכָל הַפִּסּוּלִים שֶׁמָּֽלְקוּ\hdot
מְלִיקָתָן פְּסוּלָה\hdot
וְאֵינָן מְטַמְּאוֹת בְּבֵית הַבְּלִיעָה׃}

\halacha{מָלַק וְנִמְצָא טְרֵפָה\hdot
רִבִּי מֵאִיר אוֹמֵר אֵינָהּ מְטַמְּאָה בְּבֵית הַבְּלִיעָה\hdot
רִבִּי יְהוּדָה אוֹמֵר מְטַמְּאָה בְּבֵית הַבְּלִיעָה\hdot
אָמַר רִבִּי מֵאִיר מָה אִם נִבְלַת בְּהֵמָה שֶׁהִיא מְטַמְּאָה בְּמַגַּע וּבְמַשָּׂא\hdot
שְׁחִיטָתָהּ מְטַהֶרֶת אֶת טְרֵפָתָהּ מְטֻּמְאָתָהּ\hdot
נִבְלַת הָעוֹף שֶׁאֵינָהּ מְטַמְּאָה בְּמַגַּע וּבְמַשָּׂא\hdot
אֵינוֹ דִין שֶׁתְּהֵא שְׁחִיטָתָהּ מְטַהֶרֶת אֶת טְרֵפָתָהּ מִטֻּמְאָתָהּ\hdot
מַה מָּצִינוּ בִּשְׁחִיטָתָהּ שֶׁהִיא מְכַשַׁרְתָּהּ בַּאֲכִילָה\hdot
וּמְטַהֶרֶת אֶת טְרֵפָתָהּ מִטֻּמְאָתָהּ\hdot
אַף מְלִיקָתָהּ שֶׁהִיא מְכַשַׁרְתָּהּ בַּאֲכִילָה\hdot
תְּטַהֵר אֶת טְרֵפָתָהּ מִטֻּמְאָתָהּ\hdot
רִבִּי יוֹסֵי אוֹמֵר דַּיָּהּ בְּנִבְלַת בְּהֵמָה\hdot
שְׁחִיטָתָהּ מְטַהַרָהּ\hdot
אֲבָל לֹא מְלִיקָתָהּ׃}

\perek{כָּל}{הַזְּבָחִים שֶׁנִּתְעָֽרְבוּ בְּחַטּאוֹת הַמֵּתוֹת אוֹ בְשׁוֹר הַנִּסְקָל\hdot
(אֲפִילוּ אֶחָד בְּרִבּוֹא\hdot)
יָמוּתוּ כֻלָּם\hdot
נִתְעָֽרְבוּ בְּשׁוֹר שֶׁנֶּעֱבְדָה בוֹ עֲבֵירָה\hdot
אוֹ שֶׁהֵמִית אֶת הָאָדָם\hdot
עַל פִּי עֵד אֶחָד\hdot
אוֹ עַל פִּי הַבְּעָלִים\hdot
בָּרוֹבֵעַ וּבַנִּרְבָּע וּבַמֻּקְצֶה וּבַנֶּעֱבָד וּבָאֶתְנָן וּבַמְּחִיר וּבַכִּלְאַיִם וּבַטְּרֵיפָה וּבְיוֹצֵא דֹפֶן\hdot
יִרְעוּ עַד שֶׁיִּסְתָּאֲבוּ\hdot
וְיִמָּֽכְרוּ וְיָבִיא בִּדְמֵי הַיָּפֶה שֶׁבָּהֶן מֵאוֹתוֹ הַמִּין\hdot
נִתְעָרַב בְּחוֹלִין תְּמִימִים\hdot
יִמָּֽכְרוּ הַחוֹלִין לְצָרְכֵי אוֹתוֹ הַמִּין׃}

\halacha{קָדָשִׁים בְּקָדָשִׁים מִין בְּמִינוֹ\hdot
זֶה יִקָּרֵב לְשֵׁם מִי שֶׁהוּא\hdot
וְזֶה יִקָּרֵב לְשֵׁם מִי שֶׁהוּא\hdot
קָדָשִׁים בְּקָדָשִׁים מִין בְּשְׁאֵינוֹ מִינוֹ\hdot
יִרְעוּ עַד שֶׁיִּסְתָּאֲבוּ וְיִמָּֽכְרוּ וְיָבִיא בִדְמֵי הַיָּפֶה שֶׁבָּהֶן מִמִּין זֶה\hdot
וּבִדְמֵי הַיָּפֶה שֶׁבָּהֶן מִמִּין זֶה\hdot
וְיַפְסִיד הַמּוֹתָר מִבֵּיתוֹ\hdot
נִתְעָֽרְבוּ בַבְּכוֹר וּבַמַּעֲשֵׂר\hdot
יִרְעוּ עַד שֶׁיִּסְתָּאֲבוּ\hdot
וְיֵאָֽכְלוּ כַּבְּכוֹר וְכַמַּעֲשֵׂר\hdot
הַכֹּל יְכוֹלִין לְהִתְעָרֵב\hdot
חוּץ מִן הַחַטָּאת וְהָאָשָׁם׃}

\halacha{אָשָׁם שֶׁנִּתְעָרַב בַּשְּׁלָמִים\hdot
(יִרְעוּ עַד שֶׁיִּסְתָּאֲבוּ\hdot)
רִבִּי שִׁמְעוֹן אוֹמֵר שְׁנֵיהֶם יִשָּׁחֲטוּ בַצָּפוֹן\hdot
וְיֵאָֽכְלוּ כַּחֲמוֹר שֶׁבָּהֶן\hdot
אָֽמְרוּ לוֹ אֵין מְבִיאִים קָדָשִׁים לְבֵית הַפִּסּוּל\hdot
נִתְעָֽרְבוּ חֲתִיכוֹת בַּחֲתִיכוֹת\hdot
קָדְשֵׁי קָדָשִׁים בָּקָּדָשִׁים קַלִּים\hdot
הַנֶּאֱכָלִין לְיוֹם אֶחָד בַּנֶּאֱכָלִים לִשְׁנֵי יָמִים\hdot
יֵאָֽכְלוּ כַּחֲמוֹר שֶׁבָּהֶן׃}

\halacha{אֵבָרֵי חַטָּאת שֶׁנִּתְעָֽרְבוּ בְּאֵבָרֵי עוֹלָה\hdot
רִבִּי אֱלִיעֶזֶר אוֹמֵר יִתֵּן לְמַעְלָן\hdot
וְרוֹאֶה אֲנִי אֶת בְּשַׂר הַחַטָּאת מִלְּמַעְלָן\hdot
כְּאִלּוּ הוּא עֵצִים\hdot
וַחֲכָמִים אוֹמְרִים תְּעֻבַּר צוּרָתָן\hdot
וְיֵצֵא לְבֵית הַשְּׂרֵפָה׃}

\halacha{אֵבָרִין בְּאֵבָרֵי בַעֲלֵי מוּמִין\hdot
רִבִּי אֱלִיעֶזֶר אוֹמֵר\hdot
אִם קָרֵב רֹאשׁ אֶחָד מֵהֶן\hdot
יִקָּֽרְבוּ כָל הָרָאשִׁין\hdot
כְּרָעוֹ שֶׁל אֶחָד מֵהֶן\hdot
יִקָּֽרְבוּ כָל הַכְּרָעַיִם\hdot
וַחֲכָמִים אוֹמְרִים אֲפִילוּ קָֽרְבוּ כֻלָּם\hdot
חוּץ מֵאֶחָד מֵהֶן\hdot
יֵצֵא לְבֵית הַשְּׂרֵיפָה׃}

\halacha{דָּם שֶׁנִּתְעָרַב בַּמַּיִם\hdot
אִם יֵשׁ בּוֹ מַרְאֵה דָם כָּשֵׁר\hdot
נִתְעָרַב בַּיַּיִן\hdot
רוֹאִין אוֹתוֹ כְּאִלּוּ הוּא מַיִם\hdot
נִתְעָרַב בְּדַם בְּהֵמָה אוֹ בְּדַם חַיָּה רוֹאִין אוֹתוֹ כְּאִלּוּ הוּא מַיִם\hdot
רִבִּי יְהוּדָה אוֹמֵר אֵין דָּם מְבַטֵּל דָּם׃}

\halacha{נִתְעָרַב בְּדַם פְּסוּלִין\hdot
יִשָּׁפֵךְ לָאַמָּה\hdot
בְּדַם תַּמְצִית יִשָּׁפֵךְ לָאַמָּה\hdot
רִבִּי אֱלִיעֶזֶר מַכְשִׁיר\hdot
אִם לֹא נִמְלַךְ וְנָתַן כָּשֵׁר׃}

\halacha{דַּם תְּמִימִים בְּדַם בַּעֲלֵי מוּמִין\hdot
יִשָּׁפֵךְ לָאַמָּה\hdot
כּוֹס בְּכוֹסוֹת\hdot
רִבִּי אֱלִיעֶזֶר אוֹמֵר קָרֵב כּוֹס אֶחָד\hdot
יִקָּֽרְבוּ כָל הַכּוֹסוֹת\hdot
וַחֲכָמִים אוֹמְרִים אֲפִילוּ קָרְבוּ כֻלָּם חוּץ מֵאַחַד מֵהֶן\hdot
יִשָּׁפֵךְ לָאַמָּה׃}

\halacha{הַנִּתָּנִין לְמַטָּה שֶׁנְּתְעָֽרְבוּ בַּנִּתָּנִין לְמַעְלָה\hdot
רִבִּי אֱלִיעֶזֶר אוֹמֵר יִתֵּן לְמַעְלָה\hdot
וְרוֹאֶה אֲנִי אֶת הַתַּחְתּוֹנִים לְמַעְלָה\hdot
כְּאִלּוּ הִן מַיִם\hdot
וְיַחֲזוֹר וְיִתֵּן לְמַטָּה\hdot
וַחֲכָמִים אוֹמְרִים יִשָּֽׁפְכוּ לָאַמָּה\hdot
וְאִם לֹא נִמְלַךְ וְנַתַן כָּשֵׁר׃}

\halacha{הַנִּתָּנִן מַתָּנָה אַחַת שֶׁנִּתְעָֽרְבוּ בַּנִּתָּנִין מַתָּנָה אֶחָת\hdot
יִנָּֽתְנוּ מַתָּנָה אֶחָת\hdot
מַתַּן אַרְבַּע בְּמַתַּן אַרְבַּע יִנָּֽתְנוּ מַתַּן אַרְבַּע\hdot
מַתַּן אַרְבַּע בְּמַתָּנָה אֶחָת\hdot
רִבְּי אֱלִיעֶזֶר אוֹמֵר יִנָּֽתְנוּ בְּמַתַּן אַרְבַּע\hdot
רִבִּי יְהוֹשֻׁעַ אוֹמֵר יִנָּֽתְנוּ בְּמַתָּנָה אֶחָת\hdot
אָמַר לוֹ רִבִּי אֱלִיעֶזֶר וַהֲרֵי הוּא עוֹבֵר עַל בַּל תִּגְרַע\hdot
אָמַר לוֹ רִבִּי יְהוֹשֻׁעַ וַהֲרֵי הוּא עוֹבֵר עַל בַּל תּוֹסִיף\hdot
אָמַר לוֹ רִבִּי אֱלִיעֶזֶר לֹא נֶאֱמַר בַּל תּוֹסִיף\hdot
אֶלָּא כְּשֶׁהוּא בְעַצְמוֹ\hdot
אָמַר לוֹ רִבִּי יְהוֹשֻׁעַ לֹא נֶאֱמַר בַּל תִּגְרַע\hdot
אֶלָּא כְּשֶׁהוּא בְעַצְמוֹ\hdot
וְעוֹד אָמַר לוֹ רִבִּי יְהוֹשֻׁעַ כְּשֶׁנָּתַתָּ עָבַרְתָּ עַל בַּל תּוֹסִיף\hdot
וְעָשִׂיתָ מַעֲשֶׂה בְּיָדֶךָ\hdot
וּכְשֶׁלֹּא נָתַתָּ עָבַרְתָּ עַל בַּל תִּגְרַע\hdot
וְלֹא עָשִׂיתָ מַעֲשֶׂה בְּיָדֶךָ׃}

\halacha{הַנִּתָּנִים בִּפְנִים שֶׁנִּתְעָֽרְבוּ בַּנִּתָּנִין בַּחוּץ\hdot
יִשָּׁפְכוּ לָאַמָּה\hdot
נָתַן בַּחוּץ\hdot
וְחָזַר וְנָתַן בִּפְנִים כָּשֵׁר\hdot
בִּפְנִים וְחָזַר וְנָתַן בַּחוּץ רִבִּי עֲקִיבָא פּוֹסֵל\hdot
וַחֲכָמִים מַכְשִׁירִין\hdot
שֶׁהָיָה רִבִּי עֲקִיבָא אוֹמֵר כָּל הַדָּמִים שֶׁנִּכְנְסוּ לְכַפֵּר בַּהֵיכָל פְּסוּלִין\hdot
וַחֲכָמִים אוֹמְרִים חַטָּאת בִּלְבַד\hdot
רִבִּי אֱלִיעֶזֶר אוֹמֵר אַף הָאָשָׁם\hdot
שֶׁנֶּאֱמַר כַּחַטָּאת וְכָאָשָׁם׃}

\halacha{חַטָּאת שֶׁקִּבֵּל דָּמָהּ בִּשְׁנֵי כוֹסוֹת\hdot
יָצָא אַחַד מֵהֶן לַחוּץ\hdot
הַפְּנִימִי כָּשֵׁר\hdot
נִכְנָס אַחַד מֵהֶן לִפְנִים\hdot
רִבִּי יוֹסֵי הַגְּלִילִי מַכְשִׁיר בַּחִיצוֹן\hdot
וַחֲכָמִים פּוֹסְלִין\hdot
אָמַר רִבִּי יוֹסֵי הַגּלִילִי מָה אִם בַּמָּקוֹם שֶׁהַמַּחֲשָׁבָה פּוֹסֶלֶת בַּחוּץ\hdot
לֹא עָשָׂה אֶת הַמְּשׁוֹאָר כַּיּוֹצֵא\hdot
מָקוֹם שֶׁאֵין הַמַּחֲשָׁבָה פּוֹסֶלֶת בִּפְנִים\hdot
אֵינוֹ דִין שֶׁלֹּא נַעֲשֶׂה אֶת הַמְּשׁוֹאָר כַּנִּכְנָס\hdot
נִכְנַס לְכַפֵּר אַף עַל פִּי שֶׁלֹּא כִפֵּר פָּסַל\hdot
דִּבְרֵי רִבִּי אֱלִיעֶזֶר\hdot
רִבִּי שִׁמְעוֹן אוֹמֵר עַד שֶׁיְּכַפֵּר\hdot
רִבִּי יְהוּדָה אוֹמֵר אִם הִכְנִיס שׁוֹגֵג כָּשֵׁר\hdot
כָּל הַדָּמִים הַפְּסוּלִין שֶׁנִּתְּנוּ עַל גַּבֵּי הַמִּזְבֵּחַ\hdot
לֹא הִרְצָה הַצִּיץ אֶלָּא עַל הַטָּמֵא\hdot
שֶׁהַצִּיץ מְרַצֶּה עַל הַטָּמֵא\hdot
וְאֵינוֹ מְרַצֶּה עַל הַיּוֹצֵא׃}

\perek{הַמִּזְבֵּחַ}{מְקַדֵּשׁ אֶת הָרָאוּי לוֹ\hdot
רִבִּי יְהוֹשֻׁעַ אוֹמֵר כָּל הָרָאוּי לָאִשִּׁים אִם עָלָה לֹא יֵרֵד\hdot
שֶׁנֶּאֱמַר הִיא הָעוֹלָה עַל מוֹקְדָה עַל הַמִּזְבֵּחַ\hdot
מָה עוֹלָה שֶׁהִיא רְאוּיָה לָאִשִּׁים\hdot
אִם עָֽלְתָה לֹא תֵרֵד\hdot
אַף כָּל דָּבָר שֶׁהוּא רָאוּי לָאִשִּׁים אִם עָלָה לֹא יֵרֵד\hdot
רַבָּן גַּמְלִיאֵל אוֹמֵר כָּל הָרָאוּי לַמִּזְבֵּחַ\hdot
אִם עָלָה לֹא יֵרֵד\hdot
שֶׁנֶּאֱמַר הִיא הָעוֹלָה עַל מוֹקְדָה עַל הַמִּזְבֵּחַ\hdot
מָה עוֹלָה שֶׁהִיא רְאוּיָה לַמִּזְבֵּחַ אִם עָֽלְתָה לֹא תֵרֵד\hdot
אַף כָּל דָּבָר שֶׁהוּא רָאוּי לַמִּזְבֵּחַ\hdot
אִם עָלָה לֹא יֵרֵד\hdot
אֵין בֵּין דִּבְרֵי רַבָּן גַּמְלִיאֵל לְדִבְרֵי רִבִּי יְהוֹשֻׁעַ אֶלָּא הַדָּם וְהַנְּסָכִים\hdot
שֶׁרַבָּן גַּמְלִיאֵל אוֹמֵר לֹא יֵרְדוּ\hdot
וְרִבִּי יְהוֹשֻׁעַ אוֹמֵר יֵרְדוּ\hdot
רִבִּי שִׁמְעוֹן אוֹמֵר הַזֶּבַח כָּשֵׁר וְהַנְּסָכִים פְּסוּלִין\hdot
הַנְּסָכִים כְּשֵׁרִין וְהַזֶּבַח פָּסוּל\hdot
אֲפִילוּ זֶה וְזֶה פְּסוּלִין\hdot
הַזֶּבַח לֹא יֵרֵד וְהַנְּסָכִים יֵרְדוּ׃}

\halacha{וְאֵלּוּ אִם עָלוּ לֹא יֵרְדוּ\hdot
הַלָּן\hdot
וְהַטָּמֵא\hdot
וְהַיּוֹצֵא\hdot
וְהַנִּשְׁחָט חוּץ לִזְמַנּוֹ\hdot
וְחוּץ לְמִקוֹמוֹ\hdot
וְשֶׁקִּבְּלוּ פְסוּלִים וְזָֽרְקוּ אֶת דָּמוֹ\hdot
רִבִּי יְהוּדָה אוֹמֵר שֶׁנִּשְׁחֲטָה בַלַּיְלָה\hdot
וְשֶׁנִּשְׁפַּךְ דָּמָהּ\hdot
וְשֶׁיָּצָא דָמָהּ חוּץ לַקְּלָעִים\hdot
אִם עָֽלְתָה תֵרֵד\hdot
רִבִּי שִׁמְעוֹן אוֹמֵר לֹא תֵרֵד\hdot
שֶׁהָיָה רִבִּי שִׁמְעוֹן אוֹמֵר כָּל שֶׁפִּסּוּלוֹ בַּקֹּדֶשׁ הַקֹּדֶשׁ מְקַבְּלוֹ\hdot
לֹא הָיָה פִסּוּלוֹ בַּקֹּדֶשׁ\hdot
אֵין הַקֹּדֶשׁ מְקַבְּלוֹ׃}

\halacha{אֵלּוּ לֹא הָיָה פִסּוּלָן בַּקֹּדֶשׁ\hdot
הָרוֹבֵעַ וְהַנִּרְבָּע וְהַמּוּקְצֶה וְהַנֶּעֱבָד\hdot
וְהָאֶתְנָן וְהַמְּחִיר וְהַכִּלְאַיִם וְהַטְּרֵפָה וְיוֹצֵא דֹפֶן\hdot
וּבַעֲלֵי מוּמִין\hdot
רִבִּי עֲקִיבָא מַכְשִׁיר בְּבַעֲלֵי מוּמִין\hdot
רִבִּי חֲנַנְיָא סְגַן הַכֹּהֲנִים אוֹמֵר דּוֹחֶה הָיָה אַבָּא אֶת בַּעֲלֵי מוּמִין\hdot
מֵעַל גַּבֵּי הַמִּזְבֵּחַ׃}

\halacha{כְּשֵׁם שֶׁאִם עָלוּ לֹא יֵרְדוּ\hdot
כַּךְ אִם יֵרְדוּ לֹא יַעֲלוּ\hdot
וְכֻלָּן שֶׁעָלוּ חַיִּים לְרֹאשׁ הַמִּזְבֵּחַ יֵרְדוּ\hdot
עוֹלָה שֶׁעָֽלְתָה חַיָּה לְרֹאשׁ הַמִּזְבֵּחַ תֵּרֵד\hdot
שְׁחָטָהּ בְּרֹאשׁ הַמִּזְבֵּחַ\hdot
יַפְשִׁיט וִינַתֵּחַ בִּמְקוֹמָהּ׃}

\halacha{וְאֵלּוּ אִם עָלוּ יֵרְדוּ\hdot
בְּשַׂר קָדְשֵׁי קָדָשִׁים\hdot
וּבְשַׂר קָדָשִׁים קַלִּים\hdot
וּמוֹתַר הָעֹמֶר\hdot
וּשְׁתֵּי הַלֶּחֶם\hdot
וְלֶחֶם הַפָּנִים\hdot
וְשִׁיְּרֵי מְנָחוֹת\hdot
וְהַקְּטֹרֶת\hdot
הַצֶּמֶר שֶׁבְּרָאשֵׁי כְבָשִׂים\hdot
וְהַשֵּׂעָר שְׁבִּזְקַן תַּיָּשִׁים\hdot
וְהָעֲצָמוֹת וְהַגִּידִים וְהַקְּרָנַיִם וְהָטְלָפַיִם בִּזְמָן שֶׁאֵין מְחֻבָּרִין יַעֲלוּ\hdot
שֶׁנֶּאֱמַר וְהִקְטִיר הַכֹּהֵן אֶת הַכֹּל הַמִּזְבֵּחָה\hdot
פֵּרְשׁוּ לֹא יַעֲלוּ\hdot
שֶׁנֶּאֱמַר וְעָשִׂיתָ עוֹלֹתֶיךָ הַבָּשָׂר וְהַדָּם׃}

\halacha{וְכֻלָּם שֶׁפָּֽקְעוּ מֵעַל גַּבֵּי הַמִּזְבֵּחַ לֹא יַחֲזִיר\hdot
וְכֵן גַּחֶלֶת שֶׁפָּֽקְעָה מֵעַל גַּבֵּי הַמִּזְבֵּחַ\hdot
אֵבָרִים שֶׁפָּֽקְעוּ מֵעַל גַּבֵּי הַמִּזְבֵּחַ\hdot
קֹדֶם לַחֲצוֹת\hdot
יַחֲזִיר וּמוֹעֲלִין בָּהֶן\hdot
לְאַחַר חֲצוֹת לֹא יַחֲזִיר\hdot
וְאֵין מוֹעֲלִין בָּהֶן׃}

\halacha{כְּשֵׁם שֶׁהַמִּזְבֵּחַ מְקַדֵּשׁ אֶת הָרָאוּי לוֹ\hdot
כַּךְ הַכֶּבֶשׁ מְקַדֵּשׁ\hdot
כְּשֵׁם שֶׁהַמִּזְבֵּחַ וְהַכֶּבֶשׁ\hdot
מְקַדְּשִׁין אֶת הָרָאוּי לָהֶן\hdot
כַּךְ הַכֵּלִים מְקַדְּשִׁין\hdot
כְּלֵי הַלַּח מְקַדְּשִׁין אֶת הַלַּח\hdot
וּמִדּוֹת הַיָּבֵשׁ מְקַדְּשׁוֹת אֶת הַיָּבֵשׁ\hdot
אֵין כְּלֵי הַלַּח מְקַדְּשִׁין אֶת הַיָּבֵשׁ\hdot
וְלֹא מִדּוֹת הַיָּבֵשׁ מְקַדְּשׁוֹת אֶת הַלַּח\hdot
כְּלֵי הַקֹּדֶשׁ שֶׁנִּקָּבוּ\hdot
אִם עוֹשִׂים הֵם מֵעֵין מְלַאכְתָּן שֶׁהָיוּ עוֹשִׂין וְהֵן שְׁלֵמִים\hdot
מְקַדִּשִׁים\hdot
וְאִם לָאוּ אֵין מְקַדְּשִׁים\hdot
וְכֻלָּם אֵין מְקַדְּשִׁים אֶלָּא בַקֹּדֶשׁ׃}

\perek{כָּל}{הַתָּדִיר מֵחֲבֵרוֹ קוֹדֵם אֶת חֲבֵרוֹ\hdot
הַתְּמִידִין קוֹדְמִין לַמּוּסָפִין\hdot
מוּסְפֵי שַׁבָּת קוֹדְמִין לְמוּסְפֵי רֹאשׁ חֹדֶשׁ\hdot
מוּסְפֵי רֹאשׁ חֹדֶשׁ קוֹדְמִין לְמוּסְפֵי רֹאשׁ הַשָּׁנָה\hdot
(שֶׁנֶּאֱמַר מִלְּבַד עוֹלַת הַבֹּקֶר אֲשֶׁר לְעוֹלַת הַתָּמִיד תַּעֲשׂוּ אֶת אֵלֶּה׃}

\halacha{וְכָל הַמְּקֻדָּשׁ מֵחֲבֵירוֹ קוֹדֵם אֶת חֲבֵירוֹ\hdot
דַּם חַטָּאת קוֹדֵם לְדַם עוֹלָה\hdot
מִפְּנֵי שֶׁהוּא מְרַצֶּה\hdot
אֵיבָרֵי עוֹלָה קוֹדְמִין לְאֵמוּרֵי חַטָּאת\hdot
מִפְּנֵי שֶׁהֵן כָּלִיל לָאִשִּׁים\hdot
חַטָּאת קוֹדֶמֶת לָאָשָׁם\hdot
מִפְּנֵי שֶׁדָּמָהּ נִתָּן עַל אַרְבַּע קְרָנוֹת\hdot
וְעַל הַיְּסוֹד\hdot
אָשָׁם קוֹדֵם לַתּוֹדָה וּלְאֵיל נָזִיר\hdot
מִפְּנֵי שֶׁהוּא קָדְשֵׁי קָדָשִׁים\hdot
הַתּוֹדָה וְאֵיל נָזִיר קוֹדְמִין לַשְּׁלָמִים\hdot
מִפְּנֵי שֶׁהֵן נֶאֱכָלִין לְיוֹם אֶחָד וּטְעוּנִים לֶחֶם\hdot
שְׁלָמִים קוֹדְמִין לַבְּכוֹר מִפְּנֵי שֶׁהֵן טְעוּנִים מַתָּן אַרְבַּע\hdot
וּסְמִיכָה וּנְסָכִים וּתְנוּפַת חָזֶה וְשׁוֹק׃}

\halacha{הַבְּכוֹר קוֹדֵם לַמַּעֲשֵׂר\hdot
מִפְּנֵי שֶׁקְּדוּשָׁתוֹ מֵרֶחֶם\hdot
וְנֶאֱכָל לַכֹּהֲנִים\hdot
הַמַּעֲשֵׂר קוֹדֵם לָעוֹפוֹת\hdot
מִפְּנֵי שֶׁהוּא זֶבַח\hdot
וְיֵשׁ בּוֹ קָדְשֵׁי קָדָשִׁים דָּמוֹ וְאֵמוּרָיו׃}

\halacha{הָעוֹפוֹת קוֹדְמִין לַמְּנָחוֹת\hdot
מִפְּנֵי שֶׁהֵן מִינֵי דָמִים\hdot
מִנְחַת חוֹטֵא קוֹדֶמֶת לְמִנְחַת נְדָבָה\hdot
מִפְּנֵי שֶׁהִיא בָאָה עַל חֵטְא\hdot
חַטַּאת הָעוֹף קוֹדֶמֶת לְעוֹלַת הָעוֹף\hdot
וְכֵן בְּהֶקְדֵשָׁהּ׃}

\halacha{כָּל הַחַטָּאוֹת שֶׁבַּתּוֹרָה קוֹדְמוֹת לַאֲשָׁמוֹת\hdot
חוּץ מֵאֲשַׁם מְצוֹרָע\hdot
מִפְּנֵי שֶׁהוּא בָא לְהֶכְשֵׁר\hdot
כָּל הָאֲשָׁמוֹת שֶׁבַּתּוֹרָה\hdot
בָּאִין בְּנֵי שְׁתַּיִם\hdot
וּבָאִין בְּכֶסֶף שְׁקָלִים\hdot
חוּץ מֵאֲשַׁם נָזִיר\hdot
וַאֲשַׁם מְצוֹרָע\hdot
שֶׁהֵן בָּאִין בְּנֵי שְׁנָתָן\hdot
וְאֵינָן בָּאִין בְּכֶסֶף שְׁקָלִים׃}

\halacha{כְּשֵׁם שֶׁהֵן קוֹדְמִין בְּהַקְרָבָתָן\hdot
כַּךְ הֵן קוֹדְמִין בַּאֲכִילָתָן\hdot
שְׁלָמִים שֶׁל אֶמֶשׁ וּשְׁלָמִים שֶׁל הַיּוֹם\hdot
שֶׁל אֶמֶשׁ קוֹדְמִין\hdot
שְׁלָמִים שֶׁל אֶמֶשׁ\hdot
וְחַטָּאת וְאָשָׁם שֶׁל הַיּוֹם\hdot
שֶׁל אֶמֶשׁ קוֹדְמִין\hdot
דִּבְרֵי רִבִּי מֵאִיר\hdot
וַחֲכָמִים אוֹמְרִים הַחַטָּאת קוֹדֶמֶת מִפְּנֵי שֶׁהִיא קָדְשֵׁי קָדָשִׁים׃}

\halacha{וּבְכֻלָּם הַכֹּהֲנִים רְשָׁאִין לְשַׁנּוֹת בַּאֲכִילָתָן\hdot
לְאָכְלָן צְלוּיִם שְׁלוּקִין וּמְבֻשָּׁלִין\hdot
וְלָתֵת לְתוֹכָן תְּבָלִין חוֹלִין וּתְבָלֵי תְרוּמָה\hdot
דִּבְרֵי רִבִּי שִׁמְעוֹן\hdot
רִבִּי מֵאִיר אוֹמֵר לֹא יִתֵּן לְתוֹכָן תְּבָלֵי תְרוּמָה\hdot
שֶׁלֹּא יָבִיא אֶת הַתְּרוּמָה לִידֵי פִסּוּל׃}

\halacha{אָמַר רִבִּי שִׁמְעוֹן אִם רָאִיתָ שֶׁמֶן שֶׁהוּא מִתְחַלֵּק בַּעֲזָרָה\hdot
אֵין אַתָּה צָרִיךְ לִשְׁאוֹל מַה הוּא\hdot
אֶלָּא מוֹתַר רְקִיקֵי מְנָחוֹת יִשְׂרָאֵל\hdot
וְלוֹג שֶׁמֶן שֶׁל מְצוֹרָע\hdot
אִם רָאִיתָ שֶׁמֶן שֶׁהוּא נָתוּן עַל גַּבֵּי הָאִשִּׁים\hdot
אֵין אַתָּה צָרִיךְ לִשְׁאוֹל מַה הוּא\hdot
אֶלָּא מוֹתַר רְקִיקֵי מְנָחוֹת כֹּהֲנִים\hdot
וּמִנְחַת כֹּהֵן הַמָּשִׁיחַ\hdot
שֶׁאֵין מִתְנַדְּבִים שֶׁמֶן\hdot
רִבִּי טַרְפוֹן אוֹמֵר מִתְנַדִּבִים שָׁמֶן׃}

\perek{דַּם}{חַטָּאת שׁנִּתַּז עַל הַבֶּגֶד\hdot
הֲרֵי זֶה טָעוּן כִּבּוּס\hdot
אַף עַל פִּי שֶׁאֵין הַכָּתוּב מְדַבֵּר\hdot
אֶלָּא בַּנֶּאֱכָלוֹת\hdot
שֶׁנֶּאֱמַר בְּמָקוֹם קָדוֹשׁ תֵּאָכֵל\hdot
אֶחָד הַנֶּאֱכֶלֶת וְאֶחָד הַפְּנִימִית טְעוּנוֹת כִּבּוּס\hdot
שֶׁנֶּאֱמַר תּוֹרַת הַחַטָּאת\hdot
תּוֹרָה אַחַת לְכָל הַחַטָּאוֹת׃}

\halacha{חַטָּאת פְּסוּלָה אֵין דָּמָהּ טָעוּן כִּבּוּס\hdot
בֵּין שֶׁהָיָה לָהּ שַׁעַת הֶכְשֵׁר\hdot
בֵּין שֶׁלֹּא הָיָה לָהּ שַׁעַת הֶכְשֵׁר\hdot
אֵיזוֹ הִיא שֶׁהָיָה לָהּ שַׁעַת הֶכְשֵׁר\hdot
שֶׁלָּנָה\hdot
וְשֶׁנִּטְמֵאת\hdot
וְשֶׁיָּֽצְאָה\hdot
וְאֵיזוֹ הִיא שֶׁלֹּא הָיָה לָהּ שַׁעַת הֶכְשֵׁר\hdot
שֶׁנִּשְׁחֲטָה חוּץ לִזְמָנָהּ\hdot
וְחוּץ לִמְקֹמָהּ\hdot
וְשֶׁקִּבְּלוּ פְסוּלִין\hdot
וְזָֽרְקוּ אָת דָּמָהּ׃}

\halacha{נִתַּז מִן הַצַּוָּאר עַל הַבֶּגֶד אֵינוֹ טָעוּן כִּבּוּס\hdot
מִן הַקֶּרֶן וּמִן הַיְּסוֹד\hdot
אֵינוֹ טָעוּן כִּבּוּס\hdot
נִשְׁפַּךְ עַל הָרִצְפָה וַאֲסָפוֹ\hdot
אֵינוֹ טָעוּן כִּבּוּס\hdot
(אֵינוֹ טָעוּן כִּבּוּס) אֶלָּא הַדָּם שֶׁנִּתְקַבַּל בַּכֶּלִי\hdot
וְרָאוּי לְהַזָּיָה\hdot
נִתַּז עַל הָעוֹר\hdot
עַד שֶׁלֹּא הֻפְשַׁט\hdot
אֵינוֹ טָעוּן כִּבּוּס\hdot
מִשֶּׁהֻפְשַׁט טָעוּן כִּבּוּס\hdot
דִּבְרֵי רִבִּי יְהוּדָה\hdot
רִבִּי אֱלִיעֶזֶר אוֹמֵר אַף מִשֶּׁהֻפְשַׁט אֵינוֹ טָעוּן כִּבּוּס\hdot
אֵינוֹ טָעוּן כִּבּוּס אֶלָּא מְקוֹם הַדָּם\hdot
וְדָבָר שֶׁהוּא רָאוּי לְקַבֵּל טֻמְאָה\hdot
וְרָאוּי לְכִבּוּס׃}

\halacha{אֶחָד הַבֶּגֶד\hdot
וְאֶחָד הַשָּׂק\hdot
וְאֶחָד הָעוֹר\hdot
טְעוּנִין כִּבּוּס בְּמָקוֹם קָדוֹשׁ\hdot
וּשְׁבִירַת כְּלִי חֶרֶס בְּמָקוֹם קָדוֹשׁ\hdot
וּמְרִיקָה וּשְׁטִיפָה בִּכְלִי נְחֹשֶׁת בְּמָקוֹם קָדוֹשׁ\hdot
זֶה חֹמֶר בַּחַטָּאת\hdot
מִקָּדְשֵׁי קָדָשִׁים׃}

\halacha{בֶּגֶד שֶׁיָּצָא חוּץ לַקְּלָעִים\hdot
נִכְנָס וּמְכַבְּסוֹ בְּמָקוֹם קָדוֹשׁ\hdot
נִטְמָא חוּץ לַקְּלָעִים\hdot
קוֹרְעוֹ וְנִכְנָס וּמְכַבְּסוֹ בְּמָקוֹם קָדוֹשׁ\hdot
כְּלִּי חֶרֶס שֶׁיָּצָא חוּץ לַקְּלָעִים\hdot
נִכְנָס וְשׁוֹבְרוֹ בְּמָקוֹם קָדוֹשׁ\hdot
נִטְמָא חוּץ לַקְּלָעִים\hdot
נוֹקְבוֹ וְנִכְנָס וּמְכַבְּסוֹ בְּמָקוֹם קָדוֹשׁ׃}

\halacha{כְּלִי נְחֹשֶׁת שֶׁיָּצָא חוּץ לַקְּלָעִים\hdot
נִכְנָס וּמוֹרְקוֹ וְשׁוֹטְפוֹ בְּמָקוֹם קָדוֹשׁ\hdot
נִטְמָא חוּץ לַקְּלָעִים\hdot
פּוֹחֲתוֹ\hdot
וְנִכְנָס וּמוֹרְקוֹ וְשׁוֹטְפוֹ בְּמָקוֹם קָדוֹשׁ׃}

\halacha{אֶחָד שֶׁבִּשֵּׁל בּוֹ\hdot
וְאֶחָד שֶׁעִרָה לְתוֹכוֹ רוֹתֵחַ\hdot
אֶחָד קָדְשֵׁי קָדָשִׁים\hdot
וְאֶחָד קָדָשִׁים קַלִּים\hdot
טְעוּנִין מְרִיקָה וּשְׁטִיפָה\hdot
רִבִּי שִׁמְעוֹן אוֹמֵר קָדָשִׁים קַלִּים אֵינָן טְעוּנִין מְרִיקָה וּשְׁטִיפָה\hdot
רִבִּי טַרְפוֹן אוֹמֵר אִם בִּשֵּׁל בּוֹ מִתְּחִלַּת הָרֶגֶל יְבַשֵּׁל בּוֹ אֶת כָּל הָרֶגֶל\hdot
וַחֲכָמִים אוֹמְרִים עַד זְמַן אֲכִילָה\hdot
מְרִיקָה וּשְׁטִיפָה\hdot
מְרִיקָה כִּמְרִיקַת הַכּוֹס\hdot
וּשְׁטִיפָה כִּשְׁטִיפַת הַכּוֹס\hdot
(מְרִיקָה בַחַמִּין\hdot
וּשְׁטִיפָה בְצוֹנֵן)\hdot
וְהַשְּׁפוֹד וְהָאַסְכָּלָה מַגְעִילָן בַּחַמִּין׃}

\halacha{בִּשֵּׁל בּוֹ קָדָשִׁים וְחוֹלִין\hdot
אוֹ קָדְשֵׁי קָדָשִׁים וְקָדָשִׁים קַלִּים\hdot
אִם יֵשׁ בָּהֶן בְּנוֹתֵן טַעַם\hdot
הֲרֵי הַקַּלִּים נֶאֱכָלִים כַּחֲמוּרִים\hdot
וְאֵינָן טְעוּנִין מְרִיקָה וּשְׁטִיפָה\hdot
וְאֵינָן פּוֹסְלִין בְּמַגָּע\hdot
רָקִיק שֶׁנָּגַע בְּרָקִיק\hdot
וַחֲתִיכָה בַחֲתִיכָה\hdot
לֹא כָל הָרָקִיק וְלֹא כָל הָחֲתִיכוֹת אֲסוּרִין\hdot
אֵינוֹ אָסוּר אֶלָּא מָקוֹם שֶׁבָּלַע׃}

\perek{טְבוּל}{יוֹם וּמְחֻסַּר כִּפּוּרִים\hdot
אֵינָן חוֹלְקִין בַּקֳּדָשִׁים\hdot
לֶאֱכוֹל לָעֶרֶב\hdot
אוֹנֵן נוֹגֵעַ וְאֵינוֹ מַקְרִיב\hdot
וְאֵינוֹ חוֹלֵק לֶאֱכוֹל לָעֶרֶב\hdot
בַּעֲלֵי מוּמִין בֵּין בַּעֲלֵי מוּמִין קְבוּעִין\hdot
בֵּין בַּעֲלֵי מוּמִין עוֹבְרִין\hdot
חוֹלְקִין וְאוֹכְלִין\hdot
אֲבָל לֹא מַקְרִיבִין\hdot
וְכָל שֶׁאֵינוֹ רָאוּי לַעֲבוֹדָה אֵינוֹ חוֹלֵק בַּבָּשָׂר\hdot
וְכָל שֶׁאֵין לוֹ בַבָּשָׂר אֵין לוֹ בְעוֹרוֹ\hdot
אֲפִילוּ טָמֵא בְּשַׁעַת זְרִיקַת דָּמִים\hdot
וְטָהוֹר בְּשַׁעַת הֶקְטֵר חֲלָבִים\hdot
אֵינוֹ חוֹלֵק בַּבָּשָׂר\hdot
שֶׁנֶּאֱמַר הַמַּקְרִיב אֶת דַּם הַשְּׁלָמִים וְאֶת הַחֵלֶב מִבּנֵי אַהֲרֹן\hdot
לוֹ תִהְיֶה שׁוֹק הַיָּמִין לְמָנָה׃}

\halacha{כָּל שֶׁלֹּא זָכָה הַמִּזְבֵּחַ בִּבְשָׂרָהּ לֹא זָכוּ הַכֹּהֲנִים בְּעוֹרָהּ\hdot
שֶׁנֶּאֱמַר עוֹלַת אִישׁ\hdot
עוֹלָה שֶׁעָֽלְתָה לָאִישׁ\hdot
עוֹלָה שֶׁנִּשְׁחֲטָה שֶׁלֹּא לִשְׁמָהּ\hdot
אַף עַל פִּי שֶׁלֹּא עָֽלְתָה לַבְּעָלִים\hdot
עוֹרָהּ לַכֹּהֲנִים\hdot
אֶחָד עוֹלַת הָאִישׁ\hdot
וְאֶחָד עוֹלַת הָאִשָּׁה\hdot
עוֹרוֹתֵיהֶן לַכֹּהֲנִים׃}

\halacha{עוֹרוֹת קָדָשִׁים קַלִּים לַבְּעָלִים\hdot
וְעוֹרוֹת קָדְשֵׁי קָדָשִׁים לַכֹּהֲנִים\hdot
קַל וָחֹמֶר מָה\hdot
אִם עוֹלָה שֶׁלֹּא זָכוּ בִבְשָׂרָהּ\hdot
זָכוּ בְעוֹרָהּ\hdot
קָדְשֵׁי קָדָשִׁים שֶׁזָּכוּ בִבְשָׂרָן\hdot
אֵינוֹ דִין שֶׁיִּזְכּוּ בְעוֹרָן\hdot
אֵין מִזְבֵּחַ יוֹכִיחַ\hdot
שֶׁאֵין לוֹ עוֹר מִכָּל מָקוֹם׃}

\halacha{כָּל הַקָּדָשִׁים שֶׁאִרַע בָּהֶם פִּסּוּל קֹדֶם לְהֶפְשֵׁיטָן\hdot
אֶין עוֹרוֹתֵיהֶם לַכֹּהֲנִים\hdot
לְאַחַר הֶפְשֵׁיטָן עוֹרוֹתֵיהֶם לַכֹּהֲנִים\hdot
אָמַר רִבִּי חֲנִינָא סְגַן הַכֹּהֲנִים מִיָּמַי לֹא רָאִיתִי עוֹר יוֹצֵא לְבֵית הַשְּׂרֵפָה\hdot
אָמַר רִבִּי עֲקִיבָא מִדְּבָרָיו לָמַדְנוּ\hdot
שֶׁהַמַּפְשִׁיט אֶת הַבְּכוֹר\hdot
וְנִמְצָא טְרֵפָה\hdot
שֶׁיֵּאוֹתוּ הַכֹּהֲנִים בְּעוֹרוֹ\hdot
וַחֲכָמִים אוֹמְרִים אֵין לֹא רָאִינוּ רְאָיָה\hdot
אֶלָּא יוֹצֵא לְבֵית הַשְּׂרֵפָה׃}

\halacha{פָּרִים הַנִּשְׂרָפִים וּשְׂעִירִים הַנִּשְׂרָפִים\hdot
בִּזְמָן שֶׁהֵן נִשְׂרָפִין כְּמִצְוָתָן\hdot
נִשְׂרָפִין בְּבֵית הַדֶּשֶׁן\hdot
וּמְטַמְּאִין בְּגָדִים\hdot
וְאִם אֵינָן נִשְׂרָפִין כְּמִצְוָתָן\hdot
נִשְׂרָפִין בְּבֵית הַבִּירָה\hdot
וְאֵינָם מְטַמְּאִין בְּגָדִים׃}

\halacha{הָיוּ סוֹבְלִין אוֹתָן בְּמוֹטוֹת\hdot
יָֽצְאוּ הָרִאשׁוֹנִים חוּץ לְחוֹמַת הָעֲזָרָה\hdot
וְהָאַחֲרוֹנִים לֹא יָֽצְאוּ\hdot
הָרִאשׁוֹנִים מְטַמְּאִין בְּגָדִים\hdot
וְהָאַחֲרוֹנִים אֵינָן מְטַמְּאִין בְּגָדִים\hdot
עַד שֶׁיֵּצְאוּ\hdot
יָֽצְאוּ אֵלּוּ וָאֵלּוּ\hdot
אֵלּוּ וָאֵלּוּ מְטַמְּאִין בְּגָדִים\hdot
רִבִּי שִׁמְעוֹן אוֹמֵר אֵלּוּ וָאֵלּוּ אֵינָן מְטַמְּאִין בְּגָדִים\hdot
עַד שֶׁיֻּצַּת הָאוּר בְּרֻבָּן\hdot
נִתַּךְ הַבָּשָׂר\hdot
אֵין הַשׂוֹרֵף מְטַמֵּא בְּגָדִים׃}

\perek{הַשּׁוֹחֵט}{וְהַמַּעֲלֶה בַחוּץ\hdot
חַיָּב עַל הַשְּׁחִיטָה\hdot
וְחַיָּב עַל הָעֲלִיָּה\hdot
רִבִּי יוֹסֵי הַגָּלִילִי אוֹמֵר שָׁחַט בִּפְנִים וְהֶעֱלָה בַחוּץ חַיָּב\hdot
שָׁחַט בַּחוּץ וְהֶעֱלָה בַחוּץ פָּטוּר\hdot
שֶׁלֹּא הֶעֱלָה בַחוּץ אֶלָּא דָבָר פָּסוּל\hdot
אָֽמְרוּ לוֹ אַף הַשּׁוֹחֵט בִּפְנִים וּמַעֲלֶה בַחוּץ\hdot
כֵּיוַן שֶׁהוֹצִיאוֹ פְּסָלוֹ׃}

\halacha{טָמֵא שֶׁאָכַל בֵּין קֹדֶשׁ טָמֵא\hdot
וּבֵין קֹדֵשׁ טָהוֹר חַיָּב\hdot
רִבִּי יוֹסֵי הַגָּלִילִי אוֹמֵר טָמֵא שֶׁאָכַל טָהוֹר חַיָּב\hdot
וְטָמֵא שֶׁאָכַל טָמֵא פָּטוּר\hdot
שֶׁלֹּא אָכַל אֶלָּא דָבָר טָמֵא\hdot
אָֽמְרוּ לוֹ אַף טָמֵא שֶׁאָכַל טָהוֹר\hdot
כֵּיוַן שֶׁנָּגַע בּוֹ טִמְּאָהוּ\hdot
וְטָהוֹר שֶׁאָכַל טָמֵא פָּטוּר\hdot
שֶׁאֵינוֹ חַיָּב אֶלָּא עַל טֻמְאַת הַגּוּף׃}

\halacha{חֹמֶר בַּשְּׁחִיטָה מִבְּהַעֲלִיָּה\hdot
וּבְהַעֲלִיָּה מִבַּשְּׁחִיטָה\hdot
חֹמֶר בַּשְּׁחִיטָה\hdot
שֶׁהַשּׁוֹחֵט לַהֶדְיוֹט חַיָּב\hdot
וְהַמַּעֲלֶה לַהֶדְיוֹט פָּטוּר\hdot
חֹמֶר בְּהַעֲלִיָּה\hdot
שְׁנַיִם שֶׁאַחֲזוּ בְסַכִּין וְשָׁחֲטוּ פְּטוּרִים\hdot
אַחֲזוּ בָאֵבֶר וְהֶעֱלוּהוּ חַיָּבִין\hdot
הֶעֱלָה וְחָזַר וְהֶעֱלָה וְחָזַר וְהֶעֱלָה\hdot
חַיָּב עַל כָּל הַעֲלִיָּה וְהַעֲלִיָּה\hdot
דִּבְרֵי רִבִּי שִׁמְעוֹן\hdot
רִבִּי יוֹסֵי אוֹמֵר אֵינוֹ חַיָּב אֶלָּא אֶחָת\hdot
וְאֵינוֹ חַיָּב עַד שֶׁיַּעֲלֶה לְרֹאשׁ הַמִּזְבֵּחַ\hdot
רִבִּי שִׁמְעוֹן אוֹמֵר אֲפִילוּ הֶעֱלָה עַל הַסֶּלַע אוֹ עַל הָאֶבֶן חַיָּב׃}

\halacha{אֶחָד קָדָשִׁים כְּשֵׁרִין\hdot
וְאֶחָד קָדָשִׁים פְּסוּלִין\hdot
שֶׁהָיָה פִסּוּלָן בַּקֹּדֶשׁ\hdot
וְהִקְרִיבָן בַּחוּץ חַיָּב\hdot
הַמַּעֲלֶה כַּזַּיִת מִן הָעוֹלָה\hdot
וּמִן הָאֵימוּרִין בַּחוּץ חַיָּב\hdot
הַקֹּמֶץ וְהַלְּבוֹנָה וְהַקְּטֹרֶת\hdot
וּמִנְחַת כֹּהֲנִים\hdot
וּמִנְחַת כֹּהֵן הַמָּשִׁיחַ\hdot
וּמִנְחַת נְּסָכִין\hdot
שֵׁהִקְרִיב מֵאֶחָד מֵהֶן כַּזַּיִת בַּחוּץ חַיָּב\hdot
רִבִּי אֱלִיעֶזֶר פּוֹטֵר עַד שֶׁהִקְרִיב אֶת כֻּלּוֹ\hdot
וְכֻלָּם שֶׁהִקְרִיבָן בִּפְנִים\hdot
וְשִׁיֵּר בָּהֶן כַּזַּיִת וְהִקְרִיבָן בַּחוּץ חַיָּב\hdot
וְכֻלָּם שֶׁחָֽסְרוּ כָּל שֶׁהֵן וְהִקְרִיבָן בַּחוּץ פָּטוּר׃}

\halacha{הַמַּקְרִיב קָדָשִׁים וְאֵמוּרֵיהֶם בַּחוּץ חַיָּב\hdot
מִנְחָה שֶׁלֹּא נִקְמְצָה\hdot
וְהִקְרִיבָהּ בַּחוּץ פָּטוּר\hdot
קְמָצָהּ וְחָזַר קָמְצָהּ לְתוֹכָהּ\hdot
וְהִקְרִיבָהּ בַּחוּץ חַיָּב׃}

\halacha{הַקֹּמֶץ וְהַלְּבוֹנָה שֶׁהִקְרִיב אֶת אַחַד מֵהֶן בַּחוּץ חַיָּב\hdot
רִבִּי אֱלִיעֶזֶר פּוֹטֵר עַד  שֶׁיַּקְרִיב אֶת הַשֵּׁנִי\hdot
אֶחָד בִּפְנִים וְאֶחָד בַּחוּץ חַיָּב\hdot
שְׁנֵי בָזִיכֵי לְבוֹנָה שֶׁהִקְרִיב אֶת אַחַד מֵהֶן בַּחוּץ חַיָּב\hdot
רִבִּי אֱלִיעֶזֶר פּוֹטֵר עַד שֶׁיַּקְרִיב אֶת הַשֵּׁנִי\hdot
אֶחָד בִּפְנִים וְאֶחָד בַּחוּץ חַיָּב\hdot
הַזּוֹרֵק מִקְצַת דָּמִים בַּחוּץ חַיָּב\hdot
רִבִּי אֶלְעָזָר אוֹמֵר אַף הַמְּנַסֵּךְ מֵי חָג בֶּחָג בַּחוּץ חַיָּב\hdot
רִבִּי נְחֶמְיָה אוֹמֵר שִׁיְּרֵי הַדָּם שֶׁהִקְרִיבָן בַּחוּץ חַיָּב׃}

\halacha{הַמּוֹלֵק אֶת הָעוֹף בִּפְנִים\hdot
וְהֶעֱלָה בַחוּץ חַיָּב\hdot
מָלַק בַּחוּץ וְהֶעֱלָה בַחוּץ פָּטוּר\hdot
הַשּׁוֹחֵט אֶת הָעוֹף בִּפְנִים וְהֶעֱלָה בַחוּץ פָּטוּר\hdot
שָׁחַט בַּחוּץ וְהֶעֱלָה בַחוּץ חַיָּב\hdot
נִמְצָא דֶּרֶךְ הֶכְשֵׁרוֹ מִבִּפְנִים פִּטּוּרוֹ בַחוּץ\hdot
דֶּרֶךְ הֶכְשֵׁרוֹ בַחוּץ פִּטּוּרוֹ בִּפְנִים\hdot
רִבִּי שִׁמְעוֹן אוֹמֵר כָּל שֶׁחַיָּבִין עָלָיו בַּחוּץ חַיָּבִין עַל כַּיּוֹצֵא בוֹ בִּפְנִים שֶׁהֶעֱלָהוּ בַחוּץ\hdot
חוּץ מִן הַשּׁוֹחֵט בִּפְנִים\hdot
וּמַעֲלֶה בַחוּץ׃}

\halacha{הַחַטָּאת שֶׁקִּבֵּל דָּמָהּ בְּכוֹס אֶחָד\hdot
נָתַן בַּחוּץ וְחָזַר וְנָתַן בִּפְנִים\hdot
בִּפְנִים וְחָזַר וְנָתַן בַּחוּץ חַיָּב\hdot
שֶׁכֻּלּוֹ רָאוּי לָבוֹא בִּפְנִים\hdot
קִבֵּל דָּמָהּ בִּשְׁנֵי כוֹסוֹת\hdot
(נָתַן שְׁנֵיהֶם בִּפְנִים פָּטוּר)\hdot
שְׁנֵיהֶם בַּחוּץ חַיָּב\hdot
אֶחָד בִּפְנִים וְאֶחָד בַּחוּץ פָּטוּר\hdot
אֶחָד בַּחוּץ וְאֶחָד בִּפְנִים\hdot
חַיָּב עַל הַחִיצוֹן\hdot
וְהַפְּנִימִי מְכַפֵּר\hdot
לְמָה הַדָּבָר דּוֹמֶה לְמַפְרִישׁ חַטָּאתוֹ וְנֶאֱבָדָה\hdot
וְהִפְרִישׁ אַחֶרֶת תַּחְתֶּיהָ\hdot
וְאַחַר כַּךְ נִמְצֵאת הָרִאשׁוֹנָה\hdot
וַהֲרֵי שְׁתֵּיהֶן עוֹמְדוֹת\hdot
שָׁחַט שְׁתֵּיהֶן בִּפְנִים פָּטוּר\hdot
שָׁחַט שְׁתֵּיהֶן בַּחוּץ חַיָּב\hdot
אַחַת בִּפְנִים וְאַחַת בַּחוּץ פָּטוּר\hdot
אַחַת בַּחוּץ וְאַחַת בִּפְנִים\hdot
חַיָּב עַל הַחִיצוֹנָה\hdot
וְהַפְּנִימִית מְכַפֶּרֶת\hdot
כְּשֵׁם שֶׁדָּמָהּ פּוֹטֵר אֶת בְּשָׂרָהּ\hdot
כַּךְ הוּא פּוֹטֵר אֶת בְּשַׂר חֲבֵרְתָּהּ׃}

\perek{פָּרַת}{חַטָּאת שֶׁשְּׂרָפָהּ חוּץ מִגִּתָּהּ\hdot
וְכֵן שָׂעִיר הַמִּשְׁתַּלֵּחַ שֶׁהִקְרִיבוֹ בַחוּץ פָּטוּר\hdot
שֶׁנֶּאֱמַר וְאֶל פֶּתַח אֹהֶל מוֹעֵד לֹא הֵבִיאוֹ\hdot
כָּל שֶׁאֵינוֹ רָאוּי לָבֹא אֶל פֶּתַח אֹהֶל מוֹעֵד\hdot
אֵין חַיָּבִין עָלָיו׃}

\halacha{הָרוֹבֵעַ וְהַנִּרְבָּע\hdot
וְהַמֻּקְצֶה\hdot
וְהַנֶּעֱבָד\hdot
וְהָאֶתְנָן\hdot
וְהַמְּחִיר\hdot
וְהַכִּלְאַיִם\hdot
וְהַטְּרֵפָה\hdot
וְיוֹצֵא דֹפֶן\hdot
שֶׁהִקְרִיבָן בַּחוּץ פָּטוּר\hdot
שֶׁנֶּאֱמַר לִפְנֵי מִשְׁכַּן יְיָ\hdot
כָּל שֶׁאֵינוֹ רָאוּי לָבֹא לִפְנֵי מִשְׁכַּן יְיָ\hdot
אֵין חַיָּבִין עָלָיו\hdot
בַּעֲלֵי מוּמִין בֵּין בַּעֲלֵי מוּמֵין קְבוּעִים\hdot
בֵּין בַּעֲלֵי מוּמִין עוֹבְרִים\hdot
שֶׁהִקְרִיבָן בַּחוּץ פָּטוּר\hdot
רִבִּי שִׁמְעוֹן אוֹמֵר בַּעֲלֵי מוּמִין קְבוּעִין פָּטוּר\hdot
וּבַעֲלֵי מוּמִין עוֹבְרִין\hdot
עוֹבְרִין בְּלֹא תַעֲשֶׂה\hdot
תּוֹרִים שֶׁלֹּא הִגִּיעַ זְמַנָּן\hdot
וּבְנֵי יוֹנָה שֶׁעָבַר זְמַנָּן\hdot
שֶׁהִקְרִיבָן בַּחוּץ פָּטוּר\hdot
רִבִּי שִׁמְעוֹן אוֹמֵר בְּנֵי יוֹנָה שֶׁעָבַר זְמַנָּן פָּטוּר\hdot
וְתוֹרִים שֶׁלֹּא הִגִּיעַ זְמַנָּן בְּלֹא תַעֲשֶׂה\hdot
אוֹתוֹ וְאֶת בְּנוֹ וּמְחֻסַּר זְמָן פָּטוּר\hdot
רִבִּי שִׁמְעוֹן אוֹמֵר הֲרֵי זֶה בְּלֹא תַעֲשֶׂה\hdot
שֶׁהָיָה רִבִּי שִׁמְעוֹן אוֹמֵר כָּל שֶׁהוּא רָאוּי לָבוֹא לְאַחַר זְמָן\hdot
הֲרֵי זֶה בְּלֹא תַעֲשֶׂה\hdot
וְאֵין בּוֹ כָּרֵת\hdot
וַחֲכָמִים אוֹמְרִים כָּל שֶׁאֵין בּוֹ כָּרֵת אֵין בּוֹ בְּלֹא תַעֲשֶׂה׃}

\halacha{מְחֻסַּר זְמָן\hdot
בֵּין בְּגוּפוֹ בֵּין בִּבְעָלָיו\hdot
אֵיזֶה הוּא מְחֻסַּר זְמָן בִבְעָלָיו\hdot
הַזָּב וְהַזָּבָה (וְיוֹלֶדֶת וּמְצוֹרָע) שֶׁהִקְרִיבוּ חַטָּאתָם וַאֲשָׁמָם בַּחוּץ פְּטוּרִין\hdot
עוֹלוֹתֵיהֶן וְשַׁלְמֵיהֶן בַּחוּץ חַיָּבִין\hdot
הַמַּעֲלֶה מִבְּשַׂר חַטָּאת\hdot
מִבְּשַׂר אָשָׁם\hdot
מִבְּשַׂר קָדְשֵׁי קָדָשִׁים\hdot
מִבְּשַׂר קָדָשִׁים קַלִּים\hdot
וּמוֹתַר הָעֹמֶר\hdot
וּשְׁתֵּי הַלֶּחֶם\hdot
וְלֶחֶם הַפָּנִים\hdot
וְשִׁיְּרֵי מְנָחוֹת\hdot
הַיּוֹצֵק\hdot
הַבּוֹלֵל\hdot
הַפּוֹתֵת\hdot
הַמּוֹלֵחַ\hdot
הַמֵּנִיף\hdot
הַמַּגִּישׁ\hdot
הַמְּסַדֵּר אֶת הַשֻּׁלְחָן\hdot
וְהַמֵּטִיב אֶת הַנֵּרוֹת\hdot
וְהַקּוֹמֵץ\hdot
וְהַמְּקַבֵּל דָּמִים בַּחוּץ פָּטוּר\hdot
אֵין חַיָּבִין עָלָיו\hdot
לֹא מִשּׁוּם זָרוּת\hdot
וְלֹא מִשּׁוּם טֻמְאָה\hdot
וְלֹא מִשּׁוּם מְחֻסַּר בְּגָדִים\hdot
וְלֹא מִשּׁוּם רְחוּץ יָדַיִם וְרַגְלַיִם׃}

\halacha{עַד שֶׁלֹּא הוּקַם הַמִּשְׁכָּן הָיוּ הַבָּמוֹת מוּתָרוֹת\hdot
וַעֲבוֹדָה בַבְּכוֹרוֹת\hdot
מִשֶּׁהוּקַם הַמִּשְׁכָּן\hdot
נֶאֱסְרוּ הַבָּמוֹת\hdot
וַעֲבוֹדָה בַכֹּהֲנִים\hdot
קָדְשֵׁי קָדָשִׁים נֶאֱכָלִים לִפְנִים מִן הַקְּלָעִים\hdot
קָדָשִׁים קַלִּים בְּכָל מַחֲנֵה יִשְׂרָאֵל׃}

\halacha{בָּאוּ לַגִּלְגָּל וְהֻתְּרוּ הַבָּמוֹת\hdot
קָדְשֵׁי קָדָשִׁים נֶאֱכָלִים לִפְנִים מִן הַקְּלָעִים\hdot
קָדָשִׁים קַלִּים בְּכָל מָקוֹם׃}

\tractatehead{מסכת מנחות}{menahot}

\perek{כָּל}{הַמְּנָחוֹת שֶׁנִּקְמְצוּ שֶׁלֹּא לִשְׁמָן\hdot
כְּשֵׁרוֹת\hdot
אֶלָּא שֶׁלֹּא עָלוּ לַבְּעָלִים לְשֵׁם חוֹבָה\hdot
חוּץ מִמִּנְחַת חוֹטֵא וּמִנְחַת קְנָאוֹת\hdot
מִנְחַת חוֹטֵא וּמִנְחַת קְנָאוֹת שֶׁקְּמָצָן שֶׁלֹּא לִשְׁמָן\hdot
נָתַן בַּכֶּלִי וְהָלַךְ וְהִקְטִיר שֶׁלֹּא לִשְׁמָן\hdot
אוֹ לִשְׁמָן וְשֶׁלֹּא לִשְׁמָן\hdot
אוֹ שֶׁלֹּא לִשְׁמָן וְלִשְׁמָן פְּסוּלוֹת\hdot
כֵּיצַד לִשְׁמָן\hdot
וְשֶׁלֹּא לִשְׁמָן\hdot
לְשֵׁם מִנְחַת חוֹטֵא\hdot
וּלְשֵׁם מִנְחַת נְדָבָה\hdot
אוֹ שֶׁלֹּא לִשְׁמָן וְלִשְׁמָן\hdot
לְשֵׁם מִנְחַת נְדָבָה\hdot
וּלְשֵׁם מִנְחַת חוֹטֵא׃}

\halacha{אֶחָד מִנְחַת חוֹטֵא\hdot
וְאֶחָד כָּל הַמְּנָחוֹת שֶׁקְּמָצָן זָר\hdot
אוֹנֵן\hdot
טְבוּל יוֹם\hdot
מְחֻסַּר בְּגָדִים\hdot
מְחֻסַּר כִּפּוּרִים\hdot
שֶׁלֹּא רְחוּץ יָדַיִם וְרַגְלַיִם\hdot
עָרֵל טָמֵא יוֹשֵׁב\hdot
עוֹמֵד עַל גַּבֵּי כֵלִים\hdot
עַל גַּבֵּי בְהֵמָה עַל גַּבֵּי רַגְלֵי חֲבֵרוֹ פָּסַל\hdot
קָמַץ בַּשְּׂמֹאל פָּסַל\hdot
בֵּן בְּתֵירָא אוֹמֵר יַחֲזִיר וְיַחֲזוֹר וְיִקְמוֹץ בְּיָמִין\hdot
קָמַץ וְעָלָה בְיָדוֹ צְרוֹר\hdot
אוֹ גַרְגִּיר מֶלַח\hdot
אוֹ קֹרֶט שֶׁל לְבוֹנָה פָּסַל\hdot
מִפְּנֵי שְׁאָֽמְרוּ הַקֹּמֶץ הַיָּתֵר וְהֶחָסֵר פָּסוּל\hdot
אֵיזֶהוּ הַיָּתֵר שֶׁקְּמָצוֹ מְבוֹרָץ\hdot
וְחָסֵר שֶׁקְּמָצוֹ בְּרָאשֵׁי אֶצְבְּעוֹתָיו\hdot
כֵּיצַד הוּא עוֹשֶׂה\hdot
פּוֹשֵׁט אֶת אֶצְבְּעוֹתָיו עַל פַּס יָדוֹ׃}

\halacha{רִבָּה שַׁמְנָהּ\hdot
וְחִסֵּר שַׁמְנָהּ\hdot
חִסֵּר לְבוֹנָתָהּ פְּסוּלָה\hdot
הַקּוֹמֵץ אֶת הַמִּנְחָה לֶאֱכוֹל שִׁירֶיהָ בַּחוּץ\hdot
אוֹ כַּזַּיִת מִשִּׁירֶיהָ בַּחוּץ\hdot
לְהַקְטִיר קֻמְצָהּ בַּחוּץ\hdot
אוֹ כַּזַּיִת מִקֻּמְצָהּ בַּחוּץ\hdot
אוֹ לְהַקְטִיר לְבוֹנָתָהּ בַּחוּץ\hdot
פָּסוּל וְאֵין בּוֹ כָּרֵת\hdot
לֶאֱכוֹל שִׁירֶיהָ לְמָחָר אוֹ כַּזַּיִת מִשִּׁירֶיהָ לְמָחָר\hdot
לְהַקְטִיר קֻמְצָהּ לְמָחָר אוֹ כַּזַּיִת מִקֻּמְצָהּ לְמָחָר\hdot
אוֹ לְהַקְטִיר לְבוֹנָתָהּ לְמָחָר\hdot
פּגּוּל וְחַיָּבִין עָלָיו כָּרֵת\hdot
זֶה הַכְּלָל\hdot
כָּל הַקּוֹמֵץ וְהַנּוֹתֵן בֶּכֶּלִי\hdot
וְהַמְּהַלֵּךְ וְהַמַּקְטִיר\hdot
לֶאֱכוֹל דָּבָר שֶׁדַּרְכּוֹ לֶאֱכוֹל\hdot
וּלְהַקְטִיר דָּבָר שֶׁדַּרְכּוֹ לְהַקְטִיר\hdot
חוּץ לִמְקוֹמוֹ פָּסוּל וְאֵין בּוֹ כָּרֵת\hdot
חוּץ לִזְמַנּוֹ פִּגּוּל וְחַיָּבִין עָלָיו כָּרֵת\hdot
וּבִלְבַד שֶׁיִּקָּרֵב הַמַּתִּיר כְּמִצְוָתוֹ\hdot
כֵּיצַד קָרַב הַמַּתִּיר כְּמִצְוָתוֹ\hdot
קָמַץ בִּשְׁתִּיקָה וְנָתַן בַּכֶּלִי\hdot
וְהָלַךְ לְהַקְטִיר חוּץ לִזְמַנּוֹ\hdot
אוֹ שֶׁקָּמַץ חוּץ לִזְמַנּוֹ וְנָתַן בַּכֶּלִי\hdot
וְהָלַךְ וְהִקְטִיר בִּשְׁתִּיקָה\hdot
אוֹ שֶׁקָּמַץ וְנָתַן בַּכֶּלִי\hdot
וְהָלַךְ וְהִקְטִיר חוּץ לִזְמַנּוֹ\hdot
זֶה הוּא שֶׁקָרַב הַמַּתִּיר כְּמִצְוָתוֹ׃}

\halacha{כֵּיצַד לֹא קָרַב הַמַּתִּיר כְּמִצְוָתוֹ\hdot
קָמַץ חוּץ לִמְקוֹמוֹ\hdot
וְנָתַן בַּכֶּלִי\hdot
וְהָלַךְ וְהִקְטִיר חוּץ לִזְמַנּוֹ\hdot
אוֹ שֶׁקָּמַץ חוּץ לִזְמַנּוֹ\hdot
וְנָתַן בַּכֶּלִי\hdot
וְהָלַךְ וְהִקְטִיר חוּץ לִמְקוֹמוֹ\hdot
אוֹ שֶׁקָּמַץ וְנָתַן בַּכֶּלִי\hdot
וְהָלַךְ וְהִקְטִיר חוּץ לִמְקוֹמוֹ\hdot
מִנְחַת חוֹטֵא וּמִנְחַת קְנָאוֹת שֶׁקְּמָצָן שֶׁלֹּא לִשְׁמָן\hdot
וְנָתַן בַּכֶּלִי\hdot
וְהָלַךְ וְהִקְטִיר חוּץ לִזְמַנָּן\hdot
אוֹ שֶׁקָּמַץ חוּץ לִזְמַנָּן\hdot
וְנָתַן בַּכֶּלִי\hdot
וְהָלַךְ וְהִקְטִיר שֶׁלֹּא לִשְׁמָן\hdot
אוֹ שֶׁקָּמַץ וְנָתַן בַּכֶּלִי\hdot
וְהָלַךְ וְהִקְטִיר שֶׁלֹּא לִשְׁמָן\hdot
זֶה הוּא שֶׁלֹּא קָרַב הַמַּתִּיר כְּמִצְוָתוֹ\hdot
לֶאֱכוֹל כַּזַּיִת בַּחוּץ וְכַזַּיִת לְמָחָר\hdot
כַּזַּיִת לְמָחָר וְכַזַּיִת בַּחוּץ\hdot
כַּחֲצִי זַיִת בַּחוּץ וְכַחֲצִי זַיִת לְמָחָר\hdot
כַּחֲצִי זַיִת לְמָחָר וְכַחֲצִי זַיִת בַּחוּץ\hdot
 פָּסוּל וְאֵין בּוֹ כָרֵת\hdot
אָמַר רִבִּי יְהוּדָה זֶה הַכְּלָל אִם מַחֲשֶׁבֶת הַזְּמָן קָֽדְמָה לְמַחֲשֶׁבֶת הַמָּקוֹם\hdot
פִּגּוּל וְחַיָּבִין עָלָיו כָּרֵת\hdot
וְאִם מַחֲשֶׁבֶת הַמָּקוֹם\hdot
קָֽדְמָה לְמַחֲשֶׁבֶת הַזְּמָן\hdot
פָּסוּל וְאֵין בּוֹ כָרֵת\hdot
וַחֲכָמִים אוֹמְרִים זֶה וְזֶה פָּסוּל וְאֵין בּוֹ כָרֵת\hdot
לֶאֱכוֹל כַּחֲצִי זַיִת\hdot
וּלְהַקְטִיר כַּחֲצִי זַיִת כָּשֵׁר\hdot
שֶׁאֵין אֲכִילָה וְהַקְטָרָה מִצְטָֽרְפִין׃}


\perek{הַקּוֹמֵץ}{אֶת הַמִּנְחָה לֶאֱכוֹל שִׁיְּרֶיהָ\hdot
אוֹ לְהַקְטִיר קֻמְצָהּ לְמָחָר\hdot
מוֹדֶה רִבִּי יוֹסֵי בְּזֶה\hdot
שֶׁהוּא פִגּוּל וְחַיָּבִין עָלָיו כָּרֵת\hdot
לְהַקְטִיר לְבוֹנָתָהּ לְמָחָר\hdot
רִבִּי יוֹסֵי אוֹמֵר פָּסוּל וְאֵין בּוֹ כָרֵת\hdot
וַחֲכָמִים אוֹמְרִים פִּגּוּל וְחַיָּבִין עָלָיו כָּרֵת\hdot
אָֽמְרוּ לוֹ מַה שְּׁנָא זוֹ מִן הַזֶּבַח\hdot
אָמַר לָהֶם\hdot
שֶׁהַזֶּבַח דָּמוֹ וּבְשָׂרוֹ וְאֵמוּרָיו\hdot
אֶחָד\hdot
וּלְבוֹנָה אֵינָהּ מִן הַמִּנְחָה׃}

\halacha{שָׁחַט שְׁנֵי כְבָשׂים\hdot
לֶאֱכוֹל אַחַת מִן הַחַלּוֹת לְמָחָר\hdot
הִקְטִיר שְׁנֵי בָזִיכִין\hdot
לֶאֱכוֹל אֶחָד מִן הַסְּדָרִים לְמָחָר\hdot
רִבִּי יוֹסֵי אוֹמֵר אוֹתָהּ הַחַלָּה וְאוֹתוֹ הַסֵּדֶר שֶׁחָשַׁב עָלָיו\hdot
פִּגּוּל וְחַיָּבִין עָלָיו כָּרֵת\hdot
וְהַשֵּׁנִי פָּסוּל וְאֵין בּוֹ כָרֵת\hdot
וַחֲכָמִים אוֹמְרִים זֶה וזֶה פִּגּוּל וְחַיָּבִין עָלָיו כָּרֵת\hdot
נִטְמֵאת אַחַת מִן הַחַלּוֹת אוֹ אֶחָד מִן הַסְּדָרִים\hdot
רִבִּי יְהוּדָה אוֹמֵר שְׁנֵיהֶם יֵצְאוּ לְבֵית הַשְּׂרֵיפָה\hdot
שֶׁאֵין קָרְבָּן צִבּוּר חָלוּק\hdot
וַחֲכָמִים אוֹמְרִים הַטָּמֵא בְטֻמְאָתוֹ וְהַטָּהוֹר יֵאָכֵל׃}

\halacha{הַתּוֹדָה מְפַגֶּלֶת אֶת הַלֶּחֶם\hdot
וְהַלֶּחֶם אֵינוֹ מְפַגֵּל אֶת הַתּוֹדָה\hdot
כֵּיצַד\hdot
הַשּׁוֹחֵט אֶת הַתּוֹדָה לֶאֱכוֹל מִמֶּנָּה לְמָחָר\hdot
הוּא וְהַלֶּחֶם מְפֻגָּלִין\hdot
לֶאֱכוֹל מִן הַלֶּחֶם לְמָחָר\hdot
הַלֶּחֶם מְפֻגָּל\hdot
וְהַתּוֹדָה אֵינָהּ מְפֻגֶּלֶת\hdot
הַכְּבָשִׂים מְפַגְּלִין אֶת הַלֶּחֶם\hdot
וְהַלֶּחֶם אֵינוֹ מְפַגֵּל אֶת הַכְּבָשִׂים\hdot
כֵּיצַד\hdot
הַשּׁוֹחֵט אֶת הַכְּבָשִׂים לֶאֱכוֹל מֵהֶן לְמָחָר\hdot
הֵם וְהַלֶּחֶם מְפֻגָּלִין\hdot
לֶאֱכוֹל מִן הַלֶּחֶם לְמָחָר\hdot
הַלֶּחֶם מְפֻגָּל\hdot
וְהַכְּבָשִׂים אֵינָם מְפֻגָּלִין׃}

\halacha{הַזֶּבַח מְפַגֵּל אֶת הַנְּסָכִין\hdot
מֵשֶּׁקִּדְּשׁוּ בַכֶּלִי\hdot
דִּבְרֵי רִבִּי מֵאִיר\hdot
וְהַנְּסָכִין אֵינָן מְפַגְּלִין אֶת הַזֶּבַח\hdot
כֵּיצַד\hdot
הַשּׁוֹחֵט אֶת הַזֶּבַח לֶאֱכוֹל מִמֶּנּוּ לְמָחָר\hdot
הוּא וּנְסָכָיו מְפֻגָּלִין\hdot
לְהַקְרִיב מִן הַנְּסָכִין לְמָחָר\hdot
הַנְּסָכִין מְפֻגָּלִין\hdot
וְהַזֶּבַח אֵינוֹ מְפֻגָּל}

\halacha{פִּגֵּל בַּקֹּמֶץ וְלֹא בַלְּבוֹנָה\hdot
בַּלְּבוֹנָה וְלֹא בַקֹּמֶץ\hdot
רִבִּי מֵאִיר אוֹמֵר פִּגּוּל וְחַיָּבִין עָלָיו כָּרֵת\hdot
וַחֲכָמִים אוֹמְרִים אֵין בּוֹ כָּרֵת עַד שֶׁיְּפַגֵּל אֶת כָּל הַמַּתִּיר\hdot
מוֹדִים חֲכָמִים לְרִבִּי מֵאִיר בְּמִנְחַת חוֹטֵא וּבְמִנְחַת קְנָאוֹת\hdot
שֶׁאִם פִּגֵּל בַּקֹּמֶץ\hdot
שֶׁהוּא פִּגּוּל וְחַיָּבִין עָלָיו כָּרֵת\hdot
שֶׁהַקֹּמֶץ הוּא הַמַּתִּיר\hdot
שָׁחַט אֶחָד מִן הַכְּבָשִׂים לֶאֱכוֹל שְׁתֵּי חַלּוֹת לְמָחָר\hdot
הִקְטִיר אֶחָד מִן הַבָּזִיכִין לֶאֱכוֹל שְׁנֵי סְדָרִים לְמָחָר\hdot
רִבִּי מֵאִיר אוֹמֵר פִּגּוּל וְחַיָּבִין עָלָיו כָּרֵת\hdot
וַחֲכָמִים אוֹמְרִים אֵין בּוֹ כָּרֵת עַד שֶׁיְּפַגֵּל אֶת כָּל הַמַּתִּיר\hdot
שָׁחַט אֶחָד מִן הַכְּבָשִׂים לֶאֱכוֹל מִמֶּנּוּ לְמָחָר\hdot
הוּא פִּגּוּל וַחֲבֵרוֹ כָּשֵׁר\hdot
לֶאֱכוֹל מֵחֲבֵרוֹ לְמָחָר\hdot
שְׁנֵיהֶם כְּשֵׁרִים׃}

\perek{הַקּוֹמֵץ}{אֶת הַמִּנְחָה\hdot
לֶאֱכוֹל דָּבָר שֶׁאֵין דַּרְכּוֹ לֶאֱכוֹל\hdot
לְהַקְטִיר דָּבָר שֶׁאֵין דַּרְכּוֹ לְהַקְטִיר\hdot
כָּשֵׁר\hdot
רִבִּי אֱלִיעֶזֶר פּוֹסֵל\hdot
לֶאֱכוֹל דָּבָר שֶׁדַּרְכּוֹ לֶאֱכוֹל\hdot
לְהַקְטִיר דָּבָר שֶׁדַּרְכּוֹ לְהַקְטִיר\hdot
פָּחוּת מִכַּזַּיִת כָּשֵׁר\hdot
לֶאֱכוֹל כַּחֲצִי זַיִּת וּלְהַקְטִיר כַּחֲצִי זַיִת כָּשֵׁר\hdot
שֶׁאֵין אֲכִילָה וְהַקְטָרָה מִצְטָֽרְפִין׃}

\halacha{לֹא יָצַק\hdot
לֹא בָלַל\hdot
לֹא פָתַת\hdot
לֹא מָלַח\hdot
לֹא הֵנִיף\hdot
לֹא הִגִּישׁ\hdot
אוֹ שֶׁפִּתְּתָן פִּתּוֹת מְרֻבּוֹת וְלֹא מְשָׁחָן\hdot
כּשֵׁרוֹת\hdot
נִתְעָרַב קֻמְצָהּ\hdot
בְּקֹמֶץ חֲבֵרְתָּהּ\hdot
בְּמִנְחַת כֹּהֲנִים\hdot
בְּמִנְחַת כֹּהֵן מָשִׁיחַ\hdot
בְּמִנְחַת נְסָכִין\hdot
כְּשֵׁרָה\hdot
רִבִּי יְהוּדָה אוֹמֵר בְּמִנְחַת כֹּהֵן מָשִׁיחַ\hdot
וּבְמִנְחַת נְסָכִין פְּסוּלָה\hdot
שֶׁזּוֹ בְּלִילָתָהּ עָבָה\hdot
וְזוֹ בְּלִילָתָהּ רַכָּה\hdot
וְהֵן בּוֹלְעוֹת זוֹ מִזּוֹ׃}

\halacha{שְׁתֵּי מְנָחוֹת שֶׁלֹּא נִקְמְצוּ\hdot
וְנִתְעָֽרְבוּ זוֹ בְזוֹ\hdot
אִם יָכוֹל לִקְמוֹץ מִזּוֹ בִּפְנֵי עַצְמָהּ\hdot
וּמִזּוֹ בִּפְנֵי עַצְמָהּ\hdot
כְּשֵׁרוֹת\hdot
וְאִם לָאוּ פְּסוּלוֹת\hdot
הַקֹּמֶץ שֶׁנִּתְעָרַב בְּמִנְחָה שֶׁלֹּא נִקְמָצָה\hdot
לֹא יַקְטִיר\hdot
וְאִם הִקְטִיר\hdot
זוֹ שֶׁנִּקְמְצָה עָֽלְתָה לַבְּעָלִים\hdot
וְזוֹ שֶׁלֹּא נִקְמְצָה לֹא עָֽלְתָה לַבְּעָלִים\hdot
נִתְעָרַב קֻמְצָהּ בְּשִׁיְּרֶיָה\hdot
אוֹ בְּשִׁיְּרֶיהָ שֶׁל חֲבֵרְתָּהּ\hdot
לֹא יַקְטִיר\hdot
וְאִם הִקְטִיר עָֽלְתָה לַבְּעָלִים\hdot
נִטְמָא הַקֹּמֶץ וְהִקְרִיבוֹ\hdot
הַצִּיץ מְרַצֶּה\hdot
יָצָא וְהִקְרִיבוֹ אֵין הַצִּיץ מְרַצֶּה\hdot
שֶׁהַצִּיץ מְרַצֶּה עַל הַטָּמֵא\hdot
וְאֵינוֹ מְרַצֶּה עַל הַיּוֹצֵא׃}

\halacha{נִטְמְאוּ שִׁיְּרֶיהָ\hdot
נִשְׂרְפוּ שִׁיְּרֶיהָ\hdot
נֶאֶבְדוּ שִׁיְּרֶיהָ\hdot
כְּמִדַּת רִבִּי אֱלִיעֶזֶר כְּשֵׁרָה\hdot
וּכְמִדַּת רִבִּי יְהוֹשֻׁעַ פְּסוּלָה\hdot
שֶׁלֹּא בִּכְלֵי שָׁרֵת פְּסוּלָה\hdot
רִבִּי שִׁמְעוֹן מַכְשִׁיר\hdot
הִקְטִיר קֻמְצָהּ פַּעֲמַיִם כְּשֵׁרָה׃}

\halacha{הַקֹּמֶץ מִיעוּטוֹ מְעַכֵּב אֶת רֻבּוֹ\hdot
הָעִשָּׂרוֹן מִעוּטוֹ מְעַכֵּב אֶת רֻבּוֹ\hdot
הַיַּין מִעוּטוֹ מְעַכֵּב אֶת רֻבּוֹ\hdot
הַשֶּׁמֶן מִעוּטוֹ מְעַכֵּב אֶת רֻבּוֹ\hdot
הַסֹּלֶת וְהַשֶּׁמֶן מְעַכְּבִין זֶה אֶת זֶה\hdot
הַקֹּמֶץ וְהַלְּבוֹנָה מְעַכְּבִין זֶה אֶת זֶה׃}

\halacha{שְׁנֵי שְׂעִירֵי יוֹם הַכִּפּוּרִים מְעַכְּבִין זֶה אֶת זֶה\hdot
שְׁנֵי כִבְשֵׂי עֲצֶרֶת מְעַכְּבִין זֶה אֶת זֶה\hdot
שְׁתֵּי חַלּוֹת מְעַכְּבוֹת זוֹ אֶת זוֹ\hdot
שְׁנֵי סְדָרִים מְעַכְּבִין זֶה אֶת זֶה\hdot
שְׁנֵי בָזִיכִין מְעַכְּבִין זֶה אֶת זֶה\hdot
הַסְּדָרִים וְהַבָּזִיכִין מְעַכְּבִין זֶה אֶת זֶה\hdot
שְׁנֵי מִינִים שֶׁבַּנָּזִיר\hdot
שְׁלֹשָׁה שֶׁבַּפָּרָה\hdot
אַרְבָּעָה שֶׁבַּתּוֹדָה\hdot
אַרְבָּעָה שֶׁבַּלּוּלָב\hdot
אַרְבָּעָה שֶׁבַּמְּצוֹרָע\hdot
מְעַכְּבִין זֶה אֶת זֶה\hdot
שֶׁבַע הַזָּיוֹת שֶׁבַּפָּרָה מְעַכְּבוֹת זוֹ אֶת זוֹ\hdot
שֶׁבַע הַזָּיוֹת שֶׁל בֵּין הַבָּדִים וְשֶׁעַל הַפָּרֹכֶת וְשֶׁעַל מִזְבַּח הַזָּהָב\hdot
מְעַכְּבוֹת זוֹ אֶת זוֹ׃}

\halacha{שִׁבְעָה קְנֵי מְנוֹרָה מְעַכְּבִין זֶה אֶת זֶה\hdot
שִׁבְעָה נֵרוֹתֶיהָ מְעַכְּבִין זֶה אֶת זֶה\hdot
שְׁתֵּי פָרָשִׁיּוֹת שֶׁבַּמְּזוּזָה\hdot
מְעַכְּבוֹת זוֹ אֶת זוֹ\hdot
וַאֲפִילוּ כְתָב אֶחַד מְעַכְּבָן\hdot
אַרְבַּע פָּרְשִׁיּוֹת שֶׁבַּתְּפִלִּין\hdot
מְעַכְּבוֹת זוֹ אֶת זוֹ\hdot
וַאֲפִילוּ כְתָב אֶחַד מְעַכְּבָן\hdot
אַרְבַּע צִיצִיּוֹת מְעַכְּבוֹת זוֹ אֶת זוֹ\hdot
שֶׁאַרְבַּעְתָּן מִצְוָה אֶחָת\hdot
רִבִּי יִשְׁמָעֵאל אוֹמֵר אַרְבַּעְתָן אַרְבַּע מִצְוֺת׃}

\perek{הַתְּכֵלֶת}{אֵינָהּ מְעַכֶּבֶת אֶת הַלָּבָן\hdot
וְהַלָּבָן אֵינוֹ מְעַכֵּב אֶת הַתְּכֵלֶת\hdot
תְּפִלָּה שֶׁל יָד אֵינָהּ מְעַכֶּבֶת שֶׁל רֹאשׁ\hdot
וְשֶׁל רֹאשׁ אֵינָהּ מְעַכֶּבֶת שֶׁל יָד\hdot
הַסֹּלֶת וְהַשֶּׁמֶן אֵינָן מְעַכְּבִין אֶת הַיַּיִן\hdot
וְלֹא הַיַּיִן מְעַכְּבָן\hdot
הַמַּתָּנוֹת שֶׁעַל הַמִּזְבֵּחַ הַחִיצוֹן\hdot
אֵינָן מְעַכְּבוֹת זוֹ אֶת זוֹ׃}

\halacha{הַפָּרִים וְהָאֵלִים וְהַכְּבָשִׂים\hdot
אֵינָן מְעַכְּבִין זֶה אֶת זֶה\hdot
רִבִּי שִׁמְעוֹן אוֹמֵר אִם הָיוּ לָהֶם פָּרִים מְרֻבִּים\hdot
וְלֹא הָיוּ נְסָכִים\hdot
יָבִיאוּ פָר אֶחָד וּנְסָכָיו\hdot
וְלֹא יִקָּֽרְבוּ כֻלָּן בְּלֹא נְסָכִין׃}

\halacha{הַפָּר וְהָאֵלִים וְהַכְּבָשִׂים וְהַשָּׂעִיר\hdot
אֵינָן מְעַכְּבִין אֶת הַלֶּחֶם\hdot
וְלֹא הַלֶּחֶם מְעַכְּבָן\hdot
הַלֶחֶם מְעַכֵּב אֶת הַכְּבָשִׂים\hdot
וְהַכְּבָשִׂים אֵינָן מְעַכְּבִין אֶת הַלֶּחֶם\hdot
דִּבְרֵי רִבִּי עֲקִיבָא\hdot
אָמַר רִבִּי שִׁמְעוֹן בֶּן נָנָס לֹא כִי\hdot
אֶלָּא הַכְּבָשִׂים מְעַכְּבִין (אֶת) הַלֶּחֶם\hdot
וְהַלֶּחֶם אֵינוֹ מְעַכֵּב אֶת הַכְּבָשִׂים\hdot
שֶׁכֵּן מָצִינוּ כְּשֶׁהָיוּ יִשְׂרָאֵל בַּמִּדְבָּר אַרְבָּעִים שָׁנָה\hdot
קָֽרְבוּ כְבָשִׂים בְּלֹא לֶחֶם\hdot
אַף (כָּאן) יִקָּֽרְבוּ כְבָשִׂים בְּלֹא לֶחֶם\hdot
אָמַר רִבִּי שִׁמְעוֹן הֲלָכָה כְּדִבְרֵי בֶּן נָנָס\hdot
אֲבָל אֵין הַטַּעַם כִּדְבָרָיו\hdot
שֶׁכָּל הֶאָמוּר בְּחוּמָשׁ הַפְּקוּדִים\hdot
קָרַב בַּמִּדְבָּר\hdot
וְכָל הֶאָמוּר בְּתוֹרַת כֹּהֲנִים לֹא קָרַב בַּמִּדְבָּר\hdot
כְּשֶׁבָּאוּ לָאָרֶץ קָֽרְבוּ אֵלּוּ וָאֵלּוּ\hdot
וּמִפְּנֵי מָה אֲנִי אוֹמֵר יִקָּֽרְבוּ כְבָשִׂים בְּלֹא לֶחֶם\hdot
שֶׁהַכְּבָשִׂים מַתִּירִין אֶת עַצְמָן בְּלֹא לֶחֶם\hdot
לֶחֶם בְּלֹא כְבָשִׂים אֵין לִי מִי יַתִּירֶנּוּ׃}

\halacha{הַתְּמִידִין אֵינָן מְעַכְּבִין אֶת הַמּוּסָפִים\hdot
וְלֹא הַמּוּסָפִים מְעַכְּבִין אֶת הַתְּמִידִים\hdot
וְלֹא הַמּוּסָפִין מְעַכְּבִין זֶה אֶת זֶה\hdot
לֹא הִקְרִיבוּ כֶבֶשׂ בַּבּוֹקֶר יַקְרִיבוּ בֵּין הָעַרְבָּיִם\hdot
אָמַר רִבִּי שִׁמְעוֹן אֵימָתַי בִּזְמָן שֶׁהָיוּ אֲנוּסִין אוֹ שׁוֹגְגִין\hdot
אֲבָל אִם הָיוּ מְזִידִים וְלֹא הִקְרִיבוּ כֶּבֶשׂ בַּבֹּקֶר לֹא יַקְרִיבוּ בֵּין הָעַרְבָּיִם\hdot
לֹא הִקְטִירוּ קְטֹרֶת בַּבֹּקֶר יַקְטִירוּ בֵּין הָעַרְבָּיִם\hdot
אָמַר רִבִּי שִׁמְעוֹן וְכֻלָּהּ הָֽיְתָה קְרֵבָה בֵּין הָעַרְבָּיִם\hdot
שֶׁאֵין מְחַנְּכִין אֶת מִזְבַּח הַזָּהָב אֶלָּא בִּקְטֹרֶת הַסַּמִּים\hdot
(שֶׁל בֵּין הָעַרְבָּיִם)\hdot
וְלֹא מִזְבַּח הָעוֹלָה אֶלָּא בְּתָמִיד שֶׁל שַׁחַר\hdot
וְלֹא אֶת הַשֻּׁלְחָן אֶלָּא בְּלֶחֶם הַפָּנִים בְּשַׁבָּת\hdot
וְלֹא אֶת הַמְּנוֹרָה אֶלָּא בְּשִׁבְעָה נֵרוֹתֶיהָ בֵּין הָעַרְבָּיִם׃}

\halacha{חֲבִיתֵי כֹהֵן גָּדוֹל לֹא הָיוּ בָאוֹת חֶצְיָם\hdot
אֶלָּא מֵבִיא עִשָּׂרוֹן שָׁלֵם וְחוֹצֵהוּ\hdot
וּמַקְרִיב מֶחְצָה בַבֹּקֶר וּמֶחְצָה בֵּין הָעַרְבָּיִם\hdot
וְכֹהֵן שֶׁהִקְרִיב מֶחְצָה בְּשַׁחֲרִית וּמֵת\hdot
וּמִנּוּ כֹהֵן אַחֵר תַּחְתָּיו\hdot
לֹא יָבִיא חֲצִי עִשָּׂרוֹן מִבֵּיתוֹ\hdot
וְלֹא חֲצִי עֶשְׂרוֹנוֹ שֶׁל רִאשׁוֹן\hdot
אֶלָּא מֵבִיא עִשָׂרוֹן שָׁלֵם וְחוֹצֵהוּ\hdot
וּמַקְרִיב מֶחְצָה וּמֶחְצָה אָבֵד\hdot
נִמְצְאוּ שְׁנֵי חֲצָיִם קְרֵבִין\hdot
וּשְׁנֵי חֲצָיִם אוֹבְדִין\hdot
לֹא מִנּוּ כֹהֵן אַחֵר\hdot
מִשֶּׁל מִי הָֽיְתָה קְרֵבָה\hdot
רִבִּי שִׁמְעוֹן אוֹמֵר מִשֶּׁל צִבּוּר\hdot
 רִבִּי יְהוּדָה אוֹמֵר מִשֶּׁל יוֹרְשִׁים\hdot
וּשְׁלֵימָה הָֽיְתָה קְרֵבָה׃}

\chead{יום יט}

\perek{כָּל}{הַמְּנָחוֹת בָּאוֹת מַצָּה\hdot
חוּץ מֵחָמֵץ שֶׁבַּתּוֹדָה וּשְׁתֵּי הַלֶּחֶם\hdot
שֶׁהֵן בָּאוֹת חָמֵץ\hdot
רִבִּי מֵאִיר אוֹמֵר שְׂאוֹר בּוֹדֶה לָהֶן מִתּוֹכָן וּמְחַמְּצָן\hdot
רִבִּי יְהוּדָה אוֹמֵר אַף הִיא אֵינָהּ מִן הַמֻּבְחָר\hdot
אֶלָּא מֵבִיא אֶת הַשְּׂאוֹר\hdot
וְנוֹתֵן לְתוֹךְ הַמִּדָּה\hdot
וּמְמַלֵּא אֶת הַמִּדָּה\hdot
אָֽמְרוּ לוֹ אַף הִיא הָֽיְתָה חֲסֵרָה אוֹ יְתֵירָה׃}

\halacha{כָּל הַמְּנָחוֹת נִלּוֹשׁוֹת בְּפוֹשְׁרִין\hdot
וּמְשַׁמְּרָן שֶׁלֹּא יַחְמִיצוּ\hdot
וְאִם הֶחְמִיצוּ שִׁיְּרֶיהָ\hdot
עוֹבֵר בְּלֹא תַעֲשֶׂה\hdot
שֶׁנֶּאֱמַר כָּל הַמִּנְחָה אֲשֶׁר תַּקְרִיבוּ לַיְיָ לֹא תֵעָשֶׂה חָמֵץ\hdot
וחַיַּבִים עַל לִישָׁתָהּ וְעַל עֲרִיכָתָהּ וְעַל אֲפִיָּתָהּ׃}

\halacha{יֵשׁ טְעוּנוֹת שֶׁמֶן וּלְבוֹנָה\hdot
שֶׁמֶן וְלֹא לְבוֹנָה\hdot
לְבוֹנָה וְלֹא שֶׁמֶן\hdot
לֹא שֶׁמֶן וְלֹא לְבוֹנָה\hdot
וְאֵלּוּ טְעוּנוֹת שֶׁמֶן וּלְבוֹנָה\hdot
מִנְחַת הַסֹּלֶת וְהַמַּחֲבַת וְהַמַּרְחֶשֶׁת וְהַחַלּוֹת וְהָֽרְקִיקִין\hdot
מִנְחַת כֹּהֲנִים וּמִנְחַת כֹּהֵן מָשִׁיחַ\hdot
וּמִנְחַת גּוֹיִם\hdot
וּמִנְחַת נָשִׁים\hdot
וּמִנְחַת הָעֹמֶר\hdot
מִנְחַת נְסָכִין טְעוּנָה שֶׁמֶן וְאֵינָהּ טְעוּנָה לְבוֹנָה\hdot
לֶחֶם הַפָּנִים טָעוּן לְבוֹנָה וְאֵין טָעוּן שֶׁמֶן\hdot
שְׁתֵּי הַלֶּחֶם\hdot
וּמִנְחַת חוֹטֵא\hdot
וּמִנְחַת קְנָאוֹת\hdot
לֹא שֶׁמֶן וְלֹא לְבוֹנָה׃}

\halacha{וְחַיָּב עַל הַשֶּׁמֶן בִּפְנֵי עַצְמוֹ\hdot
וְעַל הַלְּבוֹנָה בִּפְנֵי עַצְמָהּ\hdot
נָתַן עָלֶיהָ שֶׁמֶן פְּסָלָהּ\hdot
לְבוֹנָה יִלְקְטֶנָּה\hdot
נָתַן שֶׁמֶן עַל שִׁיְּרֶיהָ\hdot
אֵינוֹ עוֹבֵר בְּלֹא תַעֲשֶׂה\hdot
נָתַן כֶּלִי עַל גַּבֵּי כֶלִי לֹא פְסָלָהּ׃}

\halacha{יֵשׁ טְעוּנוֹת הַגָּשָׁה וְאֵינָן טְעוּנוֹת תְּנוּפָה\hdot
תְּנוּפָה וְלֹא הַגָּשָׁה\hdot
הַגָּשָׁה וּתְנוּפָה\hdot
לֹא תְנוּפָה וְלֹא הַגָּשָׁה\hdot
אֵלּוּ טְעוּנוֹת הַגָּשָׁה וְאֵינָן טְעוּנוֹת תְּנוּפָה\hdot
(מִנְחַת) הַסֹּלֶת וְהַמַּחֲבַת וְהַמַּרְחֶשֶׁת וְהַחַלּוֹת וְהָֽרְקִיקִין\hdot
מִנְחַת כֹּהֲנִים\hdot
מִנְחַת כֹּהֵן מָשִׁיחַ\hdot
מִנְחַת גּוֹיִם מִנְחַת נָשִׁים וּמִנְחַת חוֹטֵא\hdot
רִבִּי שִׁמְעוֹן אוֹמֵר מִנְחַת כֹּהֲנִים וּמִנְחַת כֹּהֵן מָשִׁיחַ\hdot
אֵין בָּהֶן הַגָּשָׁה\hdot
מִפְּנֵי שֶׁאֵין (בָּהֶן) קְמִיצָה\hdot
וְכָל שֶׁאֵין (בָּהֶן) קְמִיצָה\hdot
אֵין בָּהֶן הַגָּשָׁה׃}

\halacha{אֵלּוּ טְעוּנִין תְּנוּפָה וְאֵין טְעוּנִין הַגָּשָׁה\hdot
לוֹג שֶׁמֶן שֶׂל מְצוֹרָע וַאֲשָׁמוֹ\hdot
וְהַבִּכּוּרִים\hdot
כְּדִבְרֵי רִבִּי אֱלִיעֶזֶר בֶּן יַעֲקֹב\hdot
וְאֵמוּרֵי שַׁלְמֵי יָחִיד\hdot
וְחָזֶה וְשׁוֹק שֶׁלָהֶן\hdot
אֶחָד אֲנָשִׁים וְאֶחָד נָשִׁים בְּיִשְׂרָאֵל\hdot
אֲבָל לֹא בָאֲחֵרִים\hdot
וּשְׁתֵּי הַלֶּחֶם וּשְׁנֵי כִבְשֵׂי עֲצֶרֶת\hdot
כֵּיצַד הוּא עוֹשֶׂה\hdot
נוֹתֵן שְׁתֵּי הַלֶּחֶם עַל גַּבֵּי שְׁנֵי כְבָשִׁים\hdot
וּמַנִּיחַ שְׁתֵּי יָדָיו מִלְּמַטָּן\hdot
מוֹלִיךְ וּמֵבִיא\hdot
מַעֲלֶה וּמוֹרִיד\hdot
שֶׁנֶּאֱמַר אֲשֶׁר הוּנָף וַאֲשֶׁר הוּרָם\hdot
תְּנוּפָה הָֽיְתָה בַמִּזְרַח\hdot
וְהַגָּשָׁה בַמַּעֲרָב\hdot
וּתְנוּפוֹת קוֹדְמוֹת לַהַגָּשׁוֹת\hdot
מִנְחַת הָעֹמֶר\hdot
 וּמִנְחַת קְנָאוֹת\hdot
טְעוּנוֹת תְּנוּפָה וְהַגָּשָׁה\hdot
לֶחֶם הַפָּנִים וּמִנְחַת נְסָכִים\hdot
לֹא תְּנוּפָה וְלֹא הַגָּשָׁה׃}

\halacha{רִבִּי שִׁמְעוֹן אוֹמֵר שְׁלֹשָׁה מִינִים טְעוּנִים שָׁלֹשׁ מִצְוֺת\hdot
שְׁתַּיִם בְּכָל אֶחָד וְאֶחָד\hdot
וְהַשְּׁלִישִׁית אֵין בָּהֶן\hdot
וְאֵלּוּ הֵן זִבְחֵי שַׁלְמֵי יָחִיד\hdot
וְזִבְחֵי שַׁלְמֵי צִבּוּר\hdot
וַאֲשַׁם מְצוֹרָע\hdot
זִבְחֵי שַׁלְמֵי יָחִיד\hdot
טְעוּנִים סְמִיכָה חַיִּים\hdot
וּתְנוּפָה שְׁחוּטִים\hdot
וְאֵין בָּהֶן תְּנוּפָה חַיִּים\hdot
זִבְחֵי שַׁלְמֵי צִבּוּר טְעוּנִים תְּנוּפָה חַיִּים וּשְׁחוּטִים\hdot
וְאֵין בָּהֶן סְמִיכָה\hdot
וַאֲשַׁם מְצוֹרָע טָעוּן סְמִיכָה וּתְנוּפָה חַי\hdot
וְאֵין בּוֹ תְנוּפָה שָׁחוּט׃}

\halacha{הָאוֹמֵר הֲרֵי עָלַי בְּמַחֲבַת לֹא יָבִיא בְמַרְחֶשֶׁת\hdot
בְּמַרְחֶשֶׁת לֹא יָבִיא בְמַחֲבַת\hdot
וּמַה בֵּין מַחֲבַת לְמַרְחֶשֶׁת\hdot
אֶלָּא שֶׁהַמַּרְחֶשֶׁת יֵשׁ לָהּ כִּיסּוּי\hdot
וְהַמַּחֲבַת אֵין לָהּ כִּיסּוּי\hdot
דִּבְרֵי רִבִּי יוֹסֵי הַגָּלִילִי\hdot
רִבִּי חֲנַנְיָה בֶּן גַּמְלִיאֵל אוֹמֵר מַרְחֶשֶׁת עֲמוּקָה\hdot
וּמַעֲשֶׂיהָ רוֹחֲשִׁים\hdot
וּמַחֲבַת צָפָה\hdot
וּמַעֲשֶׂיהָ קָשִׁים׃}

\halacha{הָאוֹמֵר הֲרֵי עָלַי בְּתַנּוּר\hdot
לֹא יָבִיא מַאֲפֵה כֻפָּח\hdot
וּמַאֲפֵה רְעָפִים\hdot
וּמַאֲפֵה יוֹרוֹת הָעֲרָבִיִּים\hdot
רִבִּי יְהוּדָה אוֹמֵר אִם רָצָה יָבִיא מַאֲפֵה כֻפָּח\hdot
הֲרֵי עָלַי מִנְחַת מַאֲפֶה\hdot
לֹא יָבִיא מֶחְצָה חַלּוֹת וּמֶחְצָה רְקִיקִין\hdot
רִבִּי שִׁמְעוֹן מַתִּיר מִפְּנֵי שֶׁהוּא קָרְבָּן אֶחָד׃}

\perek{אֵלּוּ}{מְנָחוֹת נִקְמָצוֹת וְשִׁיְּרֵיהֶן לַכֹּהֲנִים\hdot
מִנְחַת סֹלֶת\hdot
וְהַמַּחֲבַת\hdot
וְהַמַּרְחֶשֶׁת\hdot
וְהַחַלּוֹת\hdot
וְהָֽרְקִיקִין\hdot
מִנְחַת גּוֹיִם\hdot
מִנְחַת נָשִׁים\hdot
מִנְחַת הָעֹמֶר\hdot
מִנְחַת חוֹטֵא וּמִנְחַת קְנָאוֹת\hdot
רִבִּי שִׁמְעוֹן אוֹמֵר מִנְחַת חוֹטֵא שֶׁל כֹּהֲנִים נִקְמֶצֶת\hdot
וְהַקֹּמֶץ קָרֵב לְעַצְמוֹ\hdot
וְהַשִּׁרַיִם קְרֵבִין לְעַצְמָן׃}

\halacha{מִנְחַת כֹּהֲנִים\hdot
וּמִנְחַת כֹּהֵן מָשִׁיחַ\hdot
וּמִנְחַת נְסָכִים לַמִּזְבֵּחַ\hdot
וְאֵין בָּהֶם לַכֹּהֲנִים\hdot
בָּזֶה יָפֶה כֹּחַ הַמִּזְבֵּח מִכֹּחַ הַכֹּהֲנִים\hdot
שְׁתֵּי הַלֶּחֶם וְלֶחֶם הַפָּנִים לַכֹּהֲנִים\hdot
וְאֵין בָּהֶם לַמִּזְבֵּחַ\hdot
וּבָזֶה יָפֶה כֹּחַ הַכֹּהֲנִים מִכֹּחַ הַמִּזְבֵּחַ׃}

\halacha{כָּל הַמְּנָחוֹת הַנַּעֲשׂוֹת בַּכֶּלִי\hdot
טְעוּנוֹת שָׁלֹשׁ מַתְּנוֹת שֶׁמֶן\hdot
יְצִיקָה וּבְלִילָה\hdot
וּמַתַּן שֶׁמֶן בַּכֶּלִי קוֹדֵם לַעֲשִׂיָּתָן\hdot
וְהַחַלּוֹת (בּוֹלְלָן) דִּבְרֵי רִבִּי וַחֲכָמִים אוֹמְרִים (סֹלֶת) הַחַלּוֹת טְעוּנוֹת בְּלִילָה\hdot
הָֽרְקִיקִין מְשִׁיחָה\hdot
כֵּיצַד מוֹשְׁחִין כְּמִין כִּי\hdot
וּשְׁאַר הַשֶּׁמֶן נֶאֱכָל לַכֹּהֲנִים׃}

\halacha{כָּל הַמְּנָחוֹת הַנַּעֲשׂוֹת בַּכֶּלִי טְעוּנוֹת פְּתִיתָה\hdot
מִנְחַת יִשְׂרָאֵל כּוֹפֵל אַחַת לִשְׁנַיִם\hdot
וּשְׁנַיִם לְאַרְבָּעָה וּמַבְדִּיל\hdot
מִנְחַת כֹּהֲנִים כּוֹפֵל אַחַת לִשְׁנַיִם\hdot
וּשְׁנַיִם לְאַרְבָּעָה וְאֵינוֹ מַבְדִּיל\hdot
מִנְחַת כֹּהֵן מָשִׁיחַ לֹא הָיָה מְכַפְּלָהּ\hdot
רִבִּי שִׁמְעוֹן אוֹמֵר מִנְחַת כֹּהֲנִים\hdot
וּמִנְחַת כֹּהֵן מָשִׁיחַ\hdot
אֵין בָּהֶן פְּתִיתָה\hdot
מִפְּנֵי שֶׁאֵין בָּהֶן קְמִיצָה\hdot
וְכָל שֶׁאֵין בָּהֶן קְמִיצָה\hdot
אֵין בָּהֶן פְּתִיתָה\hdot
וְכֻלָּן (פְּתִיתָתָן) כְּזֵיתִים׃}

\halacha{כָּל הַמְּנָחוֹת טְעוּנוֹת שָׁלֹשׁ מֵאוֹת שִׁיפָה\hdot
וַחֲמֵשׁ מֵאוֹת בְּעִיטָה\hdot
וְהַשִּׁיפָה וְהַבְּעִיטָה\hdot
בַּחִטִּים\hdot
רִבִּי יוֹסֵי אוֹמֵר (אַף) בַּבָּצֵק\hdot
כָּל הַמְּנָחוֹת בָּאוֹת עֶשֶׂר עֶשֶׂר\hdot
חוּץ מִלֶּחֶם הַפָּנִים\hdot
וַחֲבִיתֵי כֹהֵן גָּדוֹל\hdot
שֶׁהֵם בָּאוֹת שְׁתֵּים עֶשְׂרֵה\hdot
דִּבְרֵי רִבִּי יְהוּדָה\hdot
רִבִּי מֵאִיר אוֹמֵר כֻּלָּן בָּאוֹת שְׁתֵּים עֶשְׂרֵה\hdot
חוּץ מֵחַלּוֹת תּוֹדָה וְהַנְּזִירוּת שֶׁהֵן בָּאוֹת עֶשֶׂר עֶשֶׂר׃}

\halacha{הָעֹמֶר הָיָה בָא עִשָּׂרוֹן מִשָּׁלֹשׁ סְאִין\hdot
שְׁתֵּי הַלֶּחֶם שְׁנֵי עֶשְׂרוֹנִים מִשָּׁלֹשׁ סְאִין\hdot
לֶחֶם הַפָּנִים עֶשְׂרִים וְאַרְבָּעָה עֶשְׂרוֹנִים\hdot
מֵעֶשְׂרִים וְאַרְבָּעָה סְאִין׃}

\halacha{הָעֹמֶר הָיָה מְנֻפֶּה בִּשְׁלֹשׁ עֶשְׂרֵה נָפָה\hdot
וּשְׁתֵּי הַלֶּחֶם בִּשְׁתֵּים עֶשְׂרֵה\hdot
וְלֶחֶם הַפָּנִים בְּאַחַת עֶשְׂרֵה\hdot
רִבִּי שִׁמְעוֹן אוֹמֵר לֹא הָיָה לָהּ קִצְבָּה\hdot
אֶלָּא סֹלֶת מְנֻפָּה כָּל צָרְכָהּ הָיָה מֵבִיא\hdot
שֶׁנֶּאֱמַר וְלָקַחְתָּ סֹלֶת וְאָפִיתָ אוֹתָהּ\hdot
עַד שֶׁתְּהֵא מְנֻפָּה כָּל צָרְכָהּ׃}

\perek{הַתּוֹדָה}{הָֽיְתָה בָאָה חָמֵשׁ סְאִין יְרוּשַׁלְמִיּוֹת\hdot
שֶׁהֵן שֵׁשׁ מִדְבָּרִיּוֹת\hdot
שְׁתֵּי אֵיפוֹת\hdot
הָאֵיפָה שָׁלֹשׁ סְאִין\hdot
עֶשְׂרִים עִשָּׂרוֹן עֲשָׂרָה לְחָמֵץ\hdot
(עֲשָׂרָה לְמַצָּה) עֲשָׂרָה לְחָמֵץ עִשָּׂרוֹן לְחַלָּה\hdot
וַעֲשָׂרָה לְמַצָּה\hdot
וּבַמַּצָּה שְׁלֹשָׁה מִינִין\hdot
חַלּוֹת וּרְקִיקִין\hdot
וּרְבוּכָה\hdot
נִמְצְאוּ שְׁלֹשָׁה עֶשְׂרוֹנוֹת וְשָׁלִישׁ לְכָל מִין\hdot
שָׁלֹשׁ חַלָּוֹת לְעִשָּׂרוֹן\hdot
בַּמִּדָּה יְרוּשַׁלְמִית הָיוּ שְׁלֹשִׁים קַב\hdot
חֲמִשָּׁה עָשָׂר לְחָמֵץ וַחֲמִשָּׁה עָשָׂר לְמַצָּה\hdot
חֲמִשָּׁה עָשָׂר לְחָמֵץ\hdot
קַב וָחֵצִי לַחַלָּה\hdot
וַחֲמִשָּׁה עָשָׂר לְמַצָּה\hdot
וְהַמַּצָּה שְׁלֹשָׁה מִינִין\hdot
חַלּוֹת וּרְקִיקִין\hdot
וּרְבוּכָה\hdot
נִמְצְאוּ חֲמֵשֶׁת קַבִּין לְכָל מִין\hdot
שְׁתֵּי חַלּוֹת לַקַּב׃}

\halacha{הַמִּלּוּאִים הָיוּ בָאִים כְּמַצָּה שֶׁבַּתּוֹדָה\hdot
חַלּוֹת\hdot
וּרְקִיקִים\hdot
וּרְבוּכָה\hdot
הַנְּזִירוּת הָֽיְתָה בָאָה שְׁתֵּי יָדֹות כְּמַצָּה שֶׁבַּתּוֹדָה\hdot
חַלּוֹת וּרְקִיקִים\hdot
וְאֵין בָּהּ רְבוּכָה\hdot
(נִמְצְאוּ) עֲשֶׂרֶת קַבִּים יְרוּשַׁלְמִיּוֹת\hdot
שֶׁהֵן שִׁשָּׁה עֶשְׂרוֹנוֹת וְעִדוּיָן\hdot
וּמִכֻּלָּן הָיָה נוֹטֵל אַחַד מֵעֲשָׂרָה תְּרוּמָה\hdot
שֶׁנֶּאֱמַר וְהִקְרִיב מִמֶּנּוּ אֶחָד מִכָּל קָרְבָּן תְּרוּמָה לַיְיָ\hdot
אֶחָד שֶׁלֹּא יִטּוֹל פָּרוּס\hdot
מִכָּל קָרְבָּן שֶׁיִּהְיוּ כָל הַקָּרְבָּנוֹת שָׁוִין\hdot
וְשֶׁלֹּא יִטּוֹל מִקָּרְבָּן לַחֲבֵירוֹ\hdot
לַכֹּהֵן הַזּוֹרֵק אֶת דַּם הַשְּׁלָמִים לוֹ יִהְיֶה\hdot
וְהַשְּׁאָר נֶאֱכַל לַבְּעָלִים׃}

\halacha{הַשּׁוֹחֵט אֶת הַתּוֹדָה בִּפְנִים וְלַחְמָהּ חוּץ לְחוֹמָה\hdot
לֹא קִדֵּשׁ הַלֶּחֶם\hdot
שְׁחָטָהּ עַד שֶׁלֹּא קָֽרְמוּ בַתַּנּוּר\hdot
וַאֲפִילוּ קָֽרְמוּ כֻלָּן\hdot
חוּץ מֵאֶחָד מֵהֶן\hdot
לֹא קִדֵּשׁ הַלֶּחֶם\hdot
שְׁחָטָהּ חוּץ לִזְמָנָהּ וְחוּץ לִמְקוֹמָהּ\hdot
קִדֵּשׁ הַלֶּחֶם\hdot
שְׁחָטָהּ וְנִמְצֵאת טְרֵיפָה
לֹא קִדֵּשׁ הַלֶּחֶם\hdot
שְׁחָטָהּ וְנִמְצֵאת בַּעֲלַת מוּם\hdot
רִבִּי אֱלִיעֶזֶר אוֹמֵר קִדֵּשׁ\hdot
וַחֲכָמִים אוֹמְרִים לֹא קִדֵּשׁ\hdot
שְׁחָטָהּ שֶׁלֹּא לִשְׁמָהּ\hdot
וְכֵן אֵיל הַמִּלּוּאִים\hdot
וְכֵן שְׁנֵי כִבְשֵׂי עֲצֶרֶת\hdot
שֶׁשְּׁחָטָן שֶׁלֹּא לִשְׁמָן\hdot
לֹא קִדֵּשׁ הַלֶּחֶם׃}

\halacha{נְסָכִין שֶׁקִּדְּשׁוּ בַכֶּלִי\hdot
וְנִמְצָא הַזֶּבַח פָּסוּל\hdot
אִם יֵשׁ שָׁם זֶבַח אַחֵר יִקָּֽרְבוּ עִמּוֹ\hdot
וְאִם לָאוּ יִפָּֽסְלוּ בְלִינָה\hdot
וָלֶד תּוֹדָה וּתְמוּרָתָהּ\hdot
וְהַמַּפְרִישׁ תּוֹדָתוֹ וְנֶאֱבָדָה\hdot
וְהִפְרִישׁ אַחֶרֶת תַּחְתֶּיהָ\hdot
אֵינָן טְעוּנִים לֶחֶם שֶׁנֶּאֱמַר וְהִקְרִיב עַל זֶּבַח הַתּוֹדָה\hdot
הַתּוֹדָה טְעוּנָה לֶחֶם\hdot
וְלֹא וַלְדָהּ\hdot
וְלֹא חֲלִיפָתָהּ\hdot
וְלֹא תְמוּרָתָהּ\hdot
טְעוּנִין לֶחֶם׃}

\halacha{הָאוֹמֵר הֲרֵי עָלַי תּוֹדָה\hdot
יָבִיא הִיא וְלַחְמָהּ מִן הַחוֹלִין\hdot
תּוֹדָה מִן הַחוֹלִין וְלַחְמָהּ מִן הַמַּעֲשֵׂר\hdot
יָבִיא לַחְמָהּ מִן הַחוֹלִין\hdot
תּוֹדָה מִן הַמַּעֲשֵׂר\hdot
וְלַחְמָהּ מִן הַחוֹלִין\hdot
יָבִיא\hdot
הַתּוֹדָה הִיא וְלַחְמָהּ מִן הַמַּעֲשֵׂר\hdot
יָבִיא\hdot
וְלִא יָבִיא מֵחִטֵּי מַעֲשֵׂר שֵׁנִי\hdot
אֶלָּא מִמָּעוֹת מַעֲשֵׂר שֵׁנִי׃}

\halacha{וּמְנַיִן לָאוֹמֵר הֲרֵי עָלַי תּוֹדָה לֹא יָבִיא אֶלָּא מִן הַחוֹלִין\hdot
שֶׁנֶּאֱמַר וְזָבַחְתָּ פֶּסַח לַיְיָ אֱלֹהֶיךָ צֹאן וּבָקָר\hdot
וַהֲלֹא אֵין פֶּסַח בָּא אֶלָּא מִן הַכְּבָשִׂים וּמִן הָעִזִּים\hdot
אִם כֵּן לָמָּה נֶאֱמַר צֹאן וּבָקָר\hdot
אֶלָּא לְהָקִישׁ כָּל הַבָּא מִן הַצֹּאן וּמִן הַבָּקָר\hdot
לַפֶּסַח\hdot
מַה הַפֶּסַח שֶׁהוּא בָא בְּחוֹבָה\hdot
אֵינוֹ בָא אֶלָּא מִן הַחוֹלִין\hdot
אַף כָּל דָּבָר שֶׁהוּא בָא בְּחוֹבָה\hdot
לֹא יָבֹא אֶלָּא מִן הַחוֹלִין\hdot
(לְפִיכַךְ) הָאוֹמֵר הֲרֵי עָלַי תּוֹדָה הֲרֵי עָלַי שְׁלָמִים\hdot
הוֹאִיל וְהֵם בָּאִים חוֹבָה\hdot
לֹא יָבִיאוּ אֶלָּא מִן הַחוֹלִין\hdot
וְהַנְּסָכִים בְּכָל מָקוֹם\hdot
לֹא יָבִיאוּ אֶלָּא מִן הַחוֹלִין׃}

\perek{כָּל}{קָרְבְּנוֹת הַצִּבּוּר וְהַיָּחִיד\hdot
בָּאִים מִן הָאָרֶץ וּמֵחוּצָה לָאָרֶץ\hdot
מִן הֶחָדָשׁ וּמִן הַיָּשָׁן\hdot
חוּץ מִן הָעֹמֶר וּשְׁתֵּי הַלֶּחֶם\hdot
שֶׁאֵינָן בָּאִים אֶלָּא מִן הֶחָדָשׁ וּמִן הָאָרֶץ\hdot
וְכֻלָּן אֵינָן בָּאִין אֶלָּא מִן הַמֻּבְחָר\hdot
וְאֵיזֶהוּ מֻבְחָר\hdot
מִכְּמָס\hdot
וּמִזּוֹנִיחָה\hdot
אַלְפָא לַסֹּלֶת\hdot
שְׁנִיָּה לָהֶם חָפְרַיִם בַּבִּקְעָה\hdot
כָּל הָאֲרָצוֹת הָיוּ כְשֵׁרוֹת\hdot
אֶלָּא מִכָּאן הָיוּ מְבִיאִים׃}

\halacha{אֵין מְבִיאִין לֹא מִבֵּית הַזְּבָלִים\hdot
וְלֹא מִבֵּית הַשְּׁלָחִים\hdot
וְלֹא מִבֵּית הָאִילָן\hdot
וְאִם הֵבִיא כָּשֵׁר\hdot
כֵּיצַד הוּא עוֹשֶׂה\hdot
נִרָהּ שָׁנָה רִאשׁוֹנָה\hdot
וּבַשְּׁנִיָּה זוֹרְעָהּ קֹדֶם לַפֶּסַח שִׁבְעִים יוֹם\hdot
וְהִיא עוֹשָׂה סֹלֶת מְרֻבָּה\hdot
כֵּיצַד הוּא בוֹדֵק\hdot
הַגִּזְבָּר מַכְנִיס אֶת יָדוֹ לְתוֹכָהּ\hdot
עָלָה בָהּ אָבָק פְּסוּלָה\hdot
עַד שֶׁיְּנִיפֶנָּה\hdot
וְאִם הִתְלִיעָה פְּסוּלָה׃}

\halacha{תְּקוֹעָה אַלְפָא לֶשֶּׁמֶן\hdot
אַבָּא שָׁאוּל אוֹמֵר שְׁנִיָּה לָהּ רְגָב בְּעֵבֶר הַיַּרְדֵּן\hdot
כָּל הָאֲרָצוֹת הָיוּ כְשֵׁרוֹת\hdot
אֶלָּא מִכָּאן הָיוּ מְבִיאִים\hdot
אֵין מְבִיאִין לֹא מִבֵּית הַזְּבָלִים\hdot
וְלֹא מִבֵּית הַשְּׁלָחִים\hdot
וְלֹא מִמָּה שֶׁנִּזְרַע בֵּינֵיהֶם\hdot
וְאִם הֵבִיא כָּשֵׁר\hdot
אֵין מְבִיאִין אַנְפִיקָנוֹן\hdot
וְאִם הֵבִיא פָּסוּל\hdot
אֵין מְבִיאִין מִן הַגַּרְגְּרִים שֶׁנָּֽשְׁרוּ בַמַּיִם\hdot
וְלֹא מִן הַכְּבוּשִׁים\hdot
וְלֹא מִן הַשְּׁלוּקִין\hdot
וְאִם הֵבִיא פָּסוּל׃}

%\tractatehead{מסכת חולין}{holin}

\perek{הַכֹּל}{שׁוֹחֲטִין וּשְׁחִיטָתָן כְּשֵׁרָה\hdot
חוּץ מֵחֵרֵשׁ שׁוֹטֶה וְקָטָן שֶׁמֶא יְקַלְקְלו בִשְׁחִיטָתָן\hdot
וְכֻלָּן שֶׁשָּׁחֲטוּ\hdot
וְאַחֵרִים רוֹאִין אוֹתָן\hdot
שְׁחִיטָתָן כְּשֵׁרָה\hdot
שְׁחִיטַת נָכְרִי נְבֵילָה\hdot
וּמְטַמְּאָה בְמַשָּׁא\hdot
הַשּׁוֹחֵט בַּלַּיְלָה\hdot
וְכֵן הַסּוֹמֵא שֶׁשָּׁחַט\hdot
שְׁחִיטָתוֹ כְּשֵׁרָה\hdot
הַשּׁוֹחֵט בְּשַׁבָּת וּבְיוֹם הַכִּפּוּרִים\hdot
אַף עַל פִּי שְׁמִּתַחַיֵּב בְּנַפְשׁוֹ\hdot
שְׁחִיטָתוֹ כְּשֵׁרָה׃}

\halacha{הַשּׁוֹחֵט בְּמַגַּל יָד\hdot
בַּצּוּר וּבַקָּנֶה שְׁחִיטָתוֹ כְּשֵׁרָה\hdot
הַכֹּל שׁוֹחֲטִין\hdot
וּלְעוֹלָם שׁוֹחֲטִין\hdot
וּבַכֹּל שׁוֹחֲטִין\hdot
חוּץ מִמַּגַּל קָצִיר וְהַמְּגֵירָה וְהַשִּׁנַּיִם וְהַצִּפֹּרֶן\hdot
מִפְּנֵי שֶׁהֵן חוֹנְקִין\hdot
הַשּׁוֹחֵט בְּמַגַּל קָצִיר בְּדֶרֶךְ הֲלִיכָתָהּ\hdot
בֵּית שַׁמָּאי פּוֹסְלִין\hdot
וּבֵית הִלֵּל מַכְשִׁירִין\hdot
וְאִם הֶחֱלִיקוּ שִׁנֶּיהָ\hdot
הֲרֵי הִיא כַּסַּכִּין׃}

\halacha{הַשּׁוֹחֵט מִתּוֹךְ הַטַּבַּעַת וְשִׁיֵּר בָּהּ מְלֹא הַחוּט עַל פְּנֵי כֻלָּהּ\hdot
שְׁחִיטָתוֹ כְּשֵׁרָה\hdot
רִבִּי יוֹסֵי בַּר יְהוּדָה אוֹמֵר מְלֹא הַחוּט עַל פְּנֵי רֻבָּהּ׃}

\halacha{הַשּׁוֹחֵט מִן הַצְּדָדִין\hdot
שְׁחִיטָתוֹ כְּשֵׁרָה\hdot
הַמּוֹלֵק מִן הַצְּדָדִין\hdot
מְלִיקָתוֹ פְּסוּלָה\hdot
הַשּׁוֹחֵט מִן הָעֹרֶף\hdot
שְׁחִיטָתוֹ פְּסוּלָה\hdot
הַמּוֹלֵק מִן הַעֹרֶף מְלִיקָתוֹ כְּשֵׁרָה\hdot
הַשּׁוֹחֵט מִן הַצַּוָּאר\hdot
שְׁחִיטָתוֹ כְּשֵׁרָה\hdot
הַמּוֹלֵק מִן הַצַּוָאר\hdot
מְלִיקָתוֹ פְּסוּלָה\hdot
שֶׁכָּל הָעֹרֶף כָּשֵׁר לַמְּלִיקָה\hdot
וְכָל הַצַּוָּאר כָּשֵׁר לַשְּׁחִיטָה\hdot
נִמְצָא כָּשֵׁר בַּשְּׁחִיטָה\hdot
פָּסוּל בַּמְּלִיקָה\hdot
כָּשֵׁר בַּמְּלִיקָה\hdot
פָּסוּל בַּשְּׁחִיטָה׃}

\halacha{כָּשֵׁר בַּתּוֹרִין\hdot
פָּסוּל בִּבְנֵי יוֹנָה\hdot
כָּשֵׁר בִּבְנֵי יוֹנָה\hdot
פָּסוּל בַּתּוֹרִין\hdot
תְּחִלַּת הַצִּיהוּב בְּזֶה וּבְזֶה פָּסוּל׃}

\halacha{כָּשֵׁר בַּפָּרָה פָּסוּל בָּעֲגָלָה\hdot
כָּשֵׁר בָּעֲגָלָה פָּסוּל בַּפָּרָה\hdot
כָּשֵׁר בַּכֹּהֲנִים פָּסוּל בַּלְוִיִּם\hdot
כָּשֵׁר בַּלְוִיּם פָּסוּל בַּכֹּהֲנִים\hdot
טָהוֹר בִּכְלֵי חֶרֶשׂ טָמֵא בְּכָל הַכֵּלִים\hdot
טָהוֹר בְּכָל הַכֵּלִים טָמֵא בִּכְלֵי חֶרֶשׂ\hdot
טָהוֹר בִּכְלֵי עֵץ טָמֵא בִּכְלֵי מַתְּכוֹת\hdot
טָהוֹר בִּכְלֵי מַתְּכוֹת טָמֵא בִּכְלֵי עֵץ\hdot
הַחַיָּב בַּשְּׁקֵדִים הַמָּרִים פָּטוּר בַּמְּתוּקִים\hdot
הַחַיָּב בַּמְּתוּקִים פָּטוּר בַּמָּרִים׃}

\halacha{הַתֶּמֶד עַד שֶׁלֹּא הֶחְמִיץ\hdot
אֵינוֹ נִקָּח בְּכֶסֶף מַעֲשֵׂר\hdot
וּפוֹסֵל אֶת הַמִּקְוֶה\hdot
מִשֶּׁהֶחְמִיץ נִקָּח בּכֶסֶף מַעֲשֵׂר\hdot
וְאֵינוֹ פוֹסֵל אֶת הַמִּקְוֶה\hdot
הָאַחִים הַשֻּׁתָּפִים כְּשֶׁחַיָּבִים בַּקָּלְבּוֹן\hdot
פְּטוּרִין מִמַּעְשַׂר בְּהֵמָה\hdot
כְּשֶׁחַיָּבִין בְּמַעְשַׂר בְּהֵמָה\hdot
פְּטוּרִין מִן הַקָּלְבּוֹן\hdot
כָּל מָּקוֹם שֶׁיֵּשׁ מֶכֶר אֵין קְנָס\hdot
וְכָל מָקוֹם שֶׁיֵּשׁ קְנָס אֵין מֶכֶר\hdot
כָּל מָקוֹם שֶׁיֵּשׁ מִיאוּן אֵין חֲלִיצָה\hdot
וְכָל מָקוֹם שֶׁיֵּשׁ חֲלִיצָה אֵין מִיאוּן\hdot
כָּל מָקוֹם שֶׁיֵּשׁ תְּקִיעָה אֵין הַבְדָּלָה\hdot
וְכָל מָקוֹם שֶׁיֵּשׁ הַבְדָּלָה\hdot
אֵין תְּקִיעָה\hdot
יוֹם טוֹב שֶׁחַל לִהְיוֹת בְּעֶרֶב שַׁבָּת\hdot
תּוֹקְעִין וְלֹא מַבְדִּילִין\hdot
בְּמוֹצָאֵי שַׁבָּת מַבְדִּילִין וְלֹא תוֹקְעִין\hdot
כֵּיצַד מַבְדִּילִין\hdot
הַמַּבְדִּיל בֵּין קֹדֶשׁ לְקֹדֶשׁ\hdot
רִבִּי דוֹסָא אוֹמֵר בֵּין קֹדֶשׁ חָמוּר לְקֹדֶשׁ הַקַּל׃}

\perek{הַשּׁוֹחֵט}{אֶחָד בָּעוֹף וּשְׁנַיִם בַּבְּהֵמָה\hdot
שְׁחִיטָתוֹ כְּשֵׁרָה\hdot
וְרֻבּוֹ שֶׁל אֶחָד כָּמוֹהוֹ\hdot
רִבִּי יְהוּדָה אוֹמֵר עַד שֶׁיִּשְׁחוֹט אֶת הַוָּרִידִין\hdot
חֲצִי אֶחָד בָּעוֹף וְאֶחָד וָחֵצִי בַבְּהֵמָה\hdot
שְׁחִיטָתוֹ פְּסוּלָה\hdot
רוֹב אֶחָד בָּעוֹף וְרוֹב שְׁנַיִם בַּבְּהֵמָה\hdot
שְׁחִיטָתוֹ כְּשֵׁרָה׃}

\halacha{הַשּׁוֹחֵט שְׁנֵי רָאשִׁין כְּאֶחָד שְׁחִיטָתוֹ כְּשֵׁרָה\hdot
שְׁנַיִם אוֹחֲזִין בַּסַּכִּין וְשׁוֹחֲטִין\hdot
אֲפִילוּ אֶחָד לְמַעְלָה וְאֶחָד לְמַטָּה\hdot
שְׁחִיטָתָן כְּשֵׁרָה׃}

\halacha{הִתִּיז אֶת הָרֹאשׁ בְּבַת אַחַת פְּסוּלָה\hdot
הָיָה שׁוֹחֵט וְהִתִּיז אֶת הָרֹאשׁ בְּבַת אֶחָת\hdot
אִם יֵשׁ בַּסַּכִּין מְלֹא צַוָּאר כּשְׁרֵהָ\hdot
הָיָה שׁוֹחֵט וְהִתִּיז שְׁנֵי רָאשִׁין בְּבַת אֶחָת\hdot
אִם יֵשׁ בַּסַּכִּין מְלֹא צַוָּאר אֶחָד כְּשְׁרָה\hdot
בַּמֶּה דְּבָרִים אֲמוּרִים בִּזְמָן שֶׁהוֹלִיךְ וְלֹא הֵבִיא\hdot
אוֹ הֵבִיא וְלֹא הוֹלִיךְ\hdot
אֲבָל אִם הוֹלִיךְ וְהֵבִיא\hdot
אֲפִילוּ כָּל שֶׁהוּא\hdot
אֲפִילוּ בְּאִיזְמֵל כְּשֵׁרָה\hdot
נָֽפְלָה סַכִּין וּשְׁחָטָהּ\hdot
אַף עַל פִּי שֶׁשְּׁחֲטָה כְּדַרְכָּהּ פְּסוּלָה\hdot
שֶׁנֶּאֱמַר וְזָבַחְתָּ וְאָכַלְתָּ\hdot
מַה שֶּׁאַתָּה זוֹבֵחַ אַתָּה אוֹכֵל\hdot
נָֽפְלָה סַכִּין וְהִגְבִּיהָהּ\hdot
נָֽפְלוּ כֵלָיו וְהִגְבִּיהָן\hdot
הִשְׁחִיז אֶת הַסַּכִּין וְעַף\hdot
וּבָא חֲבֵרוֹ וְשָׁחַט\hdot
אִם שָׁהָה כְּדֵי שְׁחִיטָה פְּסוּלָה\hdot
רִבְּי שִׁמְעוֹן אוֹמֵר אִם שָׁהָה כְּדֵי בִקּוּר׃}

\halacha{שָׁחַט אֶת הַוֶּשֶׁט וּפָסַק אֶת הַגַּרְגֶּרֶת\hdot
אוֹ שָׁחַט אֶת הַגַּרְגֶּרֶת\hdot
וּפָסַק אֶת הַוֶּשֶׁט\hdot
אוֹ שֶׁשָּׁחַט אֶחָד מֵהֶן\hdot
וְהִמְתִּין לָהּ עַד שֶׁמֵּתָה\hdot
אוֹ שֶׁהֶחֱלִיד אֶת הַסַּכִּין תַּחַת הַשֵּׁנִי וּפְסָקוֹ\hdot
רִבִּי יְשֵׁבָב אוֹמֵר נְבֵלָה\hdot
רִבִּי עֲקִיבָא אוֹמֵר טְרֵפָה\hdot
כְּלָל אָמַר רִבִּי יְשֵׁבָב מִשּׁוּם רִבִּי יְהוֹשֻׁעַ\hdot
כָּל שֶׁנִּפְסְלָה בִּשְׁחִיטָתָהּ נְבֵלָה\hdot
כָּל שֶׁשְּׁחִיטָתָהּ כָּרָאוּי וְדָבָר אַחֵר גָּרַם לָהּ לִפָּסֵל טְרֵפָה\hdot
וְהוֹדָה לוֹ רִבִּי עֲקִיבָא׃}

\halacha{הַשּׁוֹחֵט בְּהֵמָה חַיָּה וְעוֹף \hdot
וְלֹא יָצָא מֵהֶן דָּם כְּשֵׁרִים\hdot
וְנֶאֱכָלִים בְּיָדַיִם מְסוּאֲבוֹת\hdot
לְפִי שֶׁלֹּא הֻכְשְׁרוּ בְדָם\hdot
רִבִּי שִׁמְעוֹן אוֹמֵר הֻכְשְׁרוּ בַשְּׁחִיטָה׃}

\halacha{הַשּׁוֹחֵט אֶת הַמְּסֻכֶּנֶת\hdot
רַבָּן (שִׁמְעוֹן בֶּן) גַּמְלִיאֵל אוֹמֵר עַד שֶׁתְּפַרְכֵּס בַּיָּד וּבָרֶגֶל\hdot
רִבִּי אֱלִיעֶזֶר אוֹמֵר דַּיָּהּ אִם זִנְּקָה\hdot
אָמַר רִבִּי שִׁמְעוֹן (אף) הַשּׁוֹחֵט בַּלַּיְלָה וּלְמָחָר הִשְׁכִּים\hdot
וּמָצָא כֳּתָלִים מְלֵאִים דָּם כְּשֵׁרָה שֶׁזִּנְּקָה\hdot
וּכְמִדַּת רִבִּי אֱלִיעֶזֶר\hdot
וַחֲכָמִים אוֹמְרִים עַד שֶׁתְּפַרְכֵּס בַּיָּד וּבָרֶגֶל\hdot
אוֹ עַד שֶׁתְּכַשְׁכֵּשׁ בִּזְנָבָהּ\hdot
אֶחָד בְּהֵמָה דַקָּה\hdot
וְאֶחָד בְּהֵמָה גַסָּה\hdot
בְּהֵמָה דַקָּה שֶׁפָּֽשְׁטָה יָדָהּ וְלֹא הֶחֱזְירָהּ פְּסוּלָה\hdot
שֶׁאֵינָהּ אֶלָּא הוֹצָאַת נֶפֶשׁ בִּלְבָד\hdot
בַּמֶּה דְבָרִים אֲמוּרִים שֶׁהָֽיְתָה בְחֶזְקַת מְסֻכֶּנֶת\hdot
אֲבָל אִם הָֽיְתָה בְחֶזְקַת בְּרִיאָה\hdot
אֲפִילוּ אֵין בָּהּ אֶחָד מִכָּל הַסִּימָנִים הַלָּלוּ כְּשֵׁרָה׃}

%\input{05_04_bechorot}
%\input{05_05_arachin}
%\input{05_06_temura}
%\input{05_07_keritot}
%\input{05_08_meila}
%\input{05_09_tamid}
%\input{05_10_midot}
%\input{05_11_kinim}

\clearpage

\thispagestyle{empty}
\centerlastline
\small
\parskip 1pt
{\centering
  \scriptsize
   אחר השלמת כל מסכתא ומסכתא יאמר זה שלש פעמים\\ ויועיל לשכחה בעזרת השם יתברך\hdot
}

הַדְרָן עֲלָךְ מַסֶּכֶת {\scriptsize(פלוני)} וְהַדְרָךְ עֲלָן\hdot
דַּעְתָן עֲלָךְ מַסֶּכֶת {\scriptsize(פלוני)} וְדַעְתָּך עֲלָן\hdot
לָא נִתְנְשֵׁי מִינָךְ מַסֶּכֶת {\scriptsize(פלוני}) וְלָא תִתְנְשֵׁי מִינָן
לָא בְעַלְמָא הַדֵּין וְלָא בְעַלְמָא דְאָתֵי׃

\vspace*{1em}

{\scriptsize ואחר כל יאמר}

יְהִי רָצוֹן מִלְּפָנֶיךָ יְיָ אֱלֹהֵינוּ וֵאלֹהֵי אֲבוֹתֵינוּ שֶׁתְּהֵא תוֹרָֽתְךָ אוּמָנוּתֵנוּ בָּעוֹלָם הַזֶּה וּתְהֵא עִמָּנו לְעוֹלָם הַבָּא\hdot
חֲנַנְיָה בַּר פַּפָּא\hdot
רָמֵי בַּר פַּפָּא\hdot
נַחְמָן בַּר פַּפָּא\hdot
אָחַאי בַּר פַּפָּא\hdot
אַבָּא מָרִי בַּר פַּפָּא\hdot
רַפְרָם בַּר פַּפָּא\hdot
רָכִישׁ בַּר פַּפָּא\hdot
סֻרְחַב בַּר פַּפָּא\hdot
אָדָא בַּר פַּפָּא\hdot
דָרוֹ בַּר פַּפָּא׃

\vspace*{1em}

הַעֲרֵב נָא יְיָ אֱלֹהֵינוּ אֶת דִּבְרֵי תוֹרָֽתְךָ בְפִינוּ וּבְפִיפִיּוֹת עַמְּךָ בֵּית יִשְׂרָאֵל וְנִהְיֶה אֲנַחְנוּ כֻּלָּנוּ וְצֶאֱצָאֵינוּ וְצֶאֱצָאֵי עַמְּךָ בֵּית יִשְׂרָאֵל יִרְאֵי שְׁמֶךָ וְלוֹמְדֵי תוֹרָֽתֶךָ׃

מֵאֹיְבַי תְּחַכְּמֵנִי מִצְוֺתֶךָ כִּי לְעוֹלָם הִיא לִי׃

יְהִי לִבִּי תָמִים בְּחֻקֶּיךָ לְמַעַן לֹא אֵבוֹשׁ׃

לְעוֹלָם לֹא אֶשְׁכַּח פִּקּוּדֶיךָ כִּי בָם חִיִּיתָנִי׃

  בָּרוּךְ אַתָּה יְהוָה לַמְּדֵנִי חֻקֶּיךָ׃



%\def \sederCount{חמישי}
%\clearpage

\centerlastline
{
\small
\parskip 1pt
{\centering
  \scriptsize
   אחר השלמת כל מסכתא ומסכתא יאמר זה שלש פעמים\\ ויועיל לשכחה בעזרת השם יתברך\hdot
}

\begin{tp}{219}
הַדְרָן עֲלָךְ מַסֶּכֶת {\scriptsize(פלוני)} וְהַדְרָךְ עֲלָן\hdot
דַּעְתָן עֲלָךְ מַסֶּכֶת {\scriptsize(פלוני)} וְדַעְתָּך עֲלָן\hdot
לָא נִתְנְשֵׁי מִינָךְ מַסֶּכֶת {\scriptsize(פלוני}) וְלָא תִתְנְשֵׁי מִינָן
לָא בְעַלְמָא הַדֵּין וְלָא בְעַלְמָא דְאָתֵי׃
\end{tp}

{\scriptsize ואחר כל יאמר}

יְהִי רָצוֹן מִלְּפָנֶיךָ יְיָ אֱלֹהֵינוּ וֵאלֹהֵי אֲבוֹתֵינוּ שֶׁתְּהֵא תוֹרָֽתְךָ אוּמָנוּתֵנוּ בָּעוֹלָם הַזֶּה וּתְהֵא עִמָּנו לְעוּלָם הַבָּא\hdot
חֲנַנְיָה בַּר פַּפָּא\hdot
רָמֵי בַּר פַּפָּא\hdot
נַחְמָן בַּר פַּפָּא\hdot
אָחַאי בַּר פַּפָּא\hdot
אַבָּא מָרִי בַּר פַּפָּא\hdot
רַפְרָם בַּר פַּפָּא\hdot
רָכִישׁ בַּר פַּפָּא\hdot
סֻרְחַב בַּר פַּפָּא\hdot
אָדָא בַּר פַּפָּא\hdot
דָרוֹ בַּר פַּפָּא׃

\hbadness 64
\begin{tp}{2048}
הַעֲרֵב נָא יְיָ אֱלֹהֵינוּ אֶת דִּבְרֵי תוֹרָֽתְךָ בְפִינוּ וּבְפִיפִיּוֹת עַמְּךָ בֵּית יִשְׂרָאֵל וְנִהְיֶה אֲנַחְנוּ כֻּלָּנוּ וְצֶאֱצָאֵינוּ וְצֶאֱצָאֵי עַמְּךָ בֵּית יִשְׂרָאֵל יִרְאֵי שְׁמֶךָ וְלוֹמְדֵי תוֹרָֽתֶךָ\hdot
 מֵאֹיְבַי תְּחַכְּמֵנִי מִצְו‍ֹתֶךָ כִּי לְעוֹלָם הִיא לִי\hdot
 יְהִי לִבִּי תָמִים בְּחֻקֶּיךָ לְמַעַן לֹא אֵבוֹשׁ\hdot
 לְעוֹלָם לֹא אֶשְׁכַּח פִּקּוּדֶיךָ כִּי בָם חִיִּיתָנִי\hdot
  בָּרוּךְ אַתָּה יְהוָה לַמְּדֵנִי חֻקֶּיךָ׃
\end{tp}
  
    {\centering
      \larger חתימת חלק ראשון מהמשנה

    }

    לְהוֹדוֹת לְהַלֵּל לְשַׁבֵּחַ לְפָאֵר לְרוֹמֵם לְהַדֵּר וּלְנַצֵּחַ עַל כָּל דִּבְרֵי שִׁירוֹת וְתִשְׁבָּחוֹת דָּוִד בֶּן יִשַּׁי עַבְדְּךָ מְשִׁיחֶךָ׃ בַּעֲצָֽתְךָ תַנְחֵנִי וְאַחַר כָּבוֹד תִּקָּחֵנִי׃
    וְהוּא רַחוּם יְכַפֵּר עָו‍ֹן וְלֹא יַשְׁחִית וְהִרְבָּה לְהָשִׁיב אַפּוֹ וְלֹא יָעִיר כָּל חֲמָתוֹ׃
    אַשְׁרֵי הַגֶּבֶר אֲשֶׁר תְּיַסְּרֶנּוּ יָּהּ וּמִתּוֹרָֽתְךָ תְלַמְּדֶנּוּ׃
    וַאֲנִי בְּחַסְדְּךָ בָטַחְתִּי יָגֵל לִבִּי בִּישׁוּעָתֶךָ אָשִׁירָה לַיְיָ כִּי גָמַל עָלָי׃
    שׂוֹשׂ אָשִׂישׂ בַּיְיָ תָּגֵל נַפְשִׁי בֵּאלֹהַי כִּי הִלְבִּישַׁנִי בִּגְדֵי יֶשַׁע מְעִיל צְדָקָה יְעָטָנִי כֶּחָתָן יְכַהֵן פְּאֵר וְכַכַּלָּה תַּעְדֶּה כֵלֶיהָ׃
    וְיִבְטְחוּ בְךָ יוֹדְעֵי שְׁמֶךָ כִּי לֹא עָזַבְתָּ דֹרְשֶׁיךָ יְהוָה׃
    שִׂמְחוּ בַיהוָה וְגִילוּ צַדִּיקִים וְהַרְנִינוּ כָּל יִשְׁרֵי לֵב׃
    נֵר לְרַגְלִי דְבָרֶךָ וְאוֹר לִנְתִיבָתִי׃
    אוֹדְךָ כִּי עֲנִיתָנִי וַתְּהִי לִי לִישׁוּעָה׃
    
    {\centering
    {\scriptsize    ואומר השיר שהיו הלוים אומרים על הדוכן בכל יום}
    
      \larger ביציאתו מבית המדרש אומר

    }

\begin{tp}{74}
    מוֹדֶה אֲנִי לְפָנֶיךָ יְיָ אֱלֹהַי שֶׁשַּׂמְתָּ חֶלְקִי מִיּוֹשְׁבֵי בֵית הַמֶּדְרָשׁ וְלֹא שַמְתָּ חֶלְקִי מִיּוֹשְׁבֵי קְרָנוֹת
    שֶׁאֲנִי מַשְׁכִּים וְהֵם מַשְׁכִּימִים אֲנִי מַשְׁכִּים לְדִבְרֵי תוֹרָה וְהֵם מַשְׁכִּימִים לִדְבָרֵים בְּטֵלִים 
    אֲנִי עָמֵל וְהֵם עֲמֵלִים אֲנִי עָמֵל וּמְקַבֵּל שָׂכָר וְהֵם עֲמֵלִים וְאֵינָם מְקַבְּלִים שָׂכָר
    אֲנִי רָץ וְהֵם רָצִים אֲנִי רָץ לְחַיֵּי הָעוֹלָם הַבָּא וְהֵם רָצִים לִבְאֵר שַׁחַת\hdot
    שֶׁנֶּאֱמֶר וְאַתָּה אֱלֹהִים תּוֹרִדֵם לִבְאֵר שַׁחַת אַנְשֵׁי דָמִים וּמִרְמָה לֹא יֶחֱצוּ יְמֵיהֶם וַאֲנִי אֶבְטַח בָּךְ׃
\end{tp}


{\centering
  \LARGE תם ונשלם ספר ראשון\\שבחי תודה לאל עליון
}

}

%
%{\centering
%  \LARGE תם ונשלם ספר חמישי\\שבחי תודה לצור קדשי
%}

\end{hebrew}
\end{document}
